% Options for packages loaded elsewhere
\PassOptionsToPackage{unicode}{hyperref}
\PassOptionsToPackage{hyphens}{url}
\PassOptionsToPackage{space}{xeCJK}
%
\documentclass[
  12pt,
]{article}
\usepackage{lmodern}
\usepackage{amssymb,amsmath}
\usepackage{ifxetex,ifluatex}
\ifnum 0\ifxetex 1\fi\ifluatex 1\fi=0 % if pdftex
  \usepackage[T1]{fontenc}
  \usepackage[utf8]{inputenc}
  \usepackage{textcomp} % provide euro and other symbols
\else % if luatex or xetex
  \usepackage{unicode-math}
  \defaultfontfeatures{Scale=MatchLowercase}
  \defaultfontfeatures[\rmfamily]{Ligatures=TeX,Scale=1}
  \setmainfont[]{DejaVu Serif}
  \setsansfont[]{DejaVu Sans}
  \setmonofont[]{DejaVu Sans Mono}
  \ifxetex
    \usepackage{xeCJK}
    \setCJKmainfont[]{Noto Serif CJK SC}
  \fi
  \ifluatex
    \usepackage[]{luatexja-fontspec}
    \setmainjfont[]{Noto Serif CJK SC}
  \fi
\fi
% Use upquote if available, for straight quotes in verbatim environments
\IfFileExists{upquote.sty}{\usepackage{upquote}}{}
\IfFileExists{microtype.sty}{% use microtype if available
  \usepackage[]{microtype}
  \UseMicrotypeSet[protrusion]{basicmath} % disable protrusion for tt fonts
}{}
\makeatletter
\@ifundefined{KOMAClassName}{% if non-KOMA class
  \IfFileExists{parskip.sty}{%
    \usepackage{parskip}
  }{% else
    \setlength{\parindent}{0pt}
    \setlength{\parskip}{6pt plus 2pt minus 1pt}}
}{% if KOMA class
  \KOMAoptions{parskip=half}}
\makeatother
\usepackage{xcolor}
\IfFileExists{xurl.sty}{\usepackage{xurl}}{} % add URL line breaks if available
\IfFileExists{bookmark.sty}{\usepackage{bookmark}}{\usepackage{hyperref}}
\hypersetup{
  hidelinks,
  pdfcreator={LaTeX via pandoc}}
\urlstyle{same} % disable monospaced font for URLs
\usepackage[margin=2.5cm]{geometry}
\usepackage{color}
\usepackage{fancyvrb}
\newcommand{\VerbBar}{|}
\newcommand{\VERB}{\Verb[commandchars=\\\{\}]}
\DefineVerbatimEnvironment{Highlighting}{Verbatim}{commandchars=\\\{\}}
% Add ',fontsize=\small' for more characters per line
\newenvironment{Shaded}{}{}
\newcommand{\AlertTok}[1]{\textcolor[rgb]{1.00,0.00,0.00}{\textbf{#1}}}
\newcommand{\AnnotationTok}[1]{\textcolor[rgb]{0.38,0.63,0.69}{\textbf{\textit{#1}}}}
\newcommand{\AttributeTok}[1]{\textcolor[rgb]{0.49,0.56,0.16}{#1}}
\newcommand{\BaseNTok}[1]{\textcolor[rgb]{0.25,0.63,0.44}{#1}}
\newcommand{\BuiltInTok}[1]{#1}
\newcommand{\CharTok}[1]{\textcolor[rgb]{0.25,0.44,0.63}{#1}}
\newcommand{\CommentTok}[1]{\textcolor[rgb]{0.38,0.63,0.69}{\textit{#1}}}
\newcommand{\CommentVarTok}[1]{\textcolor[rgb]{0.38,0.63,0.69}{\textbf{\textit{#1}}}}
\newcommand{\ConstantTok}[1]{\textcolor[rgb]{0.53,0.00,0.00}{#1}}
\newcommand{\ControlFlowTok}[1]{\textcolor[rgb]{0.00,0.44,0.13}{\textbf{#1}}}
\newcommand{\DataTypeTok}[1]{\textcolor[rgb]{0.56,0.13,0.00}{#1}}
\newcommand{\DecValTok}[1]{\textcolor[rgb]{0.25,0.63,0.44}{#1}}
\newcommand{\DocumentationTok}[1]{\textcolor[rgb]{0.73,0.13,0.13}{\textit{#1}}}
\newcommand{\ErrorTok}[1]{\textcolor[rgb]{1.00,0.00,0.00}{\textbf{#1}}}
\newcommand{\ExtensionTok}[1]{#1}
\newcommand{\FloatTok}[1]{\textcolor[rgb]{0.25,0.63,0.44}{#1}}
\newcommand{\FunctionTok}[1]{\textcolor[rgb]{0.02,0.16,0.49}{#1}}
\newcommand{\ImportTok}[1]{#1}
\newcommand{\InformationTok}[1]{\textcolor[rgb]{0.38,0.63,0.69}{\textbf{\textit{#1}}}}
\newcommand{\KeywordTok}[1]{\textcolor[rgb]{0.00,0.44,0.13}{\textbf{#1}}}
\newcommand{\NormalTok}[1]{#1}
\newcommand{\OperatorTok}[1]{\textcolor[rgb]{0.40,0.40,0.40}{#1}}
\newcommand{\OtherTok}[1]{\textcolor[rgb]{0.00,0.44,0.13}{#1}}
\newcommand{\PreprocessorTok}[1]{\textcolor[rgb]{0.74,0.48,0.00}{#1}}
\newcommand{\RegionMarkerTok}[1]{#1}
\newcommand{\SpecialCharTok}[1]{\textcolor[rgb]{0.25,0.44,0.63}{#1}}
\newcommand{\SpecialStringTok}[1]{\textcolor[rgb]{0.73,0.40,0.53}{#1}}
\newcommand{\StringTok}[1]{\textcolor[rgb]{0.25,0.44,0.63}{#1}}
\newcommand{\VariableTok}[1]{\textcolor[rgb]{0.10,0.09,0.49}{#1}}
\newcommand{\VerbatimStringTok}[1]{\textcolor[rgb]{0.25,0.44,0.63}{#1}}
\newcommand{\WarningTok}[1]{\textcolor[rgb]{0.38,0.63,0.69}{\textbf{\textit{#1}}}}
\usepackage{longtable,booktabs}
% Correct order of tables after \paragraph or \subparagraph
\usepackage{etoolbox}
\makeatletter
\patchcmd\longtable{\par}{\if@noskipsec\mbox{}\fi\par}{}{}
\makeatother
% Allow footnotes in longtable head/foot
\IfFileExists{footnotehyper.sty}{\usepackage{footnotehyper}}{\usepackage{footnote}}
\makesavenoteenv{longtable}
\usepackage{graphicx}
\makeatletter
\def\maxwidth{\ifdim\Gin@nat@width>\linewidth\linewidth\else\Gin@nat@width\fi}
\def\maxheight{\ifdim\Gin@nat@height>\textheight\textheight\else\Gin@nat@height\fi}
\makeatother
% Scale images if necessary, so that they will not overflow the page
% margins by default, and it is still possible to overwrite the defaults
% using explicit options in \includegraphics[width, height, ...]{}
\setkeys{Gin}{width=\maxwidth,height=\maxheight,keepaspectratio}
% Set default figure placement to htbp
\makeatletter
\def\fps@figure{htbp}
\makeatother
\setlength{\emergencystretch}{3em} % prevent overfull lines
\providecommand{\tightlist}{%
  \setlength{\itemsep}{0pt}\setlength{\parskip}{0pt}}
\setcounter{secnumdepth}{5}

\author{}
\date{}

\usepackage{amsmath}
\usepackage{amssymb}
\usepackage{longtable}
\usepackage{booktabs}
\usepackage{float}

\begin{document}

{
\setcounter{tocdepth}{3}
\tableofcontents
}
\hypertarget{muxf6biusquantumring-ux6280ux672fux6587ux6863}{%
\section{MöbiusQuantumRing
技术文档}\label{muxf6biusquantumring-ux6280ux672fux6587ux6863}}

\textbf{版本}: 1.0.0\\
\textbf{生成日期}: 2026年01月19日 14:53\\
\textbf{作者}: CursorLike 自动生成\\
\textbf{模板来源}: example\_paper.pdf

\begin{center}\rule{0.5\linewidth}{0.5pt}\end{center}

\hypertarget{ux76eeux5f55}{%
\subsection{目录}\label{ux76eeux5f55}}

\begin{itemize}
\tightlist
\item
  第1章: 001
\item
  第2章: 003
\item
  第3章: 005
\item
  第4章: 007
\item
  第5章: 009
\item
  第6章: 011
\item
  第7章: 013
\item
  第8章: 015
\item
  第9章: 017
\item
  第10章: 019
\item
  第11章: 021
\item
  第12章: 023
\item
  第13章: 025
\item
  第14章: 027
\item
  第15章: 029
\end{itemize}

\begin{center}\rule{0.5\linewidth}{0.5pt}\end{center}

\hypertarget{ux7b2c1ux7ae0-ux5f15ux8a00ux5faaux73afux795eux7ecfux7f51ux7edcux4e2dux7684ux7a33ux5b9aux6027ux4e0eux8868ux8fbeux80fdux529b}{%
\section{第1章
引言：循环神经网络中的稳定性与表达能力}\label{ux7b2c1ux7ae0-ux5f15ux8a00ux5faaux73afux795eux7ecfux7f51ux7edcux4e2dux7684ux7a33ux5b9aux6027ux4e0eux8868ux8fbeux80fdux529b}}

\hypertarget{ux7814ux7a76ux80ccux666fux4e0eux52a8ux673a}{%
\subsection{1.1
研究背景与动机}\label{ux7814ux7a76ux80ccux666fux4e0eux52a8ux673a}}

在人工智能技术发展的宏大叙事中，循环神经网络（Recurrent Neural Networks,
RNNs）及其衍生架构构成了处理序列化、时序化信息的基石。从自然语言处理中理解句子的上下文语义，到金融科技领域预测股票价格的波动趋势，再到生物信息学中分析基因序列的潜在模式，RNN类模型为解决具有时间依赖性或顺序结构的问题提供了核心的计算范式。然而，随着模型复杂度的提升和应用场景的扩展，一个长期困扰研究者和工程师的根本性挑战日益凸显：\textbf{模型内在的表达能力与其在训练和推理过程中的数值稳定性之间存在深刻的、难以调和的矛盾}。这一矛盾并非简单的超参数调整或工程优化所能解决，它根植于动态系统理论、数值优化和表示学习的交叉地带，直接关系到模型能否可靠地学习长期依赖关系，并最终决定了其实际应用的边界。本项目``莫比乌斯量子环形网络''的提出，正是为了从第一性原理出发，探索并构建一种能够从根本上调和这一矛盾的新型神经网络架构。

\textbf{表达能力}，在深度学习的语境中，通常指一个模型能够拟合或逼近复杂函数关系的能力。对于序列建模而言，高表达能力意味着模型能够捕捉数据中跨越长距离的、非线性的相互作用。这种能力是模型性能的上限，它依赖于网络结构的深度（时间展开的长度）、宽度（隐藏状态的维度）、非线性激活函数的类型以及参数空间的自由度。一个具备强大表达能力的RNN，其隐藏状态的动力学应能形成丰富的吸引子景观，足以模拟现实世界中千变万化的时序模式。

\textbf{稳定性}，则关注于信息在网络中前向传播（推理）和误差梯度反向传播（训练）过程中的动态行为。在深度RNN的训练中，最经典且棘手的问题便是\textbf{梯度消失}和\textbf{梯度爆炸}。当网络沿时间维度展开成一个很深的计算图时，损失函数关于早期时间步参数的梯度需要通过链式法则进行反向传播。如果每个时间步的循环变换（通常是一个线性变换接一个非线性激活）对梯度向量的雅可比矩阵的谱半径持续小于1，梯度幅值将呈指数级衰减，导致远距离时间步的贡献几乎为零，模型无法学习长期依赖。反之，若谱半径持续大于1，梯度幅值则会指数级增长，引发权重更新剧烈震荡、数值溢出，最终导致训练过程崩溃。

为了在实践中缓解这些问题，研究者们提出了诸如长短期记忆网络（LSTM）和门控循环单元（GRU）等架构。它们通过引入精心设计的门控机制（如输入门、遗忘门、输出门），显式地控制信息流的保留与遗忘，从而在经验上显著改善了模型处理长序列的能力。然而，这些门控机制的稳定性和收敛性保证大多基于启发式设计和广泛的实验验证，缺乏严格的数学理论支撑。同时，门控结构也引入了额外的参数和计算开销。

另一条更具理论深度的路径是从线性代数的等距变换概念中汲取灵感。一个\textbf{正交矩阵}（在实数域）或\textbf{酉矩阵}（在复数域）作用于向量时，能够完美保持向量的L2范数不变。将这一性质应用于RNN的循环变换，意味着在纯线性部分，信号的能量在传播过程中既不会被放大也不会被衰减，从而从根源上消除了因线性权重矩阵引起的梯度爆炸或消失问题。基于此，\textbf{正交RNN}和\textbf{酉循环神经网络}应运而生。这类模型通过特殊的参数化方法（如Cayley变换、指数映射或直接在Stiefel流形/酉群上进行优化），将循环权重矩阵严格约束在等距变换的集合内。这种方法在数学上非常优雅，提供了坚实的稳定性理论保证，并在一些需要处理超长序列依赖的任务上展现了潜力。

然而，追求极致的稳定性往往需要付出代价。严格的等距约束可能过度限制了权重矩阵的灵活性，使其难以实现信息在隐藏空间中的充分混合与复杂重组。在紧致的李群（如酉群）上进行优化比在无约束的欧几里得空间中更为复杂和低效。更重要的是，一个完全保持距离的变换，其``混合能力''可能存在理论上的上限，这可能阻碍模型学习某些需要高度非线性扭曲的复杂模式。因此，我们面临一个根本性的设计困境：无约束的权重矩阵（如传统RNN）赋予模型最大的表达自由度，但极易导致训练不稳定；而完全约束的等距矩阵（如严格正交RNN）确保了卓越的稳定性，却可能以牺牲部分表达丰富性为代价。

此外，主流的RNN架构普遍遵循一种``有向、顺序的前馈计算图''范式。信息从输入层流入，沿着网络层或时间步单向、逐层流动，经过一系列确定性变换，最终到达输出层。这种范式与自然界中许多系统达到平衡的过程，或某些分布式计算模型存在哲学上的差异。例如，在物理系统（如热力学平衡）、生物神经网络（如全局活动模式的稳定）或共识算法中，系统的全局稳态往往是通过大量局部单元基于固定规则持续相互作用，并在外部输入的驱动下，经过一个\textbf{弛豫}或\textbf{迭代收敛}过程而自发涌现的。这种``达到平衡''的计算模式，暗示了一种潜在的、不同于传统前馈架构的新型网络动力学。

\hypertarget{ux6838ux5fc3ux77dbux76feux8868ux8fbeux80fdux529bux4e0eux7a33ux5b9aux6027ux7684ux6743ux8861}{%
\subsection{1.2
核心矛盾：表达能力与稳定性的权衡}\label{ux6838ux5fc3ux77dbux76feux8868ux8fbeux80fdux529bux4e0eux7a33ux5b9aux6027ux7684ux6743ux8861}}

为了更深入地理解表达能力与稳定性之间的权衡，我们需要从数学和动力系统的角度进行剖析。

考虑一个最简单的循环神经网络单元，其隐藏状态 ( \textbf{h}\_t )
的更新公式为： {[} \textbf{h}\emph{t = \sigma(\textbf{W}
\textbf{h}}\{t-1\} + \textbf{U} \textbf{x}\_t + \textbf{b}) {]} 其中，(
\textbf{W} ) 是循环权重矩阵，( \sigma )
是逐点非线性激活函数（如tanh或ReLU），( \textbf{x}\_t ) 是t时刻的输入，(
\textbf{U} ) 是输入权重矩阵，( \textbf{b} ) 是偏置项。

在反向传播通过时间（BPTT）算法中，损失函数 ( L ) 关于早期状态 (
\textbf{h}\emph{k ) 的梯度可以写为： {[}
\frac{\partial L}{\partial \textbf{h}_k} =
\frac{\partial L}{\partial \textbf{h}_T} \prod}\{t=k+1\}\^{}\{T\}
\frac{\partial \textbf{h}_t}{\partial \textbf{h}_{t-1}} =
\frac{\partial L}{\partial \textbf{h}_T} \prod\_\{t=k+1\}\^{}\{T\}
\textbf{D}\_t \textbf{W}\^{}\top {]} 这里，( \textbf{D}\emph{t )
是一个对角矩阵，其对角线元素是激活函数 ( \sigma ) 在 (
\textbf{h}}\{t-1\} ) 处的导数。乘积项 ( \prod \textbf{D}\_t
\textbf{W}\^{}\top ) 决定了梯度传播的长期行为。

\textbf{梯度消失/爆炸的根源}：该乘积项的范数（或谱半径）决定了梯度是收缩还是膨胀。如果
( \textbf{W} ) 的最大奇异值（谱范数）显著大于1，且激活函数导数 (
\textbf{D}\_t )
不提供足够的收缩，梯度就可能爆炸。反之，如果最大奇异值显著小于1，梯度就会消失。即使
( \textbf{W} ) 是正交的，保证了线性部分的稳定性，但非线性激活函数 (
\sigma ) 的导数（例如tanh在饱和区的导数接近0）仍可能导致梯度收缩。

\textbf{表达能力的根源}：模型的表达能力很大程度上取决于 ( \textbf{W} )
所能实现的变换的丰富性。一个任意的、稠密的 ( \textbf{W} )
可以表示高维空间中的任意线性变换，包括旋转、缩放和剪切。这种灵活性对于学习复杂的序列模式至关重要。然而，正是这种任意性导致了谱范数难以控制，进而引发稳定性问题。

\textbf{正交/酉约束的利弊}：将 ( \textbf{W} )
约束为正交或酉矩阵，强制其所有奇异值为1，完美解决了线性部分的梯度范数保持问题。但这意味着
( \textbf{W} )
只能表示旋转（和反射）变换，失去了缩放和剪切的能力。从信息混合的角度看，一个纯粹的旋转可能无法充分``打乱''隐藏状态中的信息，限制了其表示复杂函数的能力。此外，在训练中保持矩阵的正交性需要专门的优化算法，增加了计算复杂性。

因此，设计一个理想的循环架构，本质上是在寻找一个高维参数空间中的``甜蜜点''，该点既能提供足够的变换灵活性以捕获复杂模式，又能确保动力系统的长期行为是良定的、可训练的。这要求我们超越简单的权重矩阵约束，从网络动力学的全局视角出发，重新思考循环计算的根本形式。

\hypertarget{ux8fc8ux5411ux65b0ux8303ux5f0fux4eceux987aux5e8fux524dux9988ux5230ux73afux5f62ux5f1bux8c6b}{%
\subsection{1.3
迈向新范式：从顺序前馈到环形弛豫}\label{ux8fc8ux5411ux65b0ux8303ux5f0fux4eceux987aux5e8fux524dux9988ux5230ux73afux5f62ux5f1bux8c6b}}

传统RNN的``顺序前馈''范式将时间维度显式地展开为计算图的深度。这种视角自然但存在局限：它强制信息沿单一方向流动，并且将``时间步''与计算步骤严格绑定。另一种等价的视角是将循环网络视为一个具有反馈连接的动态系统，其隐藏状态根据一个差分或微分方程演化。在推理时，我们实际上是在迭代求解这个动态系统的稳态或轨迹。

受物理学中弛豫过程和分布式计算中共识算法的启发，我们可以构想一种不同的网络范式：\textbf{环形弛豫网络}。在这种范式中：
1.
\textbf{网络结构}：一组神经元（或处理单元）排列成一个环状拓扑结构，每个神经元与环上的其他神经元通过固定的连接规则（如双随机连接矩阵）相互作用。
2.
\textbf{计算过程}：推理不再是单向的前馈计算，而是一个迭代过程。系统从一个初始状态开始，根据其内部动力学（由连接矩阵定义）和外部输入驱动，进行多次迭代更新。
3.
\textbf{收敛目标}：该迭代过程被设计为收敛到一个唯一的\textbf{固定点}（平衡态）。这个固定点包含了输入信息经过网络充分混合、扩散后的全局表示。
4.
\textbf{读出机制}：最终，从收敛后的稳态环上，选取局部或全局的信息进行读出，得到任务的输出。

这种``环形弛豫''范式具有几个潜在优势： *
\textbf{内在稳定性}：如果连接矩阵和更新规则被精心设计（例如，使用收缩映射），那么迭代过程保证收敛，从根本上解决了前向传播的稳定性问题。
*
\textbf{解耦时间与计算}：迭代步数不再对应序列中的时间步，而是达到平衡所需的计算量。这为处理不同长度的输入提供了统一的框架。
*
\textbf{全局信息整合}：环形结构结合弛豫过程，允许信息在环上循环传播，理论上可以实现任意距离的信息交互，有利于捕获长期依赖。
*
\textbf{生物合理性}：更接近某些生物神经网络中全局活动模式通过局部相互作用达到稳定的现象。

然而，实现这样一个范式面临关键挑战：如何设计连接矩阵和更新规则，使其既能保证收敛性（稳定性），又能实现强大的计算能力（表达能力）？连接矩阵需要具备什么样的数学性质？

\hypertarget{ux5e7aux6a21ux968fux673aux6d41ux5f62ux7a33ux5b9aux6027ux4e0eux8868ux8fbeux80fdux529bux7684ux4ea4ux6c47ux70b9}{%
\subsection{1.4
幺模随机流形：稳定性与表达能力的交汇点}\label{ux5e7aux6a21ux968fux673aux6d41ux5f62ux7a33ux5b9aux6027ux4e0eux8868ux8fbeux80fdux529bux7684ux4ea4ux6c47ux70b9}}

``莫比乌斯量子环形网络''的核心理论支柱之一是\textbf{幺模随机流形}的概念，它为我们提供了构建环形弛豫网络所需的关键数学工具。

首先，我们定义几个概念： *
\textbf{双随机矩阵}：一个非负方阵，其每一行和每一列的元素之和都等于1。双随机矩阵可以看作是一个概率转移矩阵，描述了一个马尔可夫链的状态转移。其主导特征值为1，对应的特征向量是均匀分布，这保证了在多次迭代后系统会趋向于一个均匀的稳态（如果不考虑非线性）。
* \textbf{酉矩阵}：一个复数方阵 ( U )，满足 ( U\^{}\dagger U = I )，其中
( U\^{}\dagger ) 是 ( U )
的共轭转置。酉矩阵是等距变换，保持向量的L2范数。 *
\textbf{幺模随机矩阵}：对于一个酉矩阵 ( U )，我们定义 ( H =
\textbar U\textbar\^{}2 )，即 ( H\_\{ij\} =
\textbar U\_\{ij\}\textbar\^{}2 )。可以证明，这样得到的 ( H )
是一个\textbf{双随机矩阵}。所有由N维酉矩阵通过这种``取模平方''操作得到的双随机矩阵构成的集合，称为\textbf{幺模随机流形}。

幺模随机矩阵 ( H ) 具有以下迷人性质： 1.
\textbf{双随机性保证收敛性}：作为双随机矩阵，当用于线性状态更新 (
\textbf{h} \leftarrow \textbf{h} H )
时，在一定的正则条件下（如Perron-Frobenius定理），系统会收敛到一个稳态。这为环形网络的弛豫过程提供了稳定性基础。
2. \textbf{源自酉矩阵的表达能力}：矩阵 ( H ) 的元素由酉矩阵 ( U )
的元素的模平方决定。虽然 ( H ) 本身是非负的、双随机的，但其``父''矩阵 (
U ) 存在于高维、紧致的酉群上。通过学习和优化 ( U )，我们可以间接地塑造 (
H )，从而利用酉群丰富的几何结构来增强模型的表达能力。( U )
的相位信息虽然在 ( H )
中丢失，但可以通过其他机制（如复数域动力学）重新引入。 3.
\textbf{能量守恒的隐喻}：在量子力学中，酉演化保证概率守恒。在这里，( H )
作为概率转移矩阵，其双随机性也保证了一种``概率质量''的守恒。这为网络提供了一种结构化的、可控的信息流动方式。

因此，幺模随机流形成为了连接稳定性（通过双随机性 ( H )
）和表达能力（通过酉矩阵 ( U ) ）的理想桥梁。MöbiusQuantumRing
网络的核心思想，就是学习一个参数化的酉矩阵 ( U
)，构造其对应的幺模随机连接矩阵 ( H = \textbar U\textbar\^{}2
)，以此作为环形网络上神经元相互作用的蓝图，并在此基础上构建一个具有固定点收敛保证的弛豫推理过程。

\hypertarget{ux672cux7ae0ux6982ux8981ux4e0eux6587ux6863ux7ed3ux6784}{%
\subsection{1.5
本章概要与文档结构}\label{ux672cux7ae0ux6982ux8981ux4e0eux6587ux6863ux7ed3ux6784}}

本章系统性地阐述了循环神经网络领域长期存在的核心矛盾------表达能力与稳定性的权衡，并以此作为引子，逐步推导出``莫比乌斯量子环形网络''设计理念的理论动机。我们从传统RNN的梯度问题出发，探讨了正交约束的解决方案及其局限性，进而提出了从``顺序前馈''到``环形弛豫''的范式转变的可能性。最后，我们引入了幺模随机流形这一关键数学概念，展示了它如何为构建既稳定又富有表达力的环形弛豫网络提供理论基石。

本技术文档的后续章节将围绕 MöbiusQuantumRing
的具体实现进行全面而深入的剖析。文档结构如下：

\begin{itemize}
\tightlist
\item
  \textbf{第2章
  核心架构总览}：将宏观介绍MQR网络的整体设计哲学、主要组件及其交互关系，提供一个全局视角。
\item
  \textbf{第3章 幺模随机连接矩阵的构建与优化}：深入探讨酉矩阵 ( U )
  的参数化方法（如Cayley变换）、( H = \textbar U\textbar\^{}2 )
  的性质，以及在训练中保持或优化 ( U ) 的技术细节。
\item
  \textbf{第4章
  环形动力学与固定点收敛性分析}：详细推导环形网络的弛豫更新方程，从动力系统理论的角度严格证明其收敛到唯一固定点的条件。
\item
  \textbf{第5章
  输入注入机制：LoRA式驱动}：阐述如何将外部输入信息以低秩、高效的方式注入到环形动力系统中，驱动系统达到与输入相关的平衡态。
\item
  \textbf{第6章
  输出读出机制：局部采样与投影}：说明如何从收敛后的环形稳态中提取有效信息，形成任务输出，包括局部节点采样和可学习的读出权重。
\item
  \textbf{第7章
  非线性扩展与激活函数设计}：讨论在保持收敛性的前提下，如何在注入端和环内动力学中引入非线性激活函数，以增强模型的表达能力。
\item
  \textbf{第8章 复数域动力学扩展}：探索超越实数域，在复数域利用酉矩阵 (
  U ) 本身（而非 ( H ) ）进行动力学传播的方案，以及其理论含义和实现。
\item
  \textbf{第9章
  高级特性与变体}：介绍项目支持的高级功能，如双酉矩阵解耦（世界模型/策略分离）、可学习目标平衡态、自保持混合系数等。
\item
  \textbf{第10章
  系统实现与代码剖析}：结合项目代码结构，详细讲解各核心模块的实现细节、类设计及API。
\item
  \textbf{第11章
  实验设计与性能评估}：展示MQR网络在标准序列建模任务（如像素级序列预测、语言建模等）上的实验设置、结果分析与对比。
\item
  \textbf{第12章
  收敛性实证分析与可视化}：通过实验可视化弛豫过程的收敛轨迹、环上状态分布等，实证验证理论性质。
\item
  \textbf{第13章
  消融研究与参数敏感性分析}：系统性地研究不同组件（如环大小、收敛容差、注入秩）对模型性能的影响。
\item
  \textbf{第14章
  讨论、局限与未来方向}：总结MQR架构的优势与不足，探讨其理论意义，并展望可能的改进和未来研究方向。
\item
  \textbf{第15章
  结论}：总结全文，重申MQR网络在调和循环神经网络稳定性与表达能力矛盾方面的贡献。
\end{itemize}

通过这份文档，我们希望为读者提供一份理解、复现乃至扩展 MöbiusQuantumRing
项目的完整技术蓝图。

\begin{figure}
\centering
\includegraphics{figures/diagram_1.png}
\caption{图 1}
\end{figure}

\hypertarget{ux7b2c2ux7ae0-ux5e7aux6a21ux968fux673aux6d41ux5f62ux4eceux9149ux7fa4ux5230ux53ccux968fux673aux52a8ux529bux5b66ux7684ux6570ux5b66ux57faux7840}{%
\section{第2章
幺模随机流形：从酉群到双随机动力学的数学基础}\label{ux7b2c2ux7ae0-ux5e7aux6a21ux968fux673aux6d41ux5f62ux4eceux9149ux7fa4ux5230ux53ccux968fux673aux52a8ux529bux5b66ux7684ux6570ux5b66ux57faux7840}}

\hypertarget{ux5f15ux8a00ux4eceux7a33ux5b9aux6027ux7ea6ux675fux5230ux51e0ux4f55ux7ed3ux6784}{%
\subsection{2.1
引言：从稳定性约束到几何结构}\label{ux5f15ux8a00ux4eceux7a33ux5b9aux6027ux7ea6ux675fux5230ux51e0ux4f55ux7ed3ux6784}}

在第1章中，我们深入探讨了循环神经网络中表达能力与稳定性之间的根本矛盾。传统RNN的梯度问题源于其权重矩阵在欧几里得空间中的无约束优化，而正交/酉RNN通过将权重严格约束在Stiefel流形或酉群上，从数学上保证了信号传播的等距性，从而解决了梯度爆炸与消失问题。然而，这种极致的几何约束可能以牺牲模型的表达能力和混合效率为代价。MöbiusQuantumRing网络的核心创新之一，在于引入了一个介于无约束欧氏空间与严格等距约束之间的\textbf{中间几何结构}------\textbf{幺模随机流形}。这一章将系统阐述幺模随机流形的数学定义、性质及其作为环形网络连接矩阵的理论优势，为理解MQR架构的动力学行为奠定严格的数学基础。

幺模随机流形并非凭空创造，而是从量子力学、随机过程与矩阵分析等多个数学物理领域的交叉中自然涌现的概念。它的核心思想是：与其直接约束循环权重矩阵本身，不如约束一个生成该权重矩阵的\textbf{更基本、更具良好几何性质的数学对象}。在MQR中，这个基本对象就是\textbf{酉矩阵}。通过酉矩阵的模平方运算，我们自然地导出一个具有特殊性质的连接矩阵，该矩阵不仅继承了酉矩阵的部分优良特性，还获得了对于概率分布传播至关重要的\textbf{双随机性}。这种``生成式''的约束方式，在保持稳定性的同时，为模型提供了比严格等距约束更丰富的表达能力。

\hypertarget{ux9149ux77e9ux9635ux4e0ecayleyux53c2ux6570ux5316}{%
\subsection{2.2
酉矩阵与Cayley参数化}\label{ux9149ux77e9ux9635ux4e0ecayleyux53c2ux6570ux5316}}

\hypertarget{ux9149ux77e9ux9635ux7684ux5b9aux4e49ux4e0eux57faux672cux6027ux8d28}{%
\subsubsection{2.2.1
酉矩阵的定义与基本性质}\label{ux9149ux77e9ux9635ux7684ux5b9aux4e49ux4e0eux57faux672cux6027ux8d28}}

酉矩阵是复数域上线性代数中的核心概念。一个\(N \times N\)的复方阵\(U\)被称为\textbf{酉矩阵}，如果它满足：

\[
U^\dagger U = U U^\dagger = I_N
\]

其中\(U^\dagger\)表示\(U\)的共轭转置（即Hermitian转置），\(I_N\)是\(N\)阶单位矩阵。所有\(N\)阶酉矩阵的集合构成一个李群，称为\textbf{酉群}，记作\(U(N)\)。酉矩阵具有以下关键性质：
1.
\textbf{保内积与保范数}：对于任意复向量\(x \in C^N\)，有\(\|Ux\|_2 = \|x\|_2\)。这意味着酉变换是\textbf{等距变换}，在复数域上保持向量的2-范数（模长）不变。
2.
\textbf{特征值位于单位圆上}：酉矩阵的所有特征值\(\lambda_i\)满足\(|\lambda_i| = 1\)，即它们都分布在复平面的单位圆周上。
3.
\textbf{可逆性}：酉矩阵总是可逆的，且其逆矩阵等于其共轭转置：\(U^{-1} = U^\dagger\)。
4. \textbf{行列式的模为1}：\(|\det(U)| = 1\)。

在实数域的特例下，满足\(U^T U = I\)的实矩阵称为\textbf{正交矩阵}，其集合构成正交群\(O(N)\)。正交矩阵是酉矩阵在实数域上的对应物，同样具有保内积和保范数的性质。

在RNN的语境下，使用酉矩阵作为循环变换的核心优势正在于其保范性。考虑一个简单的线性循环更新：\(h_t = U h_{t-1}\)。由于\(\|h_t\| = \|U h_{t-1}\| = \|h_{t-1}\|\)，隐藏状态的范数在时间维度上严格保持不变。这不仅在前向传播中防止了激活值的爆炸或衰减，在反向传播中，误差梯度所经过的雅可比矩阵链的奇异值也始终为1，从而彻底消除了由线性部分引起的梯度爆炸或消失问题。

\hypertarget{ux9149ux77e9ux9635ux7684ux53c2ux6570ux5316ux6311ux6218ux4e0ecayleyux53d8ux6362}{%
\subsubsection{2.2.2
酉矩阵的参数化挑战与Cayley变换}\label{ux9149ux77e9ux9635ux7684ux53c2ux6570ux5316ux6311ux6218ux4e0ecayleyux53d8ux6362}}

尽管酉矩阵在理论上完美，但在机器学习实践中直接对其进行优化面临两个主要挑战：
1.
\textbf{约束优化}：标准的梯度下降算法工作在无约束的欧几里得空间。直接在酉群\(U(N)\)这一弯曲的流形上进行优化，需要专门的流形优化技术（如Riemannian梯度下降），这通常计算复杂且不易与现有的深度学习框架集成。
2.
\textbf{表达灵活性}：并非所有酉矩阵都同样容易通过优化过程找到。我们需要一种参数化方法，能够平滑、全覆盖地生成酉群中的元素，同时允许梯度自由流动。

MQR采用\textbf{Cayley变换}作为酉矩阵的参数化工具，巧妙地解决了上述问题。Cayley变换提供了一种将\textbf{斜Hermitian矩阵}（skew-Hermitian
matrix）映射到酉矩阵的可微双射（在除去特征值-1的点外）。

\textbf{定义（斜Hermitian矩阵）}：一个复方阵\(A\)是斜Hermitian的，如果满足\(A^\dagger = -A\)。在实数域下，这对应于斜对称矩阵（\(A^T = -A\)）。

Cayley变换定义如下：

\[
U = (I + A)(I - A)^{-1}
\]

其中\(A\)是一个\(N \times N\)的斜Hermitian矩阵，\(I\)是单位矩阵。可以证明，只要\((I-A)\)可逆（即\(A\)没有特征值1），通过上述公式计算出的\(U\)一定是酉矩阵。反之，对于任意一个不含特征值-1的酉矩阵\(U\)，可以通过逆Cayley变换得到一个斜Hermitian矩阵\(A\)：

\[
A = (U - I)(U + I)^{-1}
\]

Cayley变换的几何意义可以理解为在单位圆盘（对应斜Hermitian矩阵）与单位圆周（对应酉矩阵）之间建立了一个映射。在实现中，我们并不直接存储和优化酉矩阵\(U\)，而是存储和优化一个无约束的复矩阵\(W\)，然后通过以下步骤构造斜Hermitian矩阵\(A\)：

\[
A = \frac{W - W^\dagger}{2}
\]

显然，这样构造的\(A\)满足\(A^\dagger = -A\)。随后，通过Cayley变换得到酉矩阵\(U\)。由于从\(W\)到\(U\)的整个过程由可微的矩阵运算组成（包括矩阵求逆，可通过求解线性系统稳定实现），因此我们可以通过标准的反向传播算法来优化无约束参数\(W\)，而自动确保生成的\(U\)始终位于酉群上。这种参数化方法将流形约束优化问题转化为了一个无约束优化问题，极大地简化了训练过程。

\begin{figure}
\centering
\includegraphics{figures/diagram_2.png}
\caption{图 2}
\end{figure}

\emph{图2.1：基于Cayley变换的酉矩阵参数化流程及其导出的核心性质。该流程将无约束优化与严格的几何约束相结合，是MQR架构稳定性的基石。}

\hypertarget{ux5e7aux6a21ux968fux673aux77e9ux9635ux5b9aux4e49ux6027ux8d28ux4e0eux6784ux9020}{%
\subsection{2.3
幺模随机矩阵：定义、性质与构造}\label{ux5e7aux6a21ux968fux673aux77e9ux9635ux5b9aux4e49ux6027ux8d28ux4e0eux6784ux9020}}

\hypertarget{ux4eceux9149ux77e9ux9635ux5230ux5e7aux6a21ux968fux673aux77e9ux9635}{%
\subsubsection{2.3.1
从酉矩阵到幺模随机矩阵}\label{ux4eceux9149ux77e9ux9635ux5230ux5e7aux6a21ux968fux673aux77e9ux9635}}

幺模随机矩阵是MQR网络中环形连接矩阵\(H\)的数学形态。它的定义直接源于酉矩阵：

\[
H = |U|^2 \quad \text{或等价地} \quad H_{ij} = |U_{ij}|^2
\]

其中\(|\cdot|\)表示取复数的模。也就是说，\(H\)的每个元素是酉矩阵\(U\)对应位置元素的模平方。由于\(U_{ij}\)是复数，\(|U_{ij}|^2 = U_{ij} \overline{U_{ij}}\)，这是一个非负的实数。

\textbf{定义（双随机矩阵）}：一个实方阵\(B\)被称为双随机矩阵，如果其所有元素非负，且每一行之和与每一列之和都等于1。即：
1. \(B_{ij} \ge 0\)，对所有\(i, j\)。 2.
\(\sum_{j} B_{ij} = 1\)，对所有\(i\)（行随机）。 3.
\(\sum_{i} B_{ij} = 1\)，对所有\(j\)（列随机）。

双随机矩阵在概率论中对应一个\textbf{双随机转移矩阵}，描述了一个马尔可夫链中状态之间的转移概率，且该链具有均匀的稳态分布。

\textbf{定理2.1}：由酉矩阵\(U\)通过模平方运算得到的矩阵\(H = |U|^2\)是一个双随机矩阵。

\textbf{证明}： 1.
\textbf{非负性}：由定义\(H_{ij} = |U_{ij}|^2 \ge 0\)，显然成立。 2.
\textbf{行和为1}：考虑\(H\)的第\(i\)行和：\(\sum_{j=1}^N H_{ij} = \sum_{j=1}^N U_{ij} \overline{U_{ij}}\)。这恰好是酉矩阵\(U\)的第\(i\)行向量与自身的内积。由于\(U\)是酉矩阵，其行向量构成一组标准正交基，因此该行向量的内积（即其\(L_2\)范数的平方）为1。
3.
\textbf{列和为1}：同理，考虑\(H\)的第\(j\)列和：\(\sum_{i=1}^N H_{ij} = \sum_{i=1}^N U_{ij} \overline{U_{ij}}\)，这是\(U\)的第\(j\)列向量与自身的内积。由于酉矩阵的列向量也构成一组标准正交基，因此该列和也为1。

证毕。

这种由酉矩阵导出的双随机矩阵\(H\)，被称为\textbf{幺模随机矩阵}。所有\(N\)阶幺模随机矩阵的集合构成了双随机矩阵集合的一个子集，我们称之为\textbf{幺模随机流形}。值得注意的是，并非所有双随机矩阵都是幺模随机的。幺模随机矩阵是双随机矩阵中具有额外结构的一个特殊子类。

\hypertarget{ux5e7aux6a21ux968fux673aux77e9ux9635ux7684ux51e0ux4f55ux4e0eux8c31ux6027ux8d28}{%
\subsubsection{2.3.2
幺模随机矩阵的几何与谱性质}\label{ux5e7aux6a21ux968fux673aux77e9ux9635ux7684ux51e0ux4f55ux4e0eux8c31ux6027ux8d28}}

幺模随机矩阵\(H\)继承了其母矩阵\(U\)的部分几何特性，同时也拥有作为转移矩阵的独特性质。

\begin{enumerate}
\def\labelenumi{\arabic{enumi}.}
\item
  \textbf{非负性与概率解释}：\(H\)的所有元素在\([0, 1]\)之间，且行和与列和为1。这允许我们将环形网络中的隐藏状态\(h \in R^N\)（经过适当的归一化后）解释为一个\textbf{概率分布}或\textbf{概率幅的密度}。环形动力学\(h \leftarrow h H^T\)则可以解释为这个概率分布按照转移矩阵\(H^T\)进行演化。这是MQR``弛豫到平衡''过程的核心概率图景。
\item
  \textbf{谱半径与佩龙-弗罗贝尼乌斯定理}：根据佩龙-弗罗贝尼乌斯定理，一个非负矩阵的谱半径等于其最大特征值的模，且该特征值为非负实数。对于双随机矩阵\(H\)，由于其每行和每列和均为1，向量\(\textbf{1} = (1,1,...,1)^T\)是\(H\)和\(H^T\)的右特征向量，对应的特征值为1：
  \[H \textbf{1} = \textbf{1}, \quad H^T \textbf{1} = \textbf{1}\]
  因此，\(H\)和\(H^T\)的谱半径\(\rho(H) = \rho(H^T) = 1\)。并且，特征值1是最大特征值（按模长）。
\item
  \textbf{特征值分布}：虽然\(H\)由酉矩阵\(U\)生成，但\(H\)的特征值与\(U\)的特征值没有简单的直接关系。\(H\)的特征值都是实数（因为\(H\)是实对称矩阵的相似矩阵？注意\(H\)不一定对称），且都位于复平面单位圆盘内，最大特征值为1。幺模随机约束使得\(H\)的特征值分布可能比一般的双随机矩阵更具结构性，这影响了环形动力学的收敛速度。
\item
  \textbf{与酉矩阵的关系}：\(H\)保留了\(U\)的``能量守恒''信息在概率意义上的对应。在量子力学中，\(|U_{ij}|^2\)表示从状态\(i\)跃迁到状态\(j\)的概率。在MQR中，这被解释为环形节点\(i\)上的``信息''或``激活''被传递到节点\(j\)的权重。
\end{enumerate}

\hypertarget{ux5e7aux6a21ux968fux673aux6d41ux5f62ux4f5cux4e3aux73afux5f62ux8fdeux63a5ux7684ux4f18ux8d8aux6027}{%
\subsection{2.4
幺模随机流形作为环形连接的优越性}\label{ux5e7aux6a21ux968fux673aux6d41ux5f62ux4f5cux4e3aux73afux5f62ux8fdeux63a5ux7684ux4f18ux8d8aux6027}}

在MQR架构中，选择幺模随机矩阵\(H\)作为环形节点间的固定连接矩阵，而非直接使用酉矩阵\(U\)或无约束矩阵\(W\)，是基于多方面的理论考量：

\hypertarget{ux5185ux5728ux7684ux7a33ux5b9aux6027ux4fddux8bc1}{%
\subsubsection{2.4.1
内在的稳定性保证}\label{ux5185ux5728ux7684ux7a33ux5b9aux6027ux4fddux8bc1}}

如第1章所述，稳定性是环形动力学能够有效弛豫的前提。使用\(H\)作为连接矩阵，其动力学核心步骤为线性变换\(h H^T\)。由于\(H^T\)也是双随机矩阵（\(H\)是双随机则\(H^T\)也是），其谱半径\(\rho(H^T)=1\)。在纯线性迭代\(h^{(k+1)} = h^{(k)} H^T\)中，如果初始分布\(h^{(0)}\)是一个概率分布（即元素非负且和为1），那么迭代后的\(h^{(k)}\)将始终保持为概率分布，且其范数（在\(L_1\)意义下）不变。这提供了在概率
simplex 上的稳定性。

更重要的是，当结合MQR中使用的\textbf{松弛迭代}形式（包含与注入项的凸组合）时，可以证明整个动力学系统是一个收缩映射，从而保证唯一不动点的存在和指数级收敛。幺模随机矩阵\(H\)的谱性质是这一收敛性证明的关键。

\hypertarget{ux4e30ux5bccux7684ux8868ux8fbeux80fdux529b}{%
\subsubsection{2.4.2
丰富的表达能力}\label{ux4e30ux5bccux7684ux8868ux8fbeux80fdux529b}}

虽然幺模随机流形是双随机流形的一个子集，但它仍然是一个相当丰富的空间。通过优化生成\(H\)的酉矩阵\(U\)，我们可以让\(H\)学习到复杂的节点间连接模式。与严格的正交/酉变换相比，\(H\)的模平方操作引入了非线性，破坏了严格的等距性，但换来了对概率流更强的控制能力。\(H\)可以学习到哪些节点之间应存在强连接（概率转移权重高），哪些节点之间应弱连接，从而在环形结构上形成有意义的``信息通路''。

此外，通过引入\textbf{自保持混合}参数\(\beta\)，构造\(H_{eff} = (1-\beta)I + \beta H\)，我们可以在单位矩阵（表示自循环、保持自身状态）和学习的连接矩阵\(H\)之间进行插值。这提供了一种抑制``过度混合''、保持局部信息的机制，同时\(H_{eff}\)仍然是一个双随机矩阵，不破坏稳定性。

\hypertarget{ux4e0eux7269ux7406ux548cux6570ux5b66ux6982ux5ff5ux7684ux6df1ux523bux8054ux7cfb}{%
\subsubsection{2.4.3
与物理和数学概念的深刻联系}\label{ux4e0eux7269ux7406ux548cux6570ux5b66ux6982ux5ff5ux7684ux6df1ux523bux8054ux7cfb}}

幺模随机矩阵的概念桥接了多个领域： -
\textbf{量子计算}：在量子电路中，酉矩阵表示量子门的操作，而测量概率则由酉矩阵元素的模平方给出。MQR的\(H=|U|^2\)可以看作是对量子计算中``测量''步骤的抽象。
-
\textbf{随机过程}：\(H\)作为一个转移矩阵，定义了一个有限状态马尔可夫链。环形的弛豫过程可以视为该马尔可夫链在外部输入（注入项）驱动下达到一个非平衡稳态。
-
\textbf{组合数学与优化}：双随机矩阵与置换矩阵（排列矩阵）密切相关。根据Birkhoff-von
Neumann定理，任何双随机矩阵都是置换矩阵的凸组合。幺模随机矩阵则提供了生成双随机矩阵的一种特殊而结构化的方式。

这种跨学科的连接不仅赋予了MQR架构理论上的美感，也暗示了其潜在机制的普适性和鲁棒性。

\hypertarget{ux6269ux5c55ux590dux6570ux57dfux63a8ux7406ux4e0eux53ccux9149ux5206ux89e3}{%
\subsection{2.5
扩展：复数域推理与双酉分解}\label{ux6269ux5c55ux590dux6570ux57dfux63a8ux7406ux4e0eux53ccux9149ux5206ux89e3}}

MQR架构的灵活性部分体现在其对幺模随机流形基础的扩展上。

\hypertarget{ux590dux6570ux9149ux63a8ux7406}{%
\subsubsection{2.5.1 复数酉推理}\label{ux590dux6570ux9149ux63a8ux7406}}

在基础版本中，环形动力学在实数域运行，使用\(H\)作为连接矩阵。但MQR提供了一个可选模式\texttt{dynamics\_mode=unitary}，允许在复数域进行推理。此时，环形动力学直接使用原始的酉矩阵\(U\)（或其共轭转置\(U^\dagger\)）在复数隐藏状态\(h \in C^N\)上进行传播：
\[h \leftarrow h U^\dagger \quad \text{或结合松弛形式}\]
在读出时，再对最终的复数稳态\(h^*\)取模（或采用其他测量方式）得到实数输出。这种模式允许\textbf{相位信息}参与计算。在复数域中，酉变换的保范性得以完全发挥，同时相位可能编码了振荡、干涉等更精细的动态模式，有可能提升模型对周期性或相干结构的建模能力。

\hypertarget{ux53ccux9149ux77e9ux9635ux89e3ux8026}{%
\subsubsection{2.5.2
双酉矩阵解耦}\label{ux53ccux9149ux77e9ux9635ux89e3ux8026}}

另一个高级特性是\textbf{双酉分解}：\(U_{total} = U_{policy} U_{base}\)。这里，\(U_{base}\)是一个冻结的、预定义或随机初始化的酉矩阵，充当稳定的``世界模型''或基础动力学。\(U_{policy}\)是一个可学习的酉矩阵，充当``策略''或``控制''。总的连接矩阵则为\(H = |U_{policy} U_{base}|^2\)。

这种分解具有多重优势： 1.
\textbf{稳定与可塑性分离}：\(U_{base}\)提供不变的、可靠的基础连接结构，确保动力学的基本稳定性。\(U_{policy}\)则专注于学习适应特定任务所需的调整，其更新不会破坏底层的稳定框架。
2.
\textbf{增量学习与迁移学习}：在不同的任务中，可以共享同一个\(U_{base}\)，只微调\(U_{policy}\)，这类似于预训练-微调范式。
3.
\textbf{理论分析}：分解后，\(H\)的结构可以视为在基础动力学上施加了一个可学习的``调制''。这为分析模型的行为提供了更清晰的视角。

\hypertarget{ux672cux7ae0ux5c0fux7ed3}{%
\subsection{2.6 本章小结}\label{ux672cux7ae0ux5c0fux7ed3}}

本章深入剖析了MöbiusQuantumRing网络所依托的核心数学对象------幺模随机流形。我们从酉矩阵的

\hypertarget{ux7b2c3ux7ae0-ux5e7aux6a21ux968fux673aux6d41ux5f62ux4e0eux73afux5f62ux52a8ux529bux5b66ux7684ux6570ux5b66ux57faux7840}{%
\section{第3章
幺模随机流形与环形动力学的数学基础}\label{ux7b2c3ux7ae0-ux5e7aux6a21ux968fux673aux6d41ux5f62ux4e0eux73afux5f62ux52a8ux529bux5b66ux7684ux6570ux5b66ux57faux7840}}

\hypertarget{ux4eceux9149ux7fa4ux5230ux5e7aux6a21ux968fux673aux77e9ux9635ux6838ux5fc3ux6570ux5b66ux6784ux9020}{%
\subsection{3.1
从酉群到幺模随机矩阵：核心数学构造}\label{ux4eceux9149ux7fa4ux5230ux5e7aux6a21ux968fux673aux77e9ux9635ux6838ux5fc3ux6570ux5b66ux6784ux9020}}

莫比乌斯量子环形网络（Möbius Quantum Ring
Network，简称MQR）的理论基石建立在两个紧密相关的数学概念之上：\textbf{酉群（Unitary
Group）} 与\textbf{幺模随机矩阵（Unistochastic
Matrix）}。理解这两者之间的关系，以及如何将它们转化为可学习的参数化形式，是掌握MQR架构设计哲学的关键。

\hypertarget{ux9149ux77e9ux9635ux4e0eux5e7aux6a21ux968fux673aux77e9ux9635ux7684ux5b9aux4e49ux4e0eux6027ux8d28}{%
\subsubsection{3.1.1
酉矩阵与幺模随机矩阵的定义与性质}\label{ux9149ux77e9ux9635ux4e0eux5e7aux6a21ux968fux673aux77e9ux9635ux7684ux5b9aux4e49ux4e0eux6027ux8d28}}

在复数域 (C) 上，一个 (N \times N) 的矩阵 (U)
被称为\textbf{酉矩阵（Unitary Matrix）}，当且仅当它满足条件：

{[} U\^{}\dagger U = U U\^{}\dagger = I\_N {]}

其中 (U\^{}\dagger) 表示 (U) 的共轭转置（Hermitian conjugate），(I\_N)
是 (N \times N)
的单位矩阵。酉矩阵的核心性质在于其\textbf{保范性（Norm-Preserving）}：对于任意复向量
(v \in C\^{}N)，有 (\textbar Uv\textbar\_2 =
\textbar v\textbar\_2)。这意味着酉变换是一个等距映射，不会放大或缩小输入信号的能量，从根源上避免了线性部分导致的梯度爆炸或消失问题，这正是第1章所讨论的稳定性问题的理想解决方案。

然而，直接将酉矩阵用作循环神经网络的权重矩阵面临着优化复杂性和表达能力受限的挑战。MQR采用了一个巧妙的间接策略：从酉矩阵派生出一种具有特殊概率意义的矩阵------幺模随机矩阵。

给定一个酉矩阵
(U)，我们通过对其每个元素的模（绝对值）进行平方运算，构造一个新的矩阵
(H)：

{[} H = \textbar U\textbar\^{}2, \quad \text{即} \quad H\_\{ij\} =
\textbar U\_\{ij\}\textbar\^{}2 {]}

这样得到的矩阵 (H) 具有两个至关重要的概率性质： 1.
\textbf{双随机性（Doubly Stochastic）}：矩阵 (H)
的所有元素均为非负实数，且每一行之和与每一列之和都等于1。即 (\sum\emph{j
H}\{ij\} = 1) 且 (\sum\emph{i H}\{ij\} = 1)。 2.
\textbf{幺模性（Unistochastic）}：该矩阵被明确地定义为一个酉矩阵元素的模平方。所有由这种方式生成的矩阵构成的集合，称为\textbf{幺模随机流形（Unistochastic
Manifold）}。

双随机矩阵在概率论和马尔可夫链理论中具有明确的物理解释：它可以被视作一个\textbf{有限状态、离散时间的齐次马尔可夫链的转移概率矩阵}。矩阵元素
(H\_\{ij\}) 表示从状态 (i) 转移到状态 (j)
的概率。双随机性意味着该马尔可夫链是\textbf{保测度（measure-preserving）}
的，即一个均匀分布（所有状态概率均为
(1/N)）是其平稳分布。这为环形网络中的信息传播提供了一种``能量守恒''或``概率质量守恒''的直观解释。

幺模随机性是比普通双随机性更强的约束。并非所有双随机矩阵都能表示为某个酉矩阵的模平方。幺模随机矩阵构成的集合是双随机矩阵凸集（Birkhoff
polytope）中的一个非凸子流形。这种更强的约束将酉群的几何与对称性``烙印''到了概率转移矩阵上，使得
(H)
不仅具有概率意义，还继承了酉变换部分优良的数学特性，为构建具有理论保证的稳定动力学奠定了基础。

\hypertarget{cayleyux53d8ux6362ux9149ux77e9ux9635ux7684ux5b9eux7528ux53c2ux6570ux5316}{%
\subsubsection{3.1.2
Cayley变换：酉矩阵的实用参数化}\label{cayleyux53d8ux6362ux9149ux77e9ux9635ux7684ux5b9eux7528ux53c2ux6570ux5316}}

为了在神经网络中有效地学习和优化酉矩阵
(U)，我们需要一种能够平滑参数化酉流形（即酉群
(U(N))）的方法。直接使用矩阵元素作为无约束参数进行优化，然后通过施密特正交化等后处理手段强制其满足酉性，会导致优化路径不连续且训练困难。MQR采用\textbf{Cayley变换}作为核心参数化工具，这是一种在流形优化中广泛使用的技术。

Cayley变换提供了一种将斜埃尔米特矩阵（Skew-Hermitian
Matrix）映射到酉矩阵的可逆方法。一个矩阵 (A) 是斜埃尔米特的，如果满足
(A\^{}\dagger = -A)。对于任意斜埃尔米特矩阵
(A)，以下变换定义了一个酉矩阵：

{[} U = (I + A)(I - A)\^{}\{-1\} {]}

其逆变换为 (A = (U - I)(U +
I)\^{}\{-1\})。在实现中，为了确保数值稳定性和可逆性，通常使用以下形式：

{[} U = (I + A/2)(I - A/2)\^{}\{-1\} {]}

关键点在于，斜埃尔米特矩阵 (A)
可以\textbf{自由地、无约束地参数化}。具体而言，我们可以取一个任意形状的实矩阵
(W \in R\^{}\{N \times N\})，然后构造斜埃尔米特矩阵 (A) 为：

{[} A = W - W\^{}T + i \cdot (W + W\^{}T) {]}

这里 (i) 是虚数单位。这样，通过一个无约束的实参数矩阵
(W)，我们就能生成对应的斜埃尔米特矩阵
(A)，进而通过Cayley变换得到严格满足酉性条件的矩阵
(U)。在训练过程中，优化器直接更新实参数矩阵
(W)，梯度通过Cayley变换的反向传播计算。这种方法确保了在整个优化过程中，矩阵
(U) 始终保持在酉流形上，从而保障了模型前向传播的稳定性。

\begin{figure}
\centering
\includegraphics{figures/diagram_3.png}
\caption{图 3}
\end{figure}

\emph{图3.1:
从无约束参数到幺模随机连接矩阵的构造流程。通过Cayley变换将自由参数W映射到酉矩阵U，再通过模平方运算得到具有概率解释的双随机矩阵H，作为环形动力学的基础。}

\hypertarget{ux73afux5f62ux52a8ux529bux5b66ux56faux5b9aux70b9ux677eux5f1bux4e0eux6536ux655bux6027ux5206ux6790}{%
\subsection{3.2
环形动力学：固定点松弛与收敛性分析}\label{ux73afux5f62ux52a8ux529bux5b66ux56faux5b9aux70b9ux677eux5f1bux4e0eux6536ux655bux6027ux5206ux6790}}

MQR的核心计算范式突破了传统前馈或循环网络的顺序计算图模式，代之以一种基于\textbf{固定点松弛（Fixed-Point
Relaxation）}
的动力学过程。该过程模拟了一个离散动力系统如何在外界驱动下逐渐弛豫到一个唯一的平衡态。

\hypertarget{ux57faux672cux52a8ux529bux5b66ux65b9ux7a0b}{%
\subsubsection{3.2.1
基本动力学方程}\label{ux57faux672cux52a8ux529bux5b66ux65b9ux7a0b}}

设环形网络有 (N) 个节点，其状态用一个行向量 (h \in R\^{}\{1
\times N\}) 表示（为简化，先考虑实值情况）。网络的连接由幺模随机矩阵 (H
\in R\^{}\{N \times N\})
定义。在推理的每个``时间步''（或称松弛迭代步），网络状态根据以下方程更新：

{[} h\^{}\{(t+1)\} = h\^{}\{(t)\} H\^{}T {]}

这是一个线性马尔可夫过程。由于 (H) 是双随机的，其转置 (H\^{}T)
也是双随机的。根据马尔可夫链理论，如果该链是不可约且非周期的（对于由学习得到的
(H)，这通常成立），那么无论初始状态 (h\^{}\{(0)\})
如何，在多次迭代后，状态向量都会收敛到一个唯一的\textbf{平稳分布（Stationary
Distribution）} (\pi)，满足：

{[} \pi = \pi H\^{}T, \quad \sum\_i \pi\_i = 1 {]}

对于双随机矩阵，这个平稳分布就是均匀分布：(\pi = (\frac{1}{N},
\frac{1}{N}, \ldots,
\frac{1}{N}))。这意味着，如果没有任何外部输入，环形网络中的信息最终会均匀扩散到所有节点，丢失所有初始信息。这显然不是一个有用的计算模型。

\hypertarget{ux6ce8ux5165ux9a71ux52a8ux4e0eux963bux5c3cux677eux5f1b}{%
\subsubsection{3.2.2
注入驱动与阻尼松弛}\label{ux6ce8ux5165ux9a71ux52a8ux4e0eux963bux5c3cux677eux5f1b}}

为了使网络能够对外部输入做出响应并计算有意义的表示，MQR引入了\textbf{注入算子（Injection
Operator）} (J(x))，其作用是将输入数据 (x)
编码并``注入''到环形网络中。同时，为了避免状态完全收敛到均匀分布，在动力学方程中加入了\textbf{阻尼（Damping）}
因子 (\alpha \in (0, 1{]})。完整的、带驱动的环形动力学方程如下：

{[} h\^{}\{(t+1)\} = (1 - \alpha) h\^{}\{(t)\} H\^{}T +
\alpha J(x) {]}

其中 (J(x) \in R\^{}\{1
\times N\})。这个方程可以理解为：在每一步，网络状态以 ((1-\alpha))
的权重按照内部连接 (H\^{}T) 进行扩散（马尔可夫更新），同时以 (\alpha)
的权重被外部输入 (J(x)) 所``刷新''或``驱动''。

这个方程定义了一个\textbf{线性时不变（LTI）离散动力系统}。我们可以将其重写为：

{[} h\^{}\{(t+1)\} = h\^{}\{(t)\} \left[(1-\alpha) H^T\right] +
\alpha J(x) {]}

令 (M = (1-\alpha) H\^{}T)，则方程变为 (h\^{}\{(t+1)\} = h\^{}\{(t)\} M
+ \alpha J(x))。这是一个标准的线性动力系统，其长期行为由矩阵
(M) 的谱性质决定。

\hypertarget{ux6536ux655bux6027ux8bc1ux660e}{%
\subsubsection{3.2.3 收敛性证明}\label{ux6536ux655bux6027ux8bc1ux660e}}

\textbf{定理 3.1 (环形动力学收敛性)}：对于阻尼因子 (0 \textless{}
\alpha \leq 1) 和双随机矩阵 (H)，由方程 (h\^{}\{(t+1)\} = (1 - \alpha)
h\^{}\{(t)\} H\^{}T + \alpha J(x))
定义的动力学过程，对于任意初始状态
(h\^{}\{(0)\})，都会指数收敛到一个唯一的\textbf{固定点（Fixed Point）}
(h\^{}*)，且该固定点与 (h\^{}\{(0)\}) 无关。

\textbf{证明}： 1. \textbf{压缩映射性质}：首先证明更新映射 (f(h) =
(1-\alpha) h H\^{}T + \alpha J(x))
在合适的度量下是一个压缩映射。考虑两个状态 (h\_a) 和
(h\_b)，它们的差经过一次映射后为： {[} f(h\_a) - f(h\_b) = (1-\alpha)
(h\_a - h\_b) H\^{}T {]}
对上述两边取某种范数（例如，总变差范数或与平稳分布相关的范数）。由于
(H\^{}T) 是一个随机矩阵（行和为1），其诱导的算子范数满足
(\textbar H\^{}T\textbar\_p \leq 1)（对于 (p=1) 或 (p=\infty)
范数）。因此， {[} \textbar f(h\_a) - f(h\_b)\textbar{} \leq (1-\alpha)
\textbar h\_a - h\_b\textbar{} \textbar H\^{}T\textbar{} \leq (1-\alpha)
\textbar h\_a - h\_b\textbar{} {]} 因为 (0 \textless{}
\alpha \leq 1)，所以 (0 \leq (1-\alpha) \textless{} 1)。这意味着 (f)
是一个 Lipschitz 常数为 ((1-\alpha)) 的压缩映射。

\begin{enumerate}
\def\labelenumi{\arabic{enumi}.}
\setcounter{enumi}{1}
\item
  \textbf{巴拿赫不动点定理}：根据巴拿赫不动点定理，在一个完备度量空间（此处为
  (R\^{}N) 配备上述范数）中，任何压缩映射都存在唯一的不动点
  (h\^{}\emph{)，满足 (h\^{}} = f(h\^{}*))，并且通过迭代 (h\^{}\{(t+1)\}
  = f(h\^{}\{(t)\})) 可以从任意起点指数收敛到该不动点。
\item
  \textbf{不动点显式解}：不动点 (h\^{}\emph{) 满足方程： {[} h\^{}} =
  (1-\alpha) h\textsuperscript{* H}T + \alpha J(x) {]}
  整理得： {[} h\^{}* (I - (1-\alpha) H\^{}T) = \alpha J(x)
  {]} 由于 (M = (1-\alpha)H\^{}T) 的谱半径小于1（因为
  (\alpha\textgreater0)），矩阵 ((I - M))
  是可逆的。因此，不动点有显式表达式： {[} h\^{}* =
  \alpha J(x) (I - (1-\alpha) H\textsuperscript{T)}\{-1\} {]}
  这个表达式虽然涉及矩阵求逆，在实际计算中并不直接使用，而是通过迭代松弛来逼近。但它从理论上确认了不动点的存在性、唯一性以及对输入
  (J(x)) 和矩阵 (H) 的依赖性。
\end{enumerate}

这个收敛性定理是MQR架构的理论支柱。它保证了无论网络初始化和迭代步数如何，前向推理过程总能稳定地收敛到一个确定的状态
(h\^{}\emph{)。这不仅确保了数值稳定性，还意味着我们可以将整个松弛过程视为一个``黑盒''，其输入是
(J(x)) 和 (H)，输出是
(h\^{}})。在反向传播时，我们可以利用隐函数定理或自动微分框架，直接对固定点
(h\^{}*) 求梯度，而不需要展开所有迭代步骤，从而实现高效且稳定的训练。

\hypertarget{ux975eux7ebfux6027ux6269ux5c55ux4e0eux4fddux6536ux655bux6027ux6fc0ux6d3b}{%
\subsubsection{3.2.4
非线性扩展与保收敛性激活}\label{ux975eux7ebfux6027ux6269ux5c55ux4e0eux4fddux6536ux655bux6027ux6fc0ux6d3b}}

基本的动力学方程是线性的。为了引入非线性表达能力，MQR允许在特定位置插入激活函数，同时精心设计以确保收敛性不被破坏。

\begin{enumerate}
\def\labelenumi{\arabic{enumi}.}
\item
  \textbf{注入端非线性}：注入算子 (J(x))
  本身可以是非线性的。最常见的形式是采用类似LoRA的低秩结构加激活：(J(x)
  = W\_\{up\} \cdot \sigma(W\_\{down\} \cdot x))，其中 (\sigma)
  是ReLU、GELU等常见激活函数。这发生在动力学方程外部，不影响内部松弛过程的收敛性。
\item
  \textbf{环内松弛非线性（1-Lipschitz
  激活）}：更激进的做法是在状态更新方程内部引入非线性： {[}
  h\^{}\{(t+1)\} = \sigma\left( (1-\alpha) h\^{}\{(t)\} H\^{}T +
  \alpha J(x) \right) {]}
  为了保持压缩映射性质，从而保证收敛，要求激活函数 (\sigma)
  是\textbf{1-Lipschitz}的。即对于任意输入 (a, b)，有
  (\textbar{}\sigma(a) - \sigma(b)\textbar{} \leq \textbar a -
  b\textbar)。许多常用的激活函数满足或可以调整为满足此条件，例如：

  \begin{itemize}
  \tightlist
  \item
    \textbf{双曲正切（tanh）}：导数值域在(0,1{]}，是1-Lipschitz的。
  \item
    \textbf{逐点缩放后的ReLU}：标准ReLU是1-Lipschitz的。
  \item
    \textbf{GroupSort、Householder激活等}：专门为构建利普希茨网络设计的激活函数。
    当使用1-Lipschitz激活时，更新映射 (f(h))
    的压缩性论证依然成立，因为复合一个1-Lipschitz函数不会增加Lipschitz常数。因此，收敛性定理得以保持。
  \end{itemize}
\end{enumerate}

\hypertarget{ux6709ux6548ux6df7ux5408ux4e0eux81eaux4fddux6301ux8fdeux63a5}{%
\subsection{3.3
有效混合与自保持连接}\label{ux6709ux6548ux6df7ux5408ux4e0eux81eaux4fddux6301ux8fdeux63a5}}

在基本动力学中，连接矩阵 (H)
是完全双随机的。这意味着在无驱动情况下，系统会趋向于均匀分布。即使有驱动，强烈的扩散效应也可能导致从不同输入产生的状态
(h\^{}*)
彼此过于相似，降低了模型的区分能力。这种现象可类比于过度平滑（over-smoothing）在图神经网络中的问题。

为了解决这个问题，MQR引入了\textbf{有效混合（Effective Mixing）}
机制，通过一个参数 (\beta \in [0, 1])
来调节``自保持（Self-Retention）''与``邻居混合（Neighbor
Mixing）''的平衡。具体做法是构造一个\textbf{凸组合（Convex
Combination）} 矩阵 (H\_\{\text{eff}\})：

{[} H\_\{\text{eff}\} = (1 - \beta) I + \beta H {]}

其中 (I) 是单位矩阵。然后将动力学方程中的 (H) 替换为
(H\_\{\text{eff}\})：

{[} h\^{}\{(t+1)\} = (1 - \alpha) h\^{}\{(t)\} H\_\{\text{eff}\}\^{}T +
\alpha J(x) {]}

\textbf{性质分析}： - \textbf{保持双随机性}：由于 (I) 和 (H)
都是双随机的，且 ((1-\beta)) 和 (\beta) 为非负且和为1，因此
(H\_\{\text{eff}\})
也是双随机的。这意味着所有关于收敛性的理论保证依然有效。 -
\textbf{调节扩散速度}：当 (\beta = 1) 时，(H\_\{\text{eff}\} =
H)，恢复原始的全混合模式。当 (\beta = 0) 时，(H\_\{\text{eff}\} =
I)，动力学方程变为 (h\^{}\{(t+1)\} = (1-\alpha) h\^{}\{(t)\} +
\alpha J(x))，这是一个简单的一阶低通滤波器，每个节点独立地记忆自己的注入历史，节点间无信息交换。通过调节
(\beta)，可以控制信息在环形网络中扩散的速率和范围。 -
\textbf{谱性质}：(H\_\{\text{eff}\}) 的特征值是 (H)
的特征值的缩放和平移。具体地，如果 (H) 的特征值为 (\lambda\emph{i)，则
(H}\{\text{eff}\}) 的特征值为 ((1-\beta) +
\beta \lambda\emph{i)。由于双随机矩阵的主特征值为1，对应的
(H}\{\text{eff}\}) 特征值仍为1。其他特征值则从 (\lambda\_i)
向1方向收缩，这通常会使系统更快地收敛到与驱动相关的固定点，同时保留必要的混合能力。

这种自保持混合机制为模型提供了一个重要的超参数，使其能够适应不同任务对局部性与全局性信息整合的需求。例如，对于需要强局部特征保存的任务（如图像分类），可以使用较小的
(\beta)；对于需要全局上下文充分混合的任务（如语言建模），则可以使用较大的
(\beta)。

\hypertarget{ux590dux6570ux57dfux6269ux5c55ux9149ux63a8ux7406ux4e0eux76f8ux4f4dux4fe1ux606f}{%
\subsection{3.4
复数域扩展：酉推理与相位信息}\label{ux590dux6570ux57dfux6269ux5c55ux9149ux63a8ux7406ux4e0eux76f8ux4f4dux4fe1ux606f}}

前述

\hypertarget{ux7b2c4ux7ae0-ux5e7aux6a21ux968fux673aux6d41ux5f62ux4e0eux73afux5f62ux52a8ux529bux5b66ux7406ux8bbaux57faux7840ux4e0eux6536ux655bux6027ux8bc1ux660e}{%
\section{第4章
幺模随机流形与环形动力学：理论基础与收敛性证明}\label{ux7b2c4ux7ae0-ux5e7aux6a21ux968fux673aux6d41ux5f62ux4e0eux73afux5f62ux52a8ux529bux5b66ux7406ux8bbaux57faux7840ux4e0eux6536ux655bux6027ux8bc1ux660e}}

\hypertarget{ux5e7aux6a21ux968fux673aux77e9ux9635ux4eceux9149ux7fa4ux5230ux53ccux968fux673aux6d41ux5f62}{%
\subsection{4.1
幺模随机矩阵：从酉群到双随机流形}\label{ux5e7aux6a21ux968fux673aux77e9ux9635ux4eceux9149ux7fa4ux5230ux53ccux968fux673aux6d41ux5f62}}

莫比乌斯量子环形网络（MQR）的核心数学基础建立在\textbf{幺模随机矩阵}这一概念之上。要理解其重要性，我们需要从更基础的群论和矩阵分析开始。

\hypertarget{ux9149ux77e9ux9635ux4e0eux53ccux968fux673aux77e9ux9635}{%
\subsubsection{4.1.1
酉矩阵与双随机矩阵}\label{ux9149ux77e9ux9635ux4e0eux53ccux968fux673aux77e9ux9635}}

在复向量空间中，一个\textbf{酉矩阵} ( U \in C\^{}\{N \times N\}
) 满足条件 ( U\^{}\dagger U = UU\^{}\dagger = I\_N )，其中 ( \dagger )
表示共轭转置。酉矩阵构成的集合 ( U(N) )
形成一个紧致李群，称为\textbf{酉群}。酉矩阵作用于向量时，保持向量的2-范数（欧几里得范数）不变，即对于任意向量
( v \in C\^{}N )，有 ( \textbar Uv\textbar\_2 =
\textbar v\textbar\_2
)。这一等距性质是正交/酉循环神经网络稳定性的理论基础。

另一方面，一个\textbf{双随机矩阵} ( B \in R\^{}\{N \times N\} )
满足两个条件： 1. 所有元素非负：( B\_\{ij\} \geq 0 ) 2.
行和与列和均为1：( \sum\emph{\{j=1\}\^{}N B}\{ij\} = 1 ) 且 (
\sum\emph{\{i=1\}\^{}N B}\{ij\} = 1 )

双随机矩阵可以看作是一个概率转移矩阵，描述了一个马尔可夫链的状态转移概率。根据Birkhoff-von
Neumann定理，任何双随机矩阵都可以表示为置换矩阵的凸组合。

\hypertarget{ux5e7aux6a21ux968fux673aux77e9ux9635ux7684ux5b9aux4e49ux4e0eux6027ux8d28}{%
\subsubsection{4.1.2
幺模随机矩阵的定义与性质}\label{ux5e7aux6a21ux968fux673aux77e9ux9635ux7684ux5b9aux4e49ux4e0eux6027ux8d28}}

\textbf{幺模随机矩阵}（Unistochastic
Matrix）是这两类矩阵的桥梁。给定一个酉矩阵 ( U \in U(N)
)，通过逐元素取模平方运算，我们定义对应的幺模随机矩阵：

{[} H = \textbar U\textbar\^{}2, \quad \text{其中} \quad H\_\{ij\} =
\textbar U\_\{ij\}\textbar\^{}2 {]}

幺模随机矩阵具有以下关键性质：

\begin{enumerate}
\def\labelenumi{\arabic{enumi}.}
\item
  \textbf{天然双随机性}：由于酉矩阵的行和列满足 ( \sum\emph{\{j=1\}\^{}N
  \textbar U}\{ij\}\textbar\^{}2 = 1 ) 且 ( \sum\emph{\{i=1\}\^{}N
  \textbar U}\{ij\}\textbar\^{}2 = 1 )，因此 ( H )
  自动满足双随机矩阵的条件。
\item
  \textbf{能量守恒}：对于任意概率分布向量 ( p \in R\^{}N
  )（满足 ( p\_i \geq 0, \sum\_i p\_i = 1 )），有 ( \sum\_i (Hp)\_i =
  \sum\_i p\_i = 1 )。这意味着在马尔可夫链的视角下，总概率质量保持不变。
\item
  \textbf{谱性质}：幺模随机矩阵的谱半径 ( \rho(H) = 1
  )，且1是其最大特征值，对应的特征向量是均匀分布向量 ( \textbf{1} =
  (1,1,\ldots,1)\^{}T/\sqrt{N} )。
\item
  \textbf{参数化优势}：通过酉矩阵参数化双随机矩阵，我们将优化问题约束在一个光滑的流形上，避免了直接优化双随机矩阵时面临的非负性和求和约束。
\end{enumerate}

幺模随机矩阵的集合 ( U(N) ) 是双随机矩阵集合 ( B(N)
) 的一个真子集。对于 ( N \leq 3 )，有 ( U(N) = B(N)
)，但对于 ( N \geq 4
)，幺模随机矩阵只是双随机矩阵的一个子集。这一事实意味着幺模随机流形提供了一个结构化的、可有效参数化的双随机矩阵子空间。

\hypertarget{cayleyux53c2ux6570ux5316ux5728stiefelux6d41ux5f62ux4e0aux7684ux6709ux6548ux4f18ux5316}{%
\subsection{4.2
Cayley参数化：在Stiefel流形上的有效优化}\label{cayleyux53c2ux6570ux5316ux5728stiefelux6d41ux5f62ux4e0aux7684ux6709ux6548ux4f18ux5316}}

\hypertarget{cayleyux53d8ux6362ux7684ux539fux7406}{%
\subsubsection{4.2.1
Cayley变换的原理}\label{cayleyux53d8ux6362ux7684ux539fux7406}}

在MQR中，我们采用\textbf{Cayley变换}来参数化酉矩阵。Cayley变换提供了一种将斜埃尔米特矩阵（skew-Hermitian
matrix）映射到酉矩阵的优雅方法。

给定一个斜埃尔米特矩阵 ( A \in C\^{}\{N \times N\} )，满足 (
A\^{}\dagger = -A )，Cayley变换定义为：

{[} U = (I + A)(I - A)\^{}\{-1\} {]}

可以验证，这样定义的 ( U ) 确实是酉矩阵： {[} U\^{}\dagger U = {[}(I -
A)\^{}\{-1\}{]}\^{}\dagger (I + A)\^{}\dagger (I + A)(I - A)\^{}\{-1\} =
I {]}

逆变换为： {[} A = (U - I)(U + I)\^{}\{-1\} {]}

\hypertarget{ux5b9eux6570ux57dfux7684ux7279ux5316ux4e0eux5b9eux73b0}{%
\subsubsection{4.2.2
实数域的特化与实现}\label{ux5b9eux6570ux57dfux7684ux7279ux5316ux4e0eux5b9eux73b0}}

在实际实现中，MQR主要处理实数域的情况，此时酉矩阵退化为正交矩阵，斜埃尔米特矩阵退化为斜对称矩阵（skew-symmetric
matrix）。设 ( S \in R\^{}\{N \times N\} ) 为斜对称矩阵（(
S\^{}T = -S )），则Cayley变换简化为：

{[} Q = (I + S)(I - S)\^{}\{-1\} {]}

其中 ( Q ) 是正交矩阵（( Q\^{}T Q = I )）。

斜对称矩阵的自由度仅为 ( \frac{N(N-1)}{2} )，远小于一般矩阵的 ( N\^{}2 )
个自由度。我们可以用一个自由参数矩阵 ( W \in R\^{}\{N
\times N\} ) 来生成斜对称矩阵：

{[} S = W - W\^{}T {]}

这样，我们只需要优化自由参数 ( W )，即可通过Cayley变换得到正交矩阵 ( Q
)，进而得到双随机矩阵 ( H = \textbar Q\textbar\^{}2 = Q \odot Q
)（在实数域中，模平方退化为逐元素平方）。

\hypertarget{ux6570ux503cux7a33ux5b9aux6027ux7684ux4fddux969c}{%
\subsubsection{4.2.3
数值稳定性的保障}\label{ux6570ux503cux7a33ux5b9aux6027ux7684ux4fddux969c}}

Cayley变换在数值计算中具有良好性质。当 ( S ) 的范数较小时，( (I - S) )
接近单位矩阵，其逆矩阵可以高效稳定计算。在实际实现中，MQR采用以下策略确保数值稳定性：

\begin{enumerate}
\def\labelenumi{\arabic{enumi}.}
\item
  \textbf{缩放控制}：对生成的斜对称矩阵 ( S )
  进行适当的缩放，确保其谱范数小于1，避免 ( (I - S) ) 接近奇异。
\item
  \textbf{迭代求逆}：对于大规模矩阵，采用迭代法或Cholesky分解求解线性系统
  ( (I - S)x = b )，而非显式计算逆矩阵。
\item
  \textbf{梯度计算}：通过自动微分框架（如PyTorch）隐式计算Cayley变换的梯度，避免了手动推导复杂梯度公式可能引入的数值误差。
\end{enumerate}

\begin{figure}
\centering
\includegraphics{figures/diagram_4.png}
\caption{图 4}
\end{figure}

上图展示了MQR中从参数化到推理的完整流程，突出了幺模随机矩阵的核心地位。

\hypertarget{ux73afux5f62ux52a8ux529bux5b66ux7cfbux7edfux56faux5b9aux70b9ux677eux5f1bux4e0eux6536ux655bux6027ux5206ux6790}{%
\subsection{4.3
环形动力学系统：固定点松弛与收敛性分析}\label{ux73afux5f62ux52a8ux529bux5b66ux7cfbux7edfux56faux5b9aux70b9ux677eux5f1bux4e0eux6536ux655bux6027ux5206ux6790}}

\hypertarget{ux57faux672cux52a8ux529bux5b66ux65b9ux7a0b-1}{%
\subsubsection{4.3.1
基本动力学方程}\label{ux57faux672cux52a8ux529bux5b66ux65b9ux7a0b-1}}

MQR摒弃了传统RNN的序列展开范式，采用了一种\textbf{环形动力学系统}。系统的状态由环上的
( N ) 个节点表示，每个节点对应一个实数值 ( h\_i \in R
)。整个环的状态向量记为 ( h \in R\^{}N )。

系统的动力学由以下离散时间方程描述：

{[} h\^{}\{(t+1)\} = (1-\alpha)h\textsuperscript{\{(t)\}H}T +
\alphaJ(x) {]}

其中： - ( h\^{}\{(t)\} ) 是第 ( t ) 次迭代时的环状态 - ( H
\in R\^{}\{N \times N\} ) 是幺模随机连接矩阵（( H =
\textbar U\textbar\^{}2 )） - ( \alpha \in (0,1{]} )
是弛豫率，控制外部注入与内部传播的平衡 - ( J(x):
R\^{}\{d\_\{in\}\} \to R\^{}N ) 是注入算子，将输入 ( x
) 映射为环上的驱动信号

\hypertarget{ux56faux5b9aux70b9ux5b58ux5728ux6027ux4e0eux552fux4e00ux6027}{%
\subsubsection{4.3.2
固定点存在性与唯一性}\label{ux56faux5b9aux70b9ux5b58ux5728ux6027ux4e0eux552fux4e00ux6027}}

\textbf{定理4.1（固定点存在性与唯一性）}：对于任意输入 ( x )
和任意初始状态 ( h\^{}\{(0)\} )，上述动力学系统存在唯一全局吸引固定点 (
h\^{}* )，满足：

{[} h\^{}* = (1-\alpha)h\^{}*H\^{}T + \alphaJ(x) {]}

\textbf{证明}： 将动力学方程重写为： {[} h\^{}\{(t+1)\} =
T(h\^{}\{(t)\}) {]} 其中 ( T(h) = (1-\alpha)hH\^{}T +
\alphaJ(x) )。

我们需要证明 ( T ) 是一个收缩映射。考虑任意两个状态向量 ( h, g
\in R\^{}N )： {[}

\begin{aligned}
\|T(h) - T(g)\|_1 &= \|(1-\alpha)(h-g)H^T\|_1 \\
&\leq (1-\alpha)\|h-g\|_1 \cdot \|H^T\|_1
\end{aligned}

{]}

对于双随机矩阵 ( H )，其1-范数（列和的最大值）为1，因此： {[}
\textbar T(h) - T(g)\textbar\_1 \leq (1-\alpha)\textbar h-g\textbar\_1
{]}

由于 ( \alpha \in (0,1{]} )，有 ( 0 \leq 1-\alpha \textless{} 1 )，故 (
T ) 是收缩系数为 ( (1-\alpha) ) 的收缩映射。根据Banach不动点定理，( T )
在完备度量空间 ( (R\^{}N, \textbar{}\cdot\textbar\_1) )
上存在唯一不动点 ( h\^{}*
)，且从任意初始点出发的迭代都会指数收敛到该不动点。

收敛速率由 ( (1-\alpha) ) 控制：( \alpha )
越大，收敛越快，但系统对外部注入的响应也越强；( \alpha )
越小，收敛越慢，但系统更依赖于内部传播动力学。

\hypertarget{ux7a33ux6001ux7684ux663eux5f0fux89e3}{%
\subsubsection{4.3.3
稳态的显式解}\label{ux7a33ux6001ux7684ux663eux5f0fux89e3}}

不动点方程可以显式求解。将方程重排： {[} h\^{}* -
(1-\alpha)h\^{}\emph{H\^{}T = \alphaJ(x) {]} {[} h\^{}}(I -
(1-\alpha)H\^{}T) = \alphaJ(x) {]}

由于 ( \rho((1-\alpha)H\^{}T) = 1-\alpha \textless{} 1 )，矩阵 ( I -
(1-\alpha)H\^{}T ) 可逆，其逆可表示为Neumann级数： {[} (I -
(1-\alpha)H\textsuperscript{T)}\{-1\} = \sum\_\{k=0\}\^{}\{\infty\}
(1-\alpha)\^{}k (H\textsuperscript{T)}k {]}

因此，稳态的显式解为： {[} h\^{}* = \alphaJ(x)
\sum\_\{k=0\}\^{}\{\infty\} (1-\alpha)\^{}k (H\textsuperscript{T)}k {]}

这个解具有清晰的物理解释：稳态是无限次传播的加权和，每次传播后信号被衰减因子
( (1-\alpha) ) 缩放。这类似于一个无限冲激响应（IIR）滤波器。

\hypertarget{ux80fdux91cfux5b88ux6052ux4e0eux6982ux7387ux89e3ux91ca}{%
\subsubsection{4.3.4
能量守恒与概率解释}\label{ux80fdux91cfux5b88ux6052ux4e0eux6982ux7387ux89e3ux91ca}}

当注入信号 ( J(x) )
是概率分布（非负且和为1）时，环状态在整个动力学过程中保持概率分布的性质。

\textbf{命题4.2（概率守恒）}：如果 ( J(x) ) 是概率向量且 (
h\^{}\{(0)\} ) 是概率向量，则对所有 ( t \geq 0 )，( h\^{}\{(t)\} )
都是概率向量。

\textbf{证明}：由数学归纳法。假设 ( h\^{}\{(t)\} )
是概率向量（非负且和为1），则： 1. 非负性：由于 ( H\^{}T ) 的元素非负且
( J(x) ) 非负，因此 ( h\^{}\{(t+1)\} )
是非负组合，保持非负性。 2. 和为一：计算分量和： {[}

\begin{aligned}
   \sum_i h_i^{(t+1)} &= (1-\alpha)\sum_i \sum_j h_j^{(t)} H_{ji} + \alpha\sum_i J_i(x) \\
   &= (1-\alpha)\sum_j h_j^{(t)} \underbrace{\sum_i H_{ji}}_{=1} + \alpha\underbrace{\sum_i J_i(x)}_{=1} \\
   &= (1-\alpha) \cdot 1 + \alpha \cdot 1 = 1
   \end{aligned}

{]} 其中利用了 ( H ) 的双随机性（行和为1）和 ( J(x) )
的概率性质。

这一性质使得MQR的环状态可以解释为概率分布，为后续的采样读出提供了自然基础。

\hypertarget{ux975eux7ebfux6027ux6269ux5c55ux4e0elipschitzux8fdeux7eedux6027ux4fddux6301}{%
\subsection{4.4
非线性扩展与Lipschitz连续性保持}\label{ux975eux7ebfux6027ux6269ux5c55ux4e0elipschitzux8fdeux7eedux6027ux4fddux6301}}

\hypertarget{ux6ce8ux5165ux7aefux7684ux975eux7ebfux6027}{%
\subsubsection{4.4.1
注入端的非线性}\label{ux6ce8ux5165ux7aefux7684ux975eux7ebfux6027}}

基本动力学方程中的注入算子 ( J(x) )
可以引入非线性。MQR支持两种形式的注入算子：

\begin{enumerate}
\def\labelenumi{\arabic{enumi}.}
\tightlist
\item
  \textbf{线性注入}：( J(x) = W\_\{up\}W\_\{down\}x
  )，这是典型的LoRA式低秩注入。
\item
  \textbf{非线性注入}：( J(x) = W\_\{up\},\sigma(W\_\{down\}x)
  )，其中 ( \sigma ) 是激活函数（如ReLU、GELU等）。
\end{enumerate}

非线性注入增加了模型的表达能力，但需要注意确保注入输出的范围与环状态的解释相容（如保持非负性用于概率解释）。

\hypertarget{ux73afux5185ux677eux5f1bux7684ux975eux7ebfux6027}{%
\subsubsection{4.4.2
环内松弛的非线性}\label{ux73afux5185ux677eux5f1bux7684ux975eux7ebfux6027}}

更一般地，MQR允许在环内松弛过程中引入非线性：

{[} h\^{}\{(t+1)\} = \sigma\left((1-\alpha)h\textsuperscript{\{(t)\}H}T
+ \alphaJ(x)\right) {]}

其中 ( \sigma: R \to R )
是逐元素应用的激活函数。为了保持收敛性，我们需要对 ( \sigma ) 施加约束。

\textbf{定理4.3（带非线性激活的收敛性）}：如果激活函数 ( \sigma )
是1-Lipschitz连续的，即对于所有 ( a,b \in R )，有 (
\textbar{}\sigma(a) - \sigma(b)\textbar{} \leq \textbar a-b\textbar{}
)，且 ( \sigma(0) = 0 )，则带非线性的动力学系统仍收敛到唯一固定点。

\textbf{证明}：定义非线性算子 ( T\_\sigma(h) =
\sigma\left((1-\alpha)hH\^{}T + \alphaJ(x)\right) )。对于任意
( h,g \in R\^{}N )： {[}

\begin{aligned}
\|T_\sigma(h) - T_\sigma(g)\|_1 &= \|\sigma(Ah + b) - \sigma(Ag + b)\|_1 \\
&\leq \|Ah - Ag\|_1 \quad (\text{由 } \sigma \text{ 的1-Lipschitz性}) \\
&= \|(1-\alpha)(h-g)H^T\|_1 \\
&\leq (1-\alpha)\|h-g\|_1
\end{aligned}

{]} 其中 ( A = (1-\alpha)H\^{}T )，( b = \alphaJ(x) )。因此 (
T\_\sigma ) 仍是收缩映射，收敛性得证。

常用的1-Lipschitz激活函数包括： - \textbf{ReLU}：( \sigma(x) = \max(0,x)
)，是1-Lipschitz的 - \textbf{tanh}：导数的绝对值不超过1，是1-Lipschitz的
- \textbf{Sigmoid}：导数的最大值是0.25，是0.25-Lipschitz的，需要适当缩放
- \textbf{Leaky ReLU}：( \sigma(x) = \max(\beta x, x) ) 其中 ( 0
\leq \beta \leq 1 )，是1-Lipschitz的

\hypertarget{ux81eaux4fddux6301ux6df7ux5408ux4e0eux8c31ux504fux79fb}{%
\subsubsection{4.4.3
自保持混合与谱偏移}\label{ux81eaux4fddux6301ux6df7ux5408ux4e0eux8c31ux504fux79fb}}

为了抑制双随机传播可能导致的``过度平均化''，MQR引入了\textbf{自保持混合}机制：

{[} H\_\{eff\} = (1-\beta)I + \beta H {]}

其中 ( \beta \in [0,1] ) 控制混合强度。对应的动力学方程为：

{[} h\^{}\{(t+1)\} = (1-\alpha)h\textsuperscript{\{(t)\}H\_\{eff\}}T +
\alphaJ(x) {]}

\textbf{命题4.4}：( H\_\{eff\} )
仍然是双随机矩阵，且动力学系统的收敛性保持不变。

\textbf{证明}：双随机性验证： 1. 非负性：当 ( \beta \in [0,1] ) 且 ( H )
非负时，( H\_\{eff\} ) 非负。 2. 行和：( \sum\_j {[}(1-\beta)I\_\{ij\} +
\beta H\_\{ij\}{]} = (1-\beta) \cdot 1 + \beta \cdot 1 = 1 )。

收敛性：( H\_\{eff\} ) 的1-范数仍为1，因此收缩系数仍为 ( (1-\alpha) )。

自保持混合的谱意义：( H\_\{eff\} ) 的特征值是 ( H )
的特征值的仿射变换：( \lambda\_\{eff\} = (1-\beta) + \beta\lambda\_H
)。特别地，最大特征值仍为1，但次大特征值从 ( \lambda\_2(H) )

\hypertarget{ux7b2c5ux7ae0-ux5e7aux6a21ux968fux673aux6d41ux5f62ux4eceux9149ux7fa4ux5230ux53ccux968fux673aux52a8ux529bux5b66ux7684ux6570ux5b66ux57faux7840}{%
\section{第5章
幺模随机流形：从酉群到双随机动力学的数学基础}\label{ux7b2c5ux7ae0-ux5e7aux6a21ux968fux673aux6d41ux5f62ux4eceux9149ux7fa4ux5230ux53ccux968fux673aux52a8ux529bux5b66ux7684ux6570ux5b66ux57faux7840}}

\hypertarget{ux5f15ux8a00ux8d85ux8d8aux6b63ux4ea4ux7ea6ux675fux7684ux7a33ux5b9aux6027ux6846ux67b6}{%
\subsection{5.1
引言：超越正交约束的稳定性框架}\label{ux5f15ux8a00ux8d85ux8d8aux6b63ux4ea4ux7ea6ux675fux7684ux7a33ux5b9aux6027ux6846ux67b6}}

在循环神经网络的发展脉络中，正交约束（Orthogonal
Constraint）为解决梯度消失与爆炸问题提供了一条优雅的理论路径。通过将循环权重矩阵严格限制在正交群O(N)或酉群U(N)上，可以保证线性变换的等距性，从而在理论上消除因线性部分引起的梯度幅值畸变。然而，如第1章所讨论，严格的等距性可能对模型的表达能力构成限制。莫比乌斯量子环形网络（Möbius
Quantum Ring,
MQR）的核心创新之一，在于引入了一个更为丰富且同样具有严格稳定性保证的数学结构------\textbf{幺模随机流形}（Unistochastic
Manifold）。本章将深入探讨这一概念的数学起源、性质及其在MQR架构中的具体实现，为理解整个系统的动力学行为奠定理论基础。

幺模随机流形提供了一个介于无约束矩阵空间与严格酉群之间的``甜点''区域。它继承自酉群的优良稳定性特质，同时通过一个确定性的非线性映射------取模平方运算------引入了一种受控的``软性''非线性，这为模型提供了额外的表达灵活性。本章将首先回顾酉矩阵与双随机矩阵的基本理论，然后正式定义幺模随机矩阵与幺模随机流形，分析其关键性质，最后详细阐述MQR如何利用Cayley变换参数化酉矩阵，并通过幺模随机连接矩阵H=\textbar U\textbar²构建其核心的环形动力学。

\hypertarget{ux6570ux5b66ux9884ux5907ux77e5ux8bc6ux9149ux77e9ux9635ux4e0eux53ccux968fux673aux77e9ux9635}{%
\subsection{5.2
数学预备知识：酉矩阵与双随机矩阵}\label{ux6570ux5b66ux9884ux5907ux77e5ux8bc6ux9149ux77e9ux9635ux4e0eux53ccux968fux673aux77e9ux9635}}

\hypertarget{ux9149ux77e9ux9635ux4e0eux9149ux7fa4}{%
\subsubsection{5.2.1
酉矩阵与酉群}\label{ux9149ux77e9ux9635ux4e0eux9149ux7fa4}}

在复数域C上，一个N×N的矩阵U被称为\textbf{酉矩阵}（Unitary
Matrix），如果它满足：

{[} U\^{}\dagger U = U U\^{}\dagger = I\_N {]}

其中U^{\dagger}表示U的共轭转置（Hermitian
conjugate），I\_N是N阶单位矩阵。酉矩阵的集合构成一个李群，称为\textbf{酉群}（Unitary
Group），记作U(N)。酉矩阵具有以下关键性质：

\begin{enumerate}
\def\labelenumi{\arabic{enumi}.}
\tightlist
\item
  \textbf{保内积与保范数}：对于任意向量x ∈ C\^{}N，有‖Ux‖₂ =
  ‖x‖₂。这意味着酉变换是等距变换，不改变向量的长度。
\item
  \textbf{特征值位于单位圆上}：酉矩阵的所有特征值\lambda\_i满足\textbar \lambda\_i\textbar{}
  = 1。
\item
  \textbf{列向量构成标准正交基}：酉矩阵的列（和行）向量组是C\^{}N上的一组标准正交基。
\item
  \textbf{逆矩阵易求}：U^{-1} = U^{\dagger}。
\end{enumerate}

在实数域R上，对应的概念是\textbf{正交矩阵}（Orthogonal Matrix），满足U^{\mathsf{T}}U
= UU^{\mathsf{T}} =
I，其集合构成正交群O(N)。MQR架构主要工作在复数域，以利用相位信息，但实数版本作为特例包含在内。

\hypertarget{ux53ccux968fux673aux77e9ux9635}{%
\subsubsection{5.2.2 双随机矩阵}\label{ux53ccux968fux673aux77e9ux9635}}

一个N×N的实矩阵B被称为\textbf{双随机矩阵}（Doubly Stochastic
Matrix），如果其所有元素均为非负实数，且满足：

{[} \sum\emph{\{i=1\}\^{}\{N\} B}\{ij\} = 1 \quad \forall j,
\qquad \sum\emph{\{j=1\}\^{}\{N\} B}\{ij\} = 1 \quad \forall i {]}

即，矩阵的每一行之和与每一列之和都等于1。双随机矩阵具有以下重要性质：

\begin{enumerate}
\def\labelenumi{\arabic{enumi}.}
\tightlist
\item
  \textbf{保概率分布}：如果p是一个概率向量（各元素非负且和为1），那么Bp也是一个概率向量。这使得双随机矩阵天然适用于描述马尔可夫链的状态转移，其中B\_\{ij\}表示从状态j转移到状态i的概率。
\item
  \textbf{Birkhoff-von
  Neumann定理}：任何双随机矩阵都可以表示为置换矩阵的凸组合。这揭示了双随机矩阵与置换这一基本对称操作之间的深刻联系。
\item
  \textbf{谱性质}：双随机矩阵的谱半径ρ(B)=1，且总有一个特征值为1，对应的特征向量为全1向量（在适当归一化下）。
\item
  \textbf{收缩性}：对于任意向量x，在某种意义下（例如，相对于均匀分布的方差），双随机变换倾向于``平滑''或``混合''向量，减少其分量间的差异。
\end{enumerate}

双随机矩阵的收缩性质对于确保动力学的稳定性至关重要，但它也可能导致信息的过度均匀化，这正是MQR引入可调节的H\_eff和局部读出机制的原因之一。

\hypertarget{ux5e7aux6a21ux968fux673aux77e9ux9635ux5b9aux4e49ux4e0eux57faux672cux6027ux8d28}{%
\subsection{5.3
幺模随机矩阵：定义与基本性质}\label{ux5e7aux6a21ux968fux673aux77e9ux9635ux5b9aux4e49ux4e0eux57faux672cux6027ux8d28}}

\hypertarget{ux4eceux9149ux77e9ux9635ux5230ux53ccux968fux673aux77e9ux9635ux7684ux6620ux5c04}{%
\subsubsection{5.3.1
从酉矩阵到双随机矩阵的映射}\label{ux4eceux9149ux77e9ux9635ux5230ux53ccux968fux673aux77e9ux9635ux7684ux6620ux5c04}}

幺模随机矩阵的概念建立在一个简单的观察之上：给定一个酉矩阵U，如果我们对其每个元素取模平方，得到的矩阵H，其元素定义为H\_\{ij\}
= \textbar U\_\{ij\}\textbar²，会具有怎样的性质？

\textbf{定理 5.1 (幺模随机矩阵的基本性质)}：设U ∈
U(N)是一个N阶酉矩阵，定义矩阵H，其中H\_\{ij\} =
\textbar U\_\{ij\}\textbar²。则H是一个双随机矩阵。

\textbf{证明}： 由于U是酉矩阵，其列向量构成标准正交基。对于第j列向量u\_j
= (U\_\{1j\}, U\_\{2j\}, \ldots, U\_\{Nj\})^{\mathsf{T}}，其标准正交性意味着：

{[} \sum\emph{\{i=1\}\^{}\{N\} \textbar U}\{ij\}\textbar\^{}2 =
\textbar u\_j\textbar\_2\^{}2 = 1 \quad \forall j {]}

这正是H的列和条件。同理，U的行向量也构成标准正交基（因为U^{\dagger}U = I 且 UU^{\dagger} =
I），因此对于第i行向量，有：

{[} \sum\emph{\{j=1\}\^{}\{N\} \textbar U}\{ij\}\textbar\^{}2 = 1
\quad \forall i {]}

这正是H的行和条件。此外，显然有H\_\{ij\} = \textbar U\_\{ij\}\textbar² \geq
0。因此，H满足双随机矩阵的所有定义条件。证毕。

我们将通过上述映射从酉矩阵U得到的双随机矩阵H称为\textbf{幺模随机矩阵}（Unistochastic
Matrix）。所有N阶幺模随机矩阵的集合构成了N阶双随机矩阵集合的一个子集，我们称之为\textbf{幺模随机流形}（Unistochastic
Manifold），尽管严格来说它并非一个光滑流形。

\hypertarget{ux5e7aux6a21ux968fux673aux6d41ux5f62ux7684ux7279ux6027}{%
\subsubsection{5.3.2
幺模随机流形的特性}\label{ux5e7aux6a21ux968fux673aux6d41ux5f62ux7684ux7279ux6027}}

幺模随机流形具有几个值得注意的特性：

\begin{enumerate}
\def\labelenumi{\arabic{enumi}.}
\tightlist
\item
  \textbf{非凸性}：幺模随机矩阵的集合不是凸集。两个幺模随机矩阵的凸组合不一定仍是幺模随机的。这与一般的双随机矩阵集合（是一个凸多面体）形成对比。
\item
  \textbf{维数}：酉群U(N)是N²维的实流形（考虑实维数）。取模平方映射U ↦ H
  =
  \textbar U\textbar²不是单射，因为U的每个元素的相位信息在映射中丢失了。因此，幺模随机流形的实维数小于N²，但通常远大于一般双随机矩阵凸多面体的顶点（置换矩阵）数量所暗示的维度。它提供了一个高维、连续且结构丰富的双随机矩阵参数族。
\item
  \textbf{对称性}：映射U ↦
  \textbar U\textbar²具有对称性。对于任意对角相位矩阵D =
  diag(e\^{}\{i\theta₁\}, \ldots,
  e\^{}\{i\theta\_N\})，矩阵U和D·U映射到同一个幺模随机矩阵H，因为\textbar(D·U)\emph{\{ij\}\textbar²
  = \textbar e\^{}\{i\theta\_i\} U}\{ij\}\textbar² =
  \textbar U\_\{ij\}\textbar²。这种规范对称性（Gauge
  Symmetry）在优化和表示中需要考虑。
\end{enumerate}

幺模随机矩阵作为双随机矩阵，继承了后者的稳定性（收缩性），同时，由于它们源自酉矩阵，其构造过程自然地与一种保持范数的变换相关联。在MQR中，这种关联被巧妙地利用：学习过程在稳定、易于优化的酉群上进行（通过特殊的参数化），而前向推理则在具有概率解释和稳定收缩性的幺模随机流形上进行。

\hypertarget{cayleyux53d8ux6362ux9149ux77e9ux9635ux7684ux6709ux6548ux53c2ux6570ux5316}{%
\subsection{5.4
Cayley变换：酉矩阵的有效参数化}\label{cayleyux53d8ux6362ux9149ux77e9ux9635ux7684ux6709ux6548ux53c2ux6570ux5316}}

为了在神经网络中学习和优化酉矩阵U，我们需要一种有效的参数化方法，使得：
1. 对于任意参数，生成的矩阵保证是酉矩阵。 2.
参数化覆盖足够大的酉群子集（理想情况下是稠密的）。 3.
前向计算（从参数生成U）和反向传播（计算梯度）是高效且数值稳定的。

MQR采用\textbf{Cayley变换}（Cayley
Transform）作为其默认的酉矩阵参数化方法，它完美地满足了上述要求。

\hypertarget{cayleyux53d8ux6362ux7684ux5b9aux4e49ux4e0eux6027ux8d28}{%
\subsubsection{5.4.1
Cayley变换的定义与性质}\label{cayleyux53d8ux6362ux7684ux5b9aux4e49ux4e0eux6027ux8d28}}

Cayley变换建立了斜埃尔米特矩阵（Skew-Hermitian
Matrix）与酉矩阵之间的一一对应关系（在忽略一个特征值为-1的情况下）。一个矩阵A被称为斜埃尔米特矩阵，如果满足A^{\dagger}
= -A。

\textbf{定义 5.1 (Cayley变换)}：设A是一个N×N的斜埃尔米特矩阵，且I +
A是可逆的（即-1不是A的特征值）。则通过Cayley变换定义的矩阵U为：

{[} U = (I - A)(I + A)\^{}\{-1\} {]}

可以证明，这样定义的U是一个酉矩阵。反之，对于任意酉矩阵U，若I +
U可逆（即-1不是U的特征值），则可以通过逆Cayley变换得到一个斜埃尔米特矩阵A：

{[} A = (I - U)(I + U)\^{}\{-1\} {]}

\textbf{定理 5.2}：由定义5.1中Cayley变换产生的U是酉矩阵。

\textbf{证明}： 计算U^{\dagger}U： {[}

\begin{aligned}
U^\dagger U &= \left[(I + A)^{-1}\right]^\dagger (I - A)^\dagger (I - A)(I + A)^{-1} \\
&= (I + A^\dagger)^{-1} (I - A^\dagger)(I - A)(I + A)^{-1} \quad \text{利用}(X^{-1})^\dagger = (X^\dagger)^{-1} \\
&= (I - A)^{-1} (I + A)(I - A)(I + A)^{-1} \quad \text{因为} A^\dagger = -A \\
&= (I - A)^{-1} (I - A)(I + A)(I + A)^{-1} \quad \text{因为}(I+A)\text{与}(I-A)\text{可交换（都关于A的多项式）} \\
&= I \cdot I = I
\end{aligned}

{]} 类似可证UU^{\dagger} = I。证毕。

\hypertarget{mqrux4e2dux7684ux53c2ux6570ux5316ux5b9eux73b0}{%
\subsubsection{5.4.2
MQR中的参数化实现}\label{mqrux4e2dux7684ux53c2ux6570ux5316ux5b9eux73b0}}

在MQR的实现中，我们并不直接存储和优化酉矩阵U，而是存储和优化一个无约束的实矩阵W，然后从中构造斜埃尔米特矩阵A，最后通过Cayley变换得到U。具体步骤如下：

\begin{enumerate}
\def\labelenumi{\arabic{enumi}.}
\item
  \textbf{生成斜埃尔米特矩阵}：设W是一个任意N×N的实矩阵（在复数域情况下，可以是复矩阵，但为简化，MQR默认使用实参数化生成实斜对称矩阵，对应于特殊正交群SO(N)的子集，或通过复化得到酉矩阵）。我们通过以下方式构造一个斜对称（实斜埃尔米特）矩阵A：
  {[} A = W - W\^{}T {]} 显然，A^{\mathsf{T}} =
  -A，满足斜对称条件。对于复数域扩展，可以构造A = W -
  W^{\dagger}，确保斜埃尔米特性。
\item
  \textbf{应用Cayley变换}：计算酉矩阵U = (I - A)(I +
  A)^{-1}。直接求逆的复杂度是O(N³)，对于较大的N不可行。MQR采用迭代法或利用矩阵的Woodbury恒等式进行高效计算。一种稳定且高效的计算方式是求解线性系统：
  {[} (I + A) U = I - A {]} 由于I +
  A是已知的，我们可以使用共轭梯度法（Conjugate
  Gradient）或其他迭代求解器来得到U，这通常比显式求逆更高效且数值稳定，尤其当A是稀疏或具有特殊结构时。在MQR的当前实现中，对于中等规模的N，可能直接使用PyTorch的\texttt{torch.linalg.solve}函数。
\item
  \textbf{处理特征值-1问题}：理论上，如果A有特征值-1，则I+A不可逆，Cayley变换无定义。在实践中，可以通过对A添加一个小的正则化项来避免，例如使用A'
  = A -
  \epsilonI，其中\epsilon是一个很小的正数，这会使A'的特征值偏移，远离-1。MQR的代码中通常包含此类数值稳定性措施。
\end{enumerate}

这种参数化的优势在于，优化变量W存在于无约束的欧几里得空间R\^{}\{N×N\}中，我们可以使用标准的梯度下降优化器（如Adam）来更新W。梯度通过Cayley变换的反向传播自动计算。这比直接在酉流形上使用黎曼优化器（如Geodesic
SGD）更为简单和高效。

\begin{figure}
\centering
\includegraphics{figures/diagram_5.png}
\caption{图 5}
\end{figure}

\emph{图5.1：MQR中幺模随机矩阵生成与优化流程示意图。展示了从无约束参数W出发，通过Cayley变换生成酉矩阵U，进而得到幺模随机连接矩阵H的确定性过程。H作为环形动力学的核心，驱动系统达到平衡态h}，最终通过损失函数计算梯度并反向传播更新参数W，形成完整的优化闭环。*

\hypertarget{ux5e7aux6a21ux968fux673aux8fdeux63a5ux77e9ux9635hux7684ux52a8ux529bux5b66ux89d2ux8272}{%
\subsection{5.5
幺模随机连接矩阵H的动力学角色}\label{ux5e7aux6a21ux968fux673aux8fdeux63a5ux77e9ux9635hux7684ux52a8ux529bux5b66ux89d2ux8272}}

在MQR架构中，由学习到的酉矩阵U生成的幺模随机矩阵H =
\textbar U\textbar²，扮演着环形网络中``连接权重''或``状态转移矩阵''的核心角色。考虑MQR的基本松弛动力学方程（暂不考虑非线性激活和注入项）：

{[} h\^{}\{(t+1)\} = h\^{}\{(t)\} H\^{}T {]}

或者，考虑带有自保持混合的版本：

{[} h\^{}\{(t+1)\} = (1-\beta) h\^{}\{(t)\} + \beta h\^{}\{(t)\} H\^{}T
= h\^{}\{(t)\} \left[(1-\beta)I + \beta H^T\right] = h\^{}\{(t)\}
H\_\{\text{eff}\}\^{}T {]}

其中H\_eff = (1-\beta)I +
\betaH仍然是一个双随机矩阵（因为它是单位阵与双随机矩阵的凸组合）。

\hypertarget{ux6536ux7f29ux6027ux4e0eux7a33ux5b9aux6027ux5206ux6790}{%
\subsubsection{5.5.1
收缩性与稳定性分析}\label{ux6536ux7f29ux6027ux4e0eux7a33ux5b9aux6027ux5206ux6790}}

由于H（以及H\_eff）是双随机矩阵，其谱半径ρ(H) =
1，且有一个主特征值1。对于任意初始隐藏状态向量h⁽⁰⁾，迭代h⁽ᵗ⁺¹⁾ = h⁽ᵗ⁾
H^{\mathsf{T}}将收敛到一个稳态h*，这个稳态是H^{\mathsf{T}}对应于特征值1的左特征向量（或H对应于特征值1的右特征向量）。更精确地说：

\textbf{定理 5.3
(双随机迭代的收敛性)}：设H是一个本原双随机矩阵（Primitive Doubly
Stochastic
Matrix），即存在某个正整数k使得H\^{}k的所有元素为正。则对于任意初始概率向量h⁽⁰⁾（元素非负，和为1），迭代序列\{h⁽ᵗ⁾\}定义为h⁽ᵗ⁺¹⁾
= h⁽ᵗ⁾ H^{\mathsf{T}}，当t \rightarrow \infty时，以指数速度收敛到唯一的稳态概率向量h\emph{，满足h}
H^{\mathsf{T}} = h*。

\textbf{证明概要}：本原性保证了H的Perron-Frobenius特征值1是代数和几何重数均为1的，且其他特征值的模严格小于1。因此，动力学的收敛速率由第二大特征值的模\textbar \lambda₂\textbar 决定。指数收敛速度保证了MQR推理过程的效率。

在MQR中，h通常不限制为概率向量（元素可正可负，和不一定是1），但双随机矩阵作用于向量时，仍然具有收缩向量各分量相对于其均值之差的效应。这为动力学的稳定性提供了保障。即使加入满足1-Lipschitz条件的非线性激活函数\sigma(·)（如tanh），只要适当处理，仍可证明迭代的收敛性。

\hypertarget{ux4f5cux4e3aux4fe1ux606fux6df7ux5408ux5668}{%
\subsubsection{5.5.2
作为信息混合器}\label{ux4f5cux4e3aux4fe1ux606fux6df7ux5408ux5668}}

H矩阵定义了环形网络上各节点间信息交换的强度。H\_\{ij\} =
\textbar U\_\{ij\}\textbar²可以解释为从节点j到节点i的``影响权重''。由于H是双随机的，每个节点输出的总影响（行和）与接收的总影响（列和）是平衡的，均为1。这创造了一个保守的``信息流''系统，没有节点会过度放大或衰减信息。

幺模随机矩阵H的特殊性在于，它的元素不是独立学习的，而是通过酉矩阵U的模平方产生。这隐含了一种结构约束：H的行和列向量并非任意，它们必须对应于某个酉矩阵的行和列向量的模平方。这种约束可能引导模型学习到具有某种``相干性''或``互补性''的混合模式，类似于量子力学中概率幅与概率之间的关系。

\hypertarget{ux6269ux5c55ux590dux6570ux9149ux63a8ux7406ux4e0eux76f8ux4f4dux4fe1ux606f}{%
\subsection{5.6
扩展：复数酉推理与相位信息}\label{ux6269ux5c55ux590dux6570ux9149ux63a8ux7406ux4e0eux76f8ux4f4dux4fe1ux606f}}

MQR架构提供了一个可选模式\texttt{dynamics\_mode=\textquotesingle{}unitary\textquotesingle{}}，在此模式下，环形动力学直接使用酉矩阵U的共轭转置U^{\dagger}进行传播，而不是使用双随机矩阵H：

{[} h\^{}\{(t+1)\} = h\^{}\{(t)\} U\^{}\dagger {]}

这里h是复数向量。由于U^{\dagger}是酉矩阵，该变换严格保范数：‖h⁽ᵗ⁺¹⁾‖₂ =
‖h⁽ᵗ⁾‖₂。动力学将不会收敛

\hypertarget{ux7b2c6ux7ae0-ux5e7aux6a21ux968fux673aux6d41ux5f62ux4e0eux53ccux968fux673aux8fdeux63a5ux77e9ux9635ux7684ux6570ux5b66ux57faux7840}{%
\section{第6章
幺模随机流形与双随机连接矩阵的数学基础}\label{ux7b2c6ux7ae0-ux5e7aux6a21ux968fux673aux6d41ux5f62ux4e0eux53ccux968fux673aux8fdeux63a5ux77e9ux9635ux7684ux6570ux5b66ux57faux7840}}

\hypertarget{ux5f15ux8a00ux4eceux9149ux77e9ux9635ux5230ux5e7aux6a21ux968fux673aux77e9ux9635}{%
\subsection{6.1
引言：从酉矩阵到幺模随机矩阵}\label{ux5f15ux8a00ux4eceux9149ux77e9ux9635ux5230ux5e7aux6a21ux968fux673aux77e9ux9635}}

在莫比乌斯量子环网络（Möbius Quantum Ring Network,
MQR）的核心架构中，连接矩阵的数学性质决定了整个系统的动力学行为。与传统的循环神经网络不同，MQR不直接学习权重矩阵，而是通过学习一个酉矩阵（Unitary
Matrix），然后通过特定的变换得到连接矩阵。这种设计选择并非偶然，而是基于深刻的数学原理和物理类比。本章将深入探讨这一数学基础，详细阐述从酉矩阵到幺模随机矩阵的变换过程，分析其数学性质，并证明这些性质如何确保网络的稳定性和表达能力。

幺模随机矩阵（Unistochastic Matrix）是双随机矩阵（Doubly Stochastic
Matrix）的一个特殊子类，它源于量子力学中的概念。在量子系统中，酉矩阵描述的是封闭系统的演化，而幺模随机矩阵则描述了概率的转移。MQR巧妙地将这两个概念结合起来，创造了一种既保持量子系统优雅数学结构，又适用于经典机器学习任务的网络架构。

\hypertarget{ux9149ux77e9ux9635ux7684ux57faux672cux6027ux8d28ux4e0eux53c2ux6570ux5316ux65b9ux6cd5}{%
\subsection{6.2
酉矩阵的基本性质与参数化方法}\label{ux9149ux77e9ux9635ux7684ux57faux672cux6027ux8d28ux4e0eux53c2ux6570ux5316ux65b9ux6cd5}}

\hypertarget{ux9149ux77e9ux9635ux7684ux5b9aux4e49ux4e0eux57faux672cux6027ux8d28-1}{%
\subsubsection{6.2.1
酉矩阵的定义与基本性质}\label{ux9149ux77e9ux9635ux7684ux5b9aux4e49ux4e0eux57faux672cux6027ux8d28-1}}

酉矩阵是复数域上的方阵，满足以下条件：

{[} U \in C\^{}\{N \times N\}, \quad U\^{}\dagger U = U
U\^{}\dagger = I\_N {]}

其中 (U\^{}\dagger) 表示 (U) 的共轭转置，(I\_N) 是 (N \times N)
的单位矩阵。酉矩阵具有以下重要性质：

\begin{enumerate}
\def\labelenumi{\arabic{enumi}.}
\tightlist
\item
  \textbf{保范性}：对于任意向量 (x \in C\^{}N)，有
  (\textbar Ux\textbar\_2 = \textbar x\textbar\_2)，其中
  (\textbar{}\cdot\textbar\_2) 表示欧几里得范数。
\item
  \textbf{特征值位于单位圆上}：酉矩阵的所有特征值 (\lambda\_i) 满足
  (\textbar{}\lambda\_i\textbar{} = 1)。
\item
  \textbf{特征向量正交}：不同特征值对应的特征向量相互正交。
\item
  \textbf{行列式的模为1}：(\textbar{}\det(U)\textbar{} = 1)。
\end{enumerate}

这些性质使得酉矩阵成为描述保守系统演化的理想工具。在MQR中，酉矩阵 (U)
被用作生成连接矩阵的基础。

\hypertarget{cayleyux53c2ux6570ux5316ux65b9ux6cd5}{%
\subsubsection{6.2.2
Cayley参数化方法}\label{cayleyux53c2ux6570ux5316ux65b9ux6cd5}}

在深度学习中，直接优化酉矩阵面临挑战，因为酉矩阵空间（酉群
(U(N))）是一个流形而非向量空间。MQR采用Cayley变换来实现有效的参数化：

{[} U = (I + iA)(I - iA)\^{}\{-1\} {]}

其中 (A) 是一个任意的斜埃尔米特矩阵（Skew-Hermitian Matrix），满足
(A\^{}\dagger = -A)。在实际实现中，我们通常参数化一个无约束的复数矩阵 (B
\in C\^{}\{N \times N\})，然后通过以下方式构造斜埃尔米特矩阵：

{[} A = \frac{1}{2}(B - B\^{}\dagger) {]}

Cayley变换具有以下优点： 1.
\textbf{自动保持酉性}：对于任意斜埃尔米特矩阵 (A)，变换结果 (U)
总是酉矩阵。 2. \textbf{可逆性}：变换是可逆的，逆变换为 (A = i(I - U)(I
+ U)\^{}\{-1\})。 3. \textbf{数值稳定性}：当 (A)
的范数较小时，变换具有良好的数值性质。

在MQR的实现中，Cayley参数化允许我们使用标准的梯度下降方法优化无约束参数矩阵
(B)，同时确保生成的 (U) 始终保持酉性。

\hypertarget{ux5176ux4ed6ux53c2ux6570ux5316ux65b9ux6cd5ux6bd4ux8f83}{%
\subsubsection{6.2.3
其他参数化方法比较}\label{ux5176ux4ed6ux53c2ux6570ux5316ux65b9ux6cd5ux6bd4ux8f83}}

除了Cayley变换，还有其他参数化酉矩阵的方法：

\begin{enumerate}
\def\labelenumi{\arabic{enumi}.}
\tightlist
\item
  \textbf{指数映射}：(U = \exp(iH))，其中 (H)
  是埃尔米特矩阵。这种方法在理论物理中常见，但计算矩阵指数在深度学习中成本较高。
\item
  \textbf{正交矩阵的复数推广}：对于实数情况，可以使用Householder变换或Givens旋转的乘积。
\item
  \textbf{流形优化}：直接在Stiefel流形上优化，使用黎曼梯度下降。
\end{enumerate}

MQR选择Cayley变换主要是出于计算效率和数值稳定性的考虑。Cayley变换只需要矩阵求逆，而不需要计算矩阵指数，这在深度学习的迭代优化中更为实用。

\hypertarget{ux5e7aux6a21ux968fux673aux77e9ux9635ux7684ux6784ux9020ux4e0eux6027ux8d28}{%
\subsection{6.3
幺模随机矩阵的构造与性质}\label{ux5e7aux6a21ux968fux673aux77e9ux9635ux7684ux6784ux9020ux4e0eux6027ux8d28}}

\hypertarget{ux4eceux9149ux77e9ux9635ux5230ux5e7aux6a21ux968fux673aux77e9ux9635-1}{%
\subsubsection{6.3.1
从酉矩阵到幺模随机矩阵}\label{ux4eceux9149ux77e9ux9635ux5230ux5e7aux6a21ux968fux673aux77e9ux9635-1}}

幺模随机矩阵 (H) 通过以下方式从酉矩阵 (U) 构造：

{[} H = \textbar U\textbar\^{}2 {]}

其中绝对值平方运算 (\textbar{} \cdot \textbar\^{}2)
是逐元素进行的。具体来说，对于酉矩阵 (U) 的每个元素
(u\_\{ij\})，幺模随机矩阵的对应元素为：

{[} h\_\{ij\} = \textbar u\_\{ij\}\textbar\^{}2 {]}

这个变换具有深刻的物理意义：在量子力学中，如果 (U)
描述了一个量子系统的演化，那么 (H) 描述了从状态 (i) 到状态 (j)
的转移概率。这就是著名的Born规则在矩阵形式下的体现。

\hypertarget{ux5e7aux6a21ux968fux673aux77e9ux9635ux7684ux57faux672cux6027ux8d28}{%
\subsubsection{6.3.2
幺模随机矩阵的基本性质}\label{ux5e7aux6a21ux968fux673aux77e9ux9635ux7684ux57faux672cux6027ux8d28}}

幺模随机矩阵 (H) 具有以下重要性质：

\begin{enumerate}
\def\labelenumi{\arabic{enumi}.}
\item
  \textbf{双随机性}：

  \begin{itemize}
  \tightlist
  \item
    行和归一：(\sum\emph{\{j=1\}\^{}N h}\{ij\} = 1) 对所有 (i) 成立
  \item
    列和归一：(\sum\emph{\{i=1\}\^{}N h}\{ij\} = 1) 对所有 (j) 成立
  \end{itemize}

  这可以通过酉矩阵的性质证明： {[} \sum\emph{\{j=1\}\^{}N
  \textbar u}\{ij\}\textbar\^{}2 = \sum\emph{\{j=1\}\^{}N u}\{ij\}
  \overline{u_{ij}} = (UU\^{}\dagger)\_\{ii\} = 1 {]}
  类似地，列和归一性来自 (U\^{}\dagger U = I)。
\item
  \textbf{非负性}：所有元素 (h\_\{ij\} = \textbar u\_\{ij\}\textbar\^{}2
  \geq 0)。
\item
  \textbf{特征值性质}：幺模随机矩阵的特征值满足
  (\textbar{}\lambda\_i\textbar{} \leq 1)，且至少有一个特征值为1。
\item
  \textbf{Perron-Frobenius定理的应用}：由于 (H)
  是非负矩阵且不可约（在大多数情况下），根据Perron-Frobenius定理，它有一个唯一的最大特征值
  (\lambda\_1 = 1)，对应的特征向量所有分量均为正。
\end{enumerate}

双随机性在MQR中具有关键意义：它确保了信息在环形网络中的传播是``守恒的''，不会出现某些节点过度激活而其他节点激活不足的情况。这为网络的稳定性提供了数学保证。

\hypertarget{ux5e7aux6a21ux968fux673aux77e9ux9635ux4e0eux4e00ux822cux53ccux968fux673aux77e9ux9635ux7684ux533aux522b}{%
\subsubsection{6.3.3
幺模随机矩阵与一般双随机矩阵的区别}\label{ux5e7aux6a21ux968fux673aux77e9ux9635ux4e0eux4e00ux822cux53ccux968fux673aux77e9ux9635ux7684ux533aux522b}}

并非所有双随机矩阵都是幺模随机的。幺模随机矩阵是双随机矩阵的一个真子集，具有额外的结构约束。这种区别可以通过以下方式理解：

\begin{enumerate}
\def\labelenumi{\arabic{enumi}.}
\item
  \textbf{维度考虑}：所有 (N \times N)
  双随机矩阵构成的Birkhoff多面体维度为
  ((N-1)\^{}2)，而幺模随机流形的维度约为
  (N\^{}2)（与酉群维度相同）。对于 (N \textgreater{}
  2)，幺模随机矩阵只是双随机矩阵的一个低维子流形。
\item
  \textbf{可逆性}：每个幺模随机矩阵 (H) 都至少有一个酉矩阵 (U) 使得 (H =
  \textbar U\textbar\^{}2)，但这样的 (U) 通常不唯一（存在相位自由度）。
\item
  \textbf{量子类比}：幺模随机矩阵可以解释为量子信道，而一般的双随机矩阵则不一定有这种解释。
\end{enumerate}

在MQR中，使用幺模随机矩阵而非一般双随机矩阵的主要优势在于： -
通过酉矩阵的参数化，确保了矩阵的良好数值性质 -
提供了与量子系统的直观类比，便于理论分析 -
保持了与复数域操作的兼容性，支持复数推理模式

\hypertarget{ux5e7aux6a21ux968fux673aux6d41ux5f62ux7684ux51e0ux4f55ux7ed3ux6784}{%
\subsection{6.4
幺模随机流形的几何结构}\label{ux5e7aux6a21ux968fux673aux6d41ux5f62ux7684ux51e0ux4f55ux7ed3ux6784}}

\hypertarget{ux6d41ux5f62ux5b9aux4e49ux4e0eux5c40ux90e8ux5750ux6807}{%
\subsubsection{6.4.1
流形定义与局部坐标}\label{ux6d41ux5f62ux5b9aux4e49ux4e0eux5c40ux90e8ux5750ux6807}}

幺模随机流形 (U\_N) 定义为所有 (N \times N)
幺模随机矩阵的集合：

{[} U\_N = \{ H \in R\^{}\{N \times N\} : H =
\textbar U\textbar\^{}2, U \in U(N) \} {]}

这是一个紧致的光滑流形，维度为 (N\^{}2)。在点 (H =
\textbar U\textbar\^{}2) 处，切空间可以描述为：

{[} T\_H U\_N = \{ \delta H : \delta H = U \circ \delta X +
\overline{U} \circ \delta \overline{X} \text{ 对某个斜埃尔米特 }
\delta X \} {]}

其中 (\circ) 表示逐元素乘积（Hadamard积）。

\hypertarget{ux6d4bux5730ux7ebfux4e0eux8dddux79bb}{%
\subsubsection{6.4.2
测地线与距离}\label{ux6d4bux5730ux7ebfux4e0eux8dddux79bb}}

在幺模随机流形上，可以定义自然的黎曼度量。两点 (H\_1 =
\textbar U\_1\textbar\^{}2) 和 (H\_2 = \textbar U\_2\textbar\^{}2)
之间的距离可以通过它们在酉群中的提升来定义：

{[} d\_\{U\}(H\_1, H\_2) = \min\emph{\{U\_1, U\_2:
H\_i=\textbar U\_i\textbar\^{}2\} d}\{U(N)\}(U\_1, U\_2) {]}

其中 (d\_\{U(N)\})
是酉群上的标准距离。在实际优化中，MQR并不直接在幺模随机流形上进行黎曼优化，而是通过参数化酉矩阵来间接优化，这简化了实现同时保持了数学一致性。

\hypertarget{ux4e34ux754cux70b9ux4e0eux4f18ux5316ux666fux89c2}{%
\subsubsection{6.4.3
临界点与优化景观}\label{ux4e34ux754cux70b9ux4e0eux4f18ux5316ux666fux89c2}}

幺模随机流形的优化景观对于理解MQR的训练动态至关重要。损失函数 (L(H))
在流形上的梯度下降可以表示为：

{[} \frac{dH}{dt} = -\text{grad} L(H) {]}

其中 (\text{grad} L(H))
是损失函数在流形上的黎曼梯度。由于流形是紧致的，任何光滑函数在其上都有临界点。MQR的收敛性部分依赖于幺模随机流形的这种性质。

\hypertarget{ux53ccux968fux673aux8fdeux63a5ux77e9ux9635ux7684ux52a8ux529bux5b66ux610fux4e49}{%
\subsection{6.5
双随机连接矩阵的动力学意义}\label{ux53ccux968fux673aux8fdeux63a5ux77e9ux9635ux7684ux52a8ux529bux5b66ux610fux4e49}}

\hypertarget{ux9a6cux5c14ux53efux592bux94feux89e3ux91ca}{%
\subsubsection{6.5.1
马尔可夫链解释}\label{ux9a6cux5c14ux53efux592bux94feux89e3ux91ca}}

幺模随机矩阵 (H)
可以解释为一个齐次马尔可夫链的转移矩阵。对于环上的节点状态向量 (h
\in R\^{}N)，更新规则：

{[} h\^{}\{(t+1)\} = h\^{}\{(t)\} H\^{}T {]}

描述了概率分布（或激活强度）在环上的传播。由于 (H)
是双随机的，这个马尔可夫链具有以下性质：

\begin{enumerate}
\def\labelenumi{\arabic{enumi}.}
\item
  \textbf{平稳分布}：均匀分布 (\pi = (\frac{1}{N}, \frac{1}{N}, \ldots,
  \frac{1}{N})) 是一个平稳分布，因为 (\pi H = \pi)。
\item
  \textbf{可逆性}：如果 (H) 是对称的（即 (H =
  H\^{}T)），则马尔可夫链是可逆的，满足细致平衡条件。
\item
  \textbf{收敛速度}：收敛到平稳分布的速度由第二大特征值的模
  (\textbar{}\lambda\_2\textbar) 决定，这个值越小，收敛越快。
\end{enumerate}

在MQR中，动力学过程不是简单的马尔可夫链，因为存在注入项和可能的非线性激活，但双随机性确保了基本的稳定性。

\hypertarget{ux4fe1ux606fux8bbaux89c6ux89d2}{%
\subsubsection{6.5.2 信息论视角}\label{ux4fe1ux606fux8bbaux89c6ux89d2}}

从信息论的角度看，双随机矩阵保持熵的下界。对于概率分布 (p) 和双随机矩阵
(H)，有：

{[} S(pH) \geq S(p) {]}

其中 (S(\cdot))
是香农熵。这意味着通过双随机矩阵的变换不会减少不确定性，这与信息传播中的``不创造信息''原则一致。

在MQR的上下文中，这可以解释为：环形动力学不会自发地增加信息的确定性，所有的模式识别能力都来自注入项与连接矩阵的相互作用。

\hypertarget{ux4e0eux6269ux6563ux8fc7ux7a0bux7684ux7c7bux6bd4}{%
\subsubsection{6.5.3
与扩散过程的类比}\label{ux4e0eux6269ux6563ux8fc7ux7a0bux7684ux7c7bux6bd4}}

双随机连接矩阵也可以看作图上扩散过程的离散化。连续时间的扩散方程：

{[} \frac{\partial h}{\partial t} = \Delta h {]}

其中 (\Delta)
是图拉普拉斯算子。对于环形图，如果连接矩阵是双随机的，离散更新 (h
\leftarrow hH\^{}T) 近似于扩散过程的一个时间步。

这种类比有助于理解MQR的全局信息混合能力：就像扩散过程最终会使温度均匀分布一样，环形动力学最终会使信息在环上均匀混合，除非有外部输入（注入项）持续驱动系统偏离平衡。

\hypertarget{ux6709ux6548ux8fdeux63a5ux77e9ux9635ux4e0eux81eaux4fddux6301ux6df7ux5408}{%
\subsection{6.6
有效连接矩阵与自保持混合}\label{ux6709ux6548ux8fdeux63a5ux77e9ux9635ux4e0eux81eaux4fddux6301ux6df7ux5408}}

\hypertarget{ux8fc7ux5ea6ux5e73ux5747ux5316ux95eeux9898}{%
\subsubsection{6.6.1
过度平均化问题}\label{ux8fc7ux5ea6ux5e73ux5747ux5316ux95eeux9898}}

纯粹的幺模随机连接矩阵可能导致``过度平均化''问题：由于双随机性，信息在环上传播时可能过快地被均匀化，丢失了局部结构信息。这在某些任务中可能不利于保持输入的细粒度特征。

\hypertarget{ux6709ux6548ux8fdeux63a5ux77e9ux9635ux7684ux6784ux9020}{%
\subsubsection{6.6.2
有效连接矩阵的构造}\label{ux6709ux6548ux8fdeux63a5ux77e9ux9635ux7684ux6784ux9020}}

为了解决这个问题，MQR引入了有效连接矩阵的概念：

{[} H\_\{\text{eff}\} = (1-\beta)I + \beta H {]}

其中 (\beta \in [0, 1]) 是混合参数，(I)
是单位矩阵。这种构造被称为``自保持混合''，因为它让每个节点保留一部分自身状态，只将剩余部分与邻居混合。

\hypertarget{ux6570ux5b66ux6027ux8d28ux5206ux6790}{%
\subsubsection{6.6.3
数学性质分析}\label{ux6570ux5b66ux6027ux8d28ux5206ux6790}}

有效连接矩阵 (H\_\{\text{eff}\}) 具有以下性质：

\begin{enumerate}
\def\labelenumi{\arabic{enumi}.}
\item
  \textbf{保持双随机性}：如果 (H) 是双随机的，那么 (H\_\{\text{eff}\})
  也是双随机的。 {[}

  \begin{aligned}
  \sum_j (H_{\text{eff}})_{ij} &= \sum_j [(1-\beta)\delta_{ij} + \beta h_{ij}] \\
  &= (1-\beta) + \beta \sum_j h_{ij} = (1-\beta) + \beta \cdot 1 = 1
  \end{aligned}

  {]} 列和归一性同理。
\item
  \textbf{特征值关系}：如果 (H) 的特征值为 (\{\lambda\emph{k\})，那么
  (H}\{\text{eff}\}) 的特征值为 (\{(1-\beta) +
  \beta\lambda\_k\})。特别地，最大特征值仍为1，但其他特征值向1收缩： {[}
  \lambda\emph{k(H}\{\text{eff}\}) = (1-\beta) + \beta\lambda\_k(H) {]}
  这意味着 (\textbar{}\lambda\emph{k(H}\{\text{eff}\})\textbar{}
  \geq \textbar{}\lambda\_k(H)\textbar) 对于
  (\textbar{}\lambda\_k(H)\textbar{} \textless{}
  1)，实际上减慢了收敛速度。
\item
  \textbf{控制混合速度}：参数 (\beta) 控制了信息混合的速度。当 (\beta =
  0) 时，(H\_\{\text{eff}\} = I)，节点完全独立；当 (\beta = 1)
  时，(H\_\{\text{eff}\} = H)，完全使用原始连接矩阵。
\end{enumerate}

\hypertarget{ux52a8ux529bux5b66ux5f71ux54cd}{%
\subsubsection{6.6.4 动力学影响}\label{ux52a8ux529bux5b66ux5f71ux54cd}}

自保持混合对MQR的动力学有重要影响：

\begin{enumerate}
\def\labelenumi{\arabic{enumi}.}
\tightlist
\item
  \textbf{减缓收敛}：由于特征值更接近1，系统收敛到平衡态的速度变慢，这允许更复杂的瞬态动力学。
\item
  \textbf{保持局部信息}：每个节点保留一部分自身状态，有助于保持输入的局部特征。
\item
  \textbf{提高表达能力}：通过调整
  (\beta)，可以控制全局混合与局部保持之间的平衡，这增加了模型的灵活性。
\end{enumerate}

在实际应用中，(\beta)
可以作为超参数调整，也可以作为可学习参数，让模型自适应地决定信息混合的程度。

\hypertarget{ux590dux6570ux9149ux63a8ux7406ux6a21ux5f0f}{%
\subsection{6.7
复数酉推理模式}\label{ux590dux6570ux9149ux63a8ux7406ux6a21ux5f0f}}

\hypertarget{ux590dux6570ux52a8ux529bux5b66ux7684ux5f15ux5165}{%
\subsubsection{6.7.1
复数动力学的引入}\label{ux590dux6570ux52a8ux529bux5b66ux7684ux5f15ux5165}}

MQR支持复数推理模式，其中环形动力学直接在复数域进行：

{[} h\^{}\{(t+1)\} = h\^{}\{(t)\} U\^{}\dagger {]}

这里 (h \in C\^{}N) 是复数状态向量，(U)
是学习到的酉矩阵。这种模式充分利用了酉矩阵在复数域的性质。

\hypertarget{ux76f8ux4f4dux4fe1ux606fux7684ux4f5cux7528}{%
\subsubsection{6.7.2
相位信息的作用}\label{ux76f8ux4f4dux4fe1ux606fux7684ux4f5cux7528}}

在复数模式下，状态向量 (h)
不仅包含幅度信息，还包含相位信息。相位可以编码额外的信息，如振荡模式或干涉效应。读出时，通常取绝对值：

{[} y = W\_\{\text{readout}\} \cdot \textbar h\^{}*\textbar{} {]}

其中 (h\^{}*) 是平衡态，(\textbar{}\cdot\textbar)
是逐元素取模。这种设计让相位参与推理过程，但最终输出只依赖于幅度，确保了输出的实数值性。

\hypertarget{ux4e0eux91cfux5b50ux7cfbux7edfux7684ux7c7bux6bd4}{%
\subsubsection{6.7.3
与量子系统的类比}\label{ux4e0eux91cfux5b50ux7cfbux7edfux7684ux7c7bux6bd4}}

复数酉推理模式与量子系统的演化有直接类比： - 状态向量 (h)
类似于量子态矢量 - 酉矩阵 (U) 类似于时间演化算子 - 绝对值平方操作
(\textbar h\textbar\^{}2) 类似于量子测量得到概率分布

这种类比不仅提供了直观理解，还可能启发了未来将MQR扩展到量子机器学习领域的可能性。

\hypertarget{ux6570ux503cux8003ux8651}{%
\subsubsection{6.7.4 数值考虑}\label{ux6570ux503cux8003ux8651}}

在实现复数动力学时，需要注意数值稳定性问题。虽然酉变换理论上保持范数，但数值误差可能累积。MQR实现中采用了以下策略：
1. 定期重新正交化（可选） 2. 使用高精度浮点数 3. 监控状态向量的范数变化

\hypertarget{ux6570ux5b66ux6027ux8d28ux603bux7ed3ux4e0eux8bc1ux660e}{%
\subsection{6.8
数学性质总结与证明}\label{ux6570ux5b66ux6027ux8d28ux603bux7ed3ux4e0eux8bc1ux660e}}

\hypertarget{ux4e3bux8981ux5b9aux7406ux4e0eux8bc1ux660e}{%
\subsubsection{6.8.1
主要定理与证明}\label{ux4e3bux8981ux5b9aux7406ux4e0eux8bc1ux660e}}

\textbf{定理6.1（双随机性保持）}：设 (U \in U(N)) 是酉矩阵，(H =
\textbar U\textbar\^{}2) 是逐元素绝对值平方得到的矩阵，则 (H)
是双随机矩阵。

\textbf{证明}： 对于任意 (i \in \{1, \ldots, N\})，行和为： {[}
\sum\emph{\{j=1\}\^{}N h}\{ij\} = \sum\emph{\{j=1\}\^{}N
\textbar u}\{ij\}\textbar\^{}2 = \sum\emph{\{j=1\}\^{}N u}\{ij\}
\overline{u_{ij}} = (UU\^{}\dagger)\emph{\{ii\} {]} 由于 (U)
是酉矩阵，(UU\^{}\dagger = I)，所以 ((UU\^{}\dagger)}\{ii\} =
1)。因此，所有行和均为1。

类似地，对于任意 (j \in \{1, \ldots, N\})，列和为： {[}
\sum\emph{\{i=1\}\^{}N h}\{ij\} = \sum\emph{\{i=1\}\^{}N
\textbar u}\{ij\}\textbar\^{}2 = \sum\emph{\{i=1\}\^{}N u}\{ij\}
\overline{u_{ij}} = (U\^{}\dagger U)\_\{jj\} = 1 {]} 因此，所有列和也为

\hypertarget{ux7b2c7ux7ae0-ux5f1bux8c6bux52a8ux529bux5b66013ux677eux5f1bux7b97ux6cd5ux4e0eux6536ux655bux6027ux8bc1ux660e}{%
\section{第7章
弛豫动力学：013松弛算法与收敛性证明}\label{ux7b2c7ux7ae0-ux5f1bux8c6bux52a8ux529bux5b66013ux677eux5f1bux7b97ux6cd5ux4e0eux6536ux655bux6027ux8bc1ux660e}}

\hypertarget{ux5f15ux8a00ux4eceux987aux5e8fux8ba1ux7b97ux5230ux5f1bux8c6bux5e73ux8861}{%
\subsection{7.1
引言：从顺序计算到弛豫平衡}\label{ux5f15ux8a00ux4eceux987aux5e8fux8ba1ux7b97ux5230ux5f1bux8c6bux5e73ux8861}}

在传统的深度神经网络，特别是循环神经网络中，信息处理遵循着严格的前向传播范式：输入数据按时间步或网络层顺序流过一系列确定的变换函数，最终产生输出。这种计算图模型在工程实现上直观且高效，但其本质是一种\textbf{开环的、单向的}信息流动过程。然而，自然界和许多复杂系统中普遍存在着另一种计算模式：\textbf{弛豫动力学}。

弛豫动力学描述的是一个系统在外部驱动下，通过内部单元间的持续相互作用，逐渐演化到一个稳定平衡态的过程。这个平衡态由系统的内在结构和外部输入共同决定。例如，热力学系统最终达到温度均匀的热平衡，Hopfield网络通过能量最小化收敛到记忆模式，生物神经网络中的全局活动模式也在外部刺激下趋于稳定。

Möbius Quantum
Ring（MQR）网络的核心创新之一，正是将深度学习的推理过程重新构想为这样一种弛豫动力学。本章重点分析的\texttt{013}松弛动力学模块，是实现这一哲学转变的技术核心。它定义了一个离散时间的迭代过程，其中环形网络的状态在外部注入信号的驱动下，通过双随机连接矩阵进行局部扩散，最终收敛到唯一的稳态。

\hypertarget{ux677eux5f1bux52a8ux529bux5b66ux7684ux6570ux5b66ux8868ux8ff0}{%
\subsection{7.2
013松弛动力学的数学表述}\label{ux677eux5f1bux52a8ux529bux5b66ux7684ux6570ux5b66ux8868ux8ff0}}

\hypertarget{ux57faux672cux52a8ux529bux5b66ux65b9ux7a0b-2}{%
\subsubsection{7.2.1
基本动力学方程}\label{ux57faux672cux52a8ux529bux5b66ux65b9ux7a0b-2}}

设环形网络由\(N\)个节点组成，每个节点在时间步\(t\)的状态表示为行向量\(\textbf{h}^{(t)} \in R^{1 \times N}\)（实数版本）或\(C^{1 \times N}\)（复数版本）。网络内部的连接结构由双随机矩阵\(H \in R^{N \times N}\)描述，满足：
1. \(H_{ij} \geq 0\)（非负性） 2.
\(\sum_{i=1}^N H_{ij} = 1\)（列随机性） 3.
\(\sum_{j=1}^N H_{ij} = 1\)（行随机性）

外部输入通过注入算子\(J(\textbf{x})\)映射为注入向量\(\textbf{j} \in R^{1 \times N}\)，其中\(\textbf{x}\)是原始输入特征。

\texttt{013}松弛动力学的核心迭代方程定义为：

\[
\textbf{h}^{(t+1)} = (1 - \alpha) \textbf{h}^{(t)} H^T + \alpha \textbf{j}
\]

其中\(\alpha \in (0, 1]\)是\textbf{弛豫速率参数}，控制着外部注入与内部状态扩散的相对权重。当\(\alpha=1\)时，系统完全由外部输入驱动，立即达到\(\textbf{h}^{(1)} = \textbf{j}\)；当\(\alpha\)接近0时，系统更依赖于内部状态的缓慢扩散。

\hypertarget{ux53ccux968fux673aux77e9ux9635ux7684ux5173ux952eux6027ux8d28}{%
\subsubsection{7.2.2
双随机矩阵的关键性质}\label{ux53ccux968fux673aux77e9ux9635ux7684ux5173ux952eux6027ux8d28}}

在MQR中，连接矩阵\(H\)并非任意双随机矩阵，而是通过酉矩阵\(U \in U(N)\)构造而来：

\[
H = |U|^2, \quad H_{ij} = |U_{ij}|^2
\]

这种构造方式赋予了\(H\)几个重要数学性质：

\textbf{性质7.1（特征值界限）}：对于通过\(H = |U|^2\)构造的双随机矩阵，其特征值\(\lambda_i\)满足\(|\lambda_i| \leq 1\)，且\(\lambda_1 = 1\)对应全1特征向量。

\textbf{证明}：由于\(H\)是双随机矩阵，根据Perron-Frobenius定理，其谱半径\(\rho(H) = 1\)，且1是其特征值。对于通过模平方构造的矩阵，可以证明所有特征值都在单位圆盘内。

\textbf{性质7.2（收缩性）}：对于任意向量\(\textbf{v} \in C^{1 \times N}\)，有\(\|\textbf{v} H\|_2 \leq \|\textbf{v}\|_2\)，当且仅当\(\textbf{v}\)是\(H\)的右特征向量且对应特征值模为1时取等号。

\textbf{性质7.3（稳态唯一性）}：如果\(H\)是遍历的（不可约且非周期），则存在唯一的稳态分布\(\boldsymbol{\pi}\)满足\(\boldsymbol{\pi} H = \boldsymbol{\pi}\)，且\(\boldsymbol{\pi}_i > 0\)对所有\(i\)成立。

\hypertarget{ux975eux7ebfux6027ux6269ux5c55}{%
\subsubsection{7.2.3 非线性扩展}\label{ux975eux7ebfux6027ux6269ux5c55}}

基本的线性动力学方程可以扩展为包含非线性激活函数的版本：

\[
\textbf{h}^{(t+1)} = \sigma\left((1 - \alpha) \textbf{h}^{(t)} H^T + \alpha \textbf{j}\right)
\]

其中\(\sigma(\cdot)\)是逐点应用的激活函数。为了保证收敛性，\(\sigma\)需要满足\textbf{1-Lipschitz}条件：

\[
\|\sigma(\textbf{x}) - \sigma(\textbf{y})\|_2 \leq \|\textbf{x} - \textbf{y}\|_2, \quad \forall \textbf{x}, \textbf{y}
\]

常见的1-Lipschitz激活函数包括： -
\textbf{Tanh}：\(\tanh(x)\)的导数的绝对值不超过1 -
\textbf{归一化函数}：\(\sigma(\textbf{x}) = \frac{\textbf{x}}{\max(1, \|\textbf{x}\|_2 / \tau)}\)，其中\(\tau > 0\)是缩放因子
- \textbf{逐元素裁剪}：\(\sigma(x) = \text{clip}(x, -c, c)\)

\hypertarget{ux6536ux655bux6027ux8bc1ux660e-1}{%
\subsection{7.3 收敛性证明}\label{ux6536ux655bux6027ux8bc1ux660e-1}}

\hypertarget{ux7ebfux6027ux60c5ux51b5ux7684ux6536ux655bux6027}{%
\subsubsection{7.3.1
线性情况的收敛性}\label{ux7ebfux6027ux60c5ux51b5ux7684ux6536ux655bux6027}}

\textbf{定理7.1（线性013动力学的收敛性）}：对于任意初始状态\(\textbf{h}^{(0)}\)，由线性方程\(\textbf{h}^{(t+1)} = (1 - \alpha) \textbf{h}^{(t)} H^T + \alpha \textbf{j}\)定义的迭代序列收敛到唯一固定点：

\[
\textbf{h}^* = \alpha \textbf{j} (I - (1-\alpha)H^T)^{-1}
\]

\textbf{证明}：我们将证明分为三个部分。

\textbf{第一部分：压缩映射的构造}

定义映射\(T: R^{1 \times N} \rightarrow R^{1 \times N}\)为：

\[
T(\textbf{h}) = (1-\alpha)\textbf{h}H^T + \alpha\textbf{j}
\]

对于任意两个状态\(\textbf{h}_1, \textbf{h}_2\)，有：

\[
\begin{aligned}
\|T(\textbf{h}_1) - T(\textbf{h}_2)\|_2 &= \|(1-\alpha)(\textbf{h}_1 - \textbf{h}_2)H^T\|_2 \\
&\leq (1-\alpha)\|\textbf{h}_1 - \textbf{h}_2\|_2 \cdot \|H^T\|_2
\end{aligned}
\]

其中\(\|H^T\|_2 = \sigma_{\max}(H)\)是\(H\)的最大奇异值。由于\(H\)是双随机矩阵，其奇异值不超过1，因此\(\|H^T\|_2 \leq 1\)。于是：

\[
\|T(\textbf{h}_1) - T(\textbf{h}_2)\|_2 \leq (1-\alpha)\|\textbf{h}_1 - \textbf{h}_2\|_2
\]

由于\(\alpha \in (0, 1]\)，有\(0 \leq 1-\alpha < 1\)，因此\(T\)是一个压缩映射，压缩因子为\(1-\alpha\)。

\textbf{第二部分：固定点的存在性与唯一性}

根据Banach不动点定理，在完备度量空间\((R^{1 \times N}, \|\cdot\|_2)\)上，压缩映射\(T\)存在唯一的不动点\(\textbf{h}^*\)满足\(T(\textbf{h}^*) = \textbf{h}^*\)。

直接求解不动点方程：

\[
\textbf{h}^* = (1-\alpha)\textbf{h}^* H^T + \alpha\textbf{j}
\]

整理得：

\[
\textbf{h}^* (I - (1-\alpha)H^T) = \alpha\textbf{j}
\]

由于\(I - (1-\alpha)H^T\)是可逆的（因为\((1-\alpha)H^T\)的谱半径小于1），可得：

\[
\textbf{h}^* = \alpha\textbf{j} (I - (1-\alpha)H^T)^{-1}
\]

\textbf{第三部分：收敛速率}

由压缩映射的性质，迭代序列以几何速率收敛：

\[
\|\textbf{h}^{(t)} - \textbf{h}^*\|_2 \leq (1-\alpha)^t \|\textbf{h}^{(0)} - \textbf{h}^*\|_2
\]

收敛所需的迭代次数与\(\log(1/\epsilon)/\log(1/(1-\alpha))\)成正比，其中\(\epsilon\)是期望的精度。

\hypertarget{ux975eux7ebfux6027ux60c5ux51b5ux7684ux6536ux655bux6027}{%
\subsubsection{7.3.2
非线性情况的收敛性}\label{ux975eux7ebfux6027ux60c5ux51b5ux7684ux6536ux655bux6027}}

\textbf{定理7.2（非线性013动力学的收敛性）}：如果激活函数\(\sigma\)是1-Lipschitz的，且\(\alpha \in (0, 1]\)，则非线性动力学\(\textbf{h}^{(t+1)} = \sigma\left((1-\alpha)\textbf{h}^{(t)}H^T + \alpha\textbf{j}\right)\)也收敛到唯一固定点。

\textbf{证明}：定义非线性映射\(T_\sigma(\textbf{h}) = \sigma\left((1-\alpha)\textbf{h}H^T + \alpha\textbf{j}\right)\)。对于任意\(\textbf{h}_1, \textbf{h}_2\)：

\[
\begin{aligned}
\|T_\sigma(\textbf{h}_1) - T_\sigma(\textbf{h}_2)\|_2 
&= \|\sigma\left((1-\alpha)\textbf{h}_1H^T + \alpha\textbf{j}\right) - \sigma\left((1-\alpha)\textbf{h}_2H^T + \alpha\textbf{j}\right)\|_2 \\
&\leq \|(1-\alpha)(\textbf{h}_1 - \textbf{h}_2)H^T\|_2 \quad (\text{1-Lipschitz性质}) \\
&\leq (1-\alpha)\|\textbf{h}_1 - \textbf{h}_2\|_2 \cdot \|H^T\|_2 \\
&\leq (1-\alpha)\|\textbf{h}_1 - \textbf{h}_2\|_2
\end{aligned}
\]

因此\(T_\sigma\)同样是压缩因子为\(1-\alpha\)的压缩映射，根据Banach不动点定理，存在唯一不动点。

\hypertarget{ux6536ux655bux6027ux7684ux7269ux7406ux89e3ux91ca}{%
\subsubsection{7.3.3
收敛性的物理解释}\label{ux6536ux655bux6027ux7684ux7269ux7406ux89e3ux91ca}}

013动力学的收敛性可以从随机游走的角度理解。将状态向量\(\textbf{h}\)视为概率分布（经过适当归一化），则迭代方程描述了一个带有重启的随机游走过程：

\begin{enumerate}
\def\labelenumi{\arabic{enumi}.}
\tightlist
\item
  \textbf{扩散项}\((1-\alpha)\textbf{h}H^T\)：表示当前分布按照转移矩阵\(H\)进行一步随机游走
\item
  \textbf{注入项}\(\alpha\textbf{j}\)：表示以概率\(\alpha\)从注入分布\(\textbf{j}\)重新开始
\end{enumerate}

这种``带有重启的随机游走''模型在PageRank算法中广泛应用，其收敛性是众所周知的。参数\(\alpha\)控制着重启概率，较大的\(\alpha\)意味着系统更容易``忘记''历史状态，更快地响应外部输入。

\hypertarget{ux52a8ux529bux5b66ux6a21ux5757ux7684ux5b9eux73b0}{%
\subsection{7.4
013动力学模块的实现}\label{ux52a8ux529bux5b66ux6a21ux5757ux7684ux5b9eux73b0}}

\hypertarget{ux7c7bux8bbeux8ba1ux4e0eux63a5ux53e3}{%
\subsubsection{7.4.1
类设计与接口}\label{ux7c7bux8bbeux8ba1ux4e0eux63a5ux53e3}}

在MQR代码库中，013动力学由\texttt{Dynamics013}类实现，其核心结构如下：

\begin{Shaded}
\begin{Highlighting}[]
\KeywordTok{class}\NormalTok{ Dynamics013(nn.Module):}
    \CommentTok{"""}
\CommentTok{    核心松弛动力学模块。}
\CommentTok{    实现迭代更新: h\_\{t+1\} = (1{-}alpha) * h\_t @ H.T + alpha * injection}
\CommentTok{    }
\CommentTok{    参数:}
\CommentTok{        alpha: 弛豫速率，控制注入强度 (0 \textless{} alpha \textless{}= 1)}
\CommentTok{        use\_complex: 是否使用复数计算}
\CommentTok{        nonlinearity: 非线性激活类型，如\textquotesingle{}tanh\textquotesingle{}, \textquotesingle{}relu\textquotesingle{}, \textquotesingle{}norm\textquotesingle{}或None}
\CommentTok{        norm\_scale: 归一化激活的缩放因子}
\CommentTok{        max\_iterations: 最大迭代次数}
\CommentTok{        tolerance: 收敛容差}
\CommentTok{        record\_trajectory: 是否记录状态轨迹（用于调试）}
\CommentTok{    """}
    
    \KeywordTok{def} \FunctionTok{\_\_init\_\_}\NormalTok{(}\VariableTok{self}\NormalTok{, alpha}\OperatorTok{=}\FloatTok{0.5}\NormalTok{, use\_complex}\OperatorTok{=}\VariableTok{False}\NormalTok{, }
\NormalTok{                 nonlinearity}\OperatorTok{=}\VariableTok{None}\NormalTok{, norm\_scale}\OperatorTok{=}\FloatTok{1.0}\NormalTok{,}
\NormalTok{                 max\_iterations}\OperatorTok{=}\DecValTok{100}\NormalTok{, tolerance}\OperatorTok{=}\FloatTok{1e{-}6}\NormalTok{,}
\NormalTok{                 record\_trajectory}\OperatorTok{=}\VariableTok{False}\NormalTok{):}
        \BuiltInTok{super}\NormalTok{().}\FunctionTok{\_\_init\_\_}\NormalTok{()}
        \VariableTok{self}\NormalTok{.alpha }\OperatorTok{=}\NormalTok{ alpha}
        \VariableTok{self}\NormalTok{.use\_complex }\OperatorTok{=}\NormalTok{ use\_complex}
        \VariableTok{self}\NormalTok{.nonlinearity }\OperatorTok{=}\NormalTok{ nonlinearity}
        \VariableTok{self}\NormalTok{.norm\_scale }\OperatorTok{=}\NormalTok{ norm\_scale}
        \VariableTok{self}\NormalTok{.max\_iterations }\OperatorTok{=}\NormalTok{ max\_iterations}
        \VariableTok{self}\NormalTok{.tolerance }\OperatorTok{=}\NormalTok{ tolerance}
        \VariableTok{self}\NormalTok{.record\_trajectory }\OperatorTok{=}\NormalTok{ record\_trajectory}
        
        \CommentTok{\# 根据nonlinearity参数选择激活函数}
        \ControlFlowTok{if}\NormalTok{ nonlinearity }\OperatorTok{==} \StringTok{\textquotesingle{}tanh\textquotesingle{}}\NormalTok{:}
            \VariableTok{self}\NormalTok{.activation }\OperatorTok{=}\NormalTok{ nn.Tanh()}
        \ControlFlowTok{elif}\NormalTok{ nonlinearity }\OperatorTok{==} \StringTok{\textquotesingle{}norm\textquotesingle{}}\NormalTok{:}
            \VariableTok{self}\NormalTok{.activation }\OperatorTok{=} \VariableTok{self}\NormalTok{.\_norm\_activation}
        \ControlFlowTok{elif}\NormalTok{ nonlinearity }\KeywordTok{is} \VariableTok{None}\NormalTok{:}
            \VariableTok{self}\NormalTok{.activation }\OperatorTok{=} \KeywordTok{lambda}\NormalTok{ x: x}
        \ControlFlowTok{else}\NormalTok{:}
            \ControlFlowTok{raise} \PreprocessorTok{ValueError}\NormalTok{(}\SpecialStringTok{f"不支持的nonlinearity: }\SpecialCharTok{\{}\NormalTok{nonlinearity}\SpecialCharTok{\}}\SpecialStringTok{"}\NormalTok{)}
    
    \KeywordTok{def}\NormalTok{ \_norm\_activation(}\VariableTok{self}\NormalTok{, x):}
        \CommentTok{"""1{-}Lipschitz归一化激活函数"""}
        \ControlFlowTok{if} \VariableTok{self}\NormalTok{.use\_complex:}
\NormalTok{            norm }\OperatorTok{=}\NormalTok{ torch.sqrt(torch.real(x }\OperatorTok{*}\NormalTok{ x.conj()).}\BuiltInTok{sum}\NormalTok{(dim}\OperatorTok{={-}}\DecValTok{1}\NormalTok{, keepdim}\OperatorTok{=}\VariableTok{True}\NormalTok{) }\OperatorTok{+} \FloatTok{1e{-}8}\NormalTok{)}
        \ControlFlowTok{else}\NormalTok{:}
\NormalTok{            norm }\OperatorTok{=}\NormalTok{ torch.norm(x, dim}\OperatorTok{={-}}\DecValTok{1}\NormalTok{, keepdim}\OperatorTok{=}\VariableTok{True}\NormalTok{)}
\NormalTok{        scale }\OperatorTok{=}\NormalTok{ torch.minimum(torch.ones\_like(norm), }
                             \VariableTok{self}\NormalTok{.norm\_scale }\OperatorTok{/}\NormalTok{ (norm }\OperatorTok{+} \FloatTok{1e{-}8}\NormalTok{))}
        \ControlFlowTok{return}\NormalTok{ x }\OperatorTok{*}\NormalTok{ scale}
    
    \KeywordTok{def}\NormalTok{ forward(}\VariableTok{self}\NormalTok{, H, injection, h0}\OperatorTok{=}\VariableTok{None}\NormalTok{):}
        \CommentTok{"""}
\CommentTok{        执行松弛动力学直到收敛}
\CommentTok{        }
\CommentTok{        参数:}
\CommentTok{            H: 连接矩阵 [batch\_size, N, N] 或 [N, N]}
\CommentTok{            injection: 注入向量 [batch\_size, N] 或 [N]}
\CommentTok{            h0: 初始状态，如果为None则使用injection}
\CommentTok{            }
\CommentTok{        返回:}
\CommentTok{            h\_star: 收敛后的稳态 [batch\_size, N]}
\CommentTok{            info: 包含收敛信息的字典}
\CommentTok{        """}
        \CommentTok{\# 确保输入维度一致}
        \ControlFlowTok{if}\NormalTok{ H.dim() }\OperatorTok{==} \DecValTok{2}\NormalTok{:}
\NormalTok{            H }\OperatorTok{=}\NormalTok{ H.unsqueeze(}\DecValTok{0}\NormalTok{)  }\CommentTok{\# [1, N, N]}
\NormalTok{            injection }\OperatorTok{=}\NormalTok{ injection.unsqueeze(}\DecValTok{0}\NormalTok{)  }\CommentTok{\# [1, N]}
        
\NormalTok{        batch\_size, N, \_ }\OperatorTok{=}\NormalTok{ H.shape}
        
        \CommentTok{\# 初始化状态}
        \ControlFlowTok{if}\NormalTok{ h0 }\KeywordTok{is} \VariableTok{None}\NormalTok{:}
\NormalTok{            h }\OperatorTok{=}\NormalTok{ injection.clone()}
        \ControlFlowTok{else}\NormalTok{:}
\NormalTok{            h }\OperatorTok{=}\NormalTok{ h0.clone()}
        
        \CommentTok{\# 记录轨迹（如果启用）}
        \ControlFlowTok{if} \VariableTok{self}\NormalTok{.record\_trajectory:}
\NormalTok{            trajectory }\OperatorTok{=}\NormalTok{ [h.clone()]}
        
        \CommentTok{\# 迭代求解}
\NormalTok{        converged }\OperatorTok{=} \VariableTok{False}
        \ControlFlowTok{for}\NormalTok{ i }\KeywordTok{in} \BuiltInTok{range}\NormalTok{(}\VariableTok{self}\NormalTok{.max\_iterations):}
\NormalTok{            h\_prev }\OperatorTok{=}\NormalTok{ h.clone()}
            
            \CommentTok{\# 核心更新步骤: h = (1{-}alpha) * h @ H\^{}T + alpha * injection}
            \ControlFlowTok{if} \VariableTok{self}\NormalTok{.use\_complex:}
                \CommentTok{\# 复数版本}
\NormalTok{                h\_next }\OperatorTok{=}\NormalTok{ (}\DecValTok{1} \OperatorTok{{-}} \VariableTok{self}\NormalTok{.alpha) }\OperatorTok{*}\NormalTok{ torch.matmul(h, H.transpose(}\OperatorTok{{-}}\DecValTok{1}\NormalTok{, }\OperatorTok{{-}}\DecValTok{2}\NormalTok{).conj()) }\OperatorTok{+} \OperatorTok{\textbackslash{}}
                         \VariableTok{self}\NormalTok{.alpha }\OperatorTok{*}\NormalTok{ injection}
            \ControlFlowTok{else}\NormalTok{:}
                \CommentTok{\# 实数版本}
\NormalTok{                h\_next }\OperatorTok{=}\NormalTok{ (}\DecValTok{1} \OperatorTok{{-}} \VariableTok{self}\NormalTok{.alpha) }\OperatorTok{*}\NormalTok{ torch.matmul(h, H.transpose(}\OperatorTok{{-}}\DecValTok{1}\NormalTok{, }\OperatorTok{{-}}\DecValTok{2}\NormalTok{)) }\OperatorTok{+} \OperatorTok{\textbackslash{}}
                         \VariableTok{self}\NormalTok{.alpha }\OperatorTok{*}\NormalTok{ injection}
            
            \CommentTok{\# 应用非线性激活}
\NormalTok{            h }\OperatorTok{=} \VariableTok{self}\NormalTok{.activation(h\_next)}
            
            \CommentTok{\# 记录轨迹}
            \ControlFlowTok{if} \VariableTok{self}\NormalTok{.record\_trajectory:}
\NormalTok{                trajectory.append(h.clone())}
            
            \CommentTok{\# 检查收敛性}
\NormalTok{            diff }\OperatorTok{=}\NormalTok{ torch.norm(h }\OperatorTok{{-}}\NormalTok{ h\_prev, dim}\OperatorTok{={-}}\DecValTok{1}\NormalTok{).}\BuiltInTok{max}\NormalTok{().item()}
            \ControlFlowTok{if}\NormalTok{ diff }\OperatorTok{\textless{}} \VariableTok{self}\NormalTok{.tolerance:}
\NormalTok{                converged }\OperatorTok{=} \VariableTok{True}
                \ControlFlowTok{break}
        
        \CommentTok{\# 准备返回信息}
\NormalTok{        info }\OperatorTok{=}\NormalTok{ \{}
            \StringTok{\textquotesingle{}converged\textquotesingle{}}\NormalTok{: converged,}
            \StringTok{\textquotesingle{}iterations\textquotesingle{}}\NormalTok{: i }\OperatorTok{+} \DecValTok{1}\NormalTok{,}
            \StringTok{\textquotesingle{}final\_diff\textquotesingle{}}\NormalTok{: diff }\ControlFlowTok{if} \StringTok{\textquotesingle{}diff\textquotesingle{}} \KeywordTok{in} \BuiltInTok{locals}\NormalTok{() }\ControlFlowTok{else} \BuiltInTok{float}\NormalTok{(}\StringTok{\textquotesingle{}inf\textquotesingle{}}\NormalTok{)}
\NormalTok{        \}}
        
        \ControlFlowTok{if} \VariableTok{self}\NormalTok{.record\_trajectory:}
\NormalTok{            info[}\StringTok{\textquotesingle{}trajectory\textquotesingle{}}\NormalTok{] }\OperatorTok{=}\NormalTok{ torch.stack(trajectory, dim}\OperatorTok{=}\DecValTok{1}\NormalTok{)}
        
        \ControlFlowTok{return}\NormalTok{ h.squeeze(}\DecValTok{0}\NormalTok{) }\ControlFlowTok{if}\NormalTok{ batch\_size }\OperatorTok{==} \DecValTok{1} \ControlFlowTok{else}\NormalTok{ h, info}
\end{Highlighting}
\end{Shaded}

\hypertarget{ux7b97ux6cd5ux6d41ux7a0bux56fe}{%
\subsubsection{7.4.2 算法流程图}\label{ux7b97ux6cd5ux6d41ux7a0bux56fe}}

以下Mermaid流程图展示了013动力学模块的前向传播过程：

\begin{figure}
\centering
\includegraphics{figures/diagram_6.png}
\caption{图 6}
\end{figure}

\hypertarget{ux5b9eux73b0ux7ec6ux8282ux4e0eux4f18ux5316}{%
\subsubsection{7.4.3
实现细节与优化}\label{ux5b9eux73b0ux7ec6ux8282ux4e0eux4f18ux5316}}

\textbf{批量处理}：\texttt{Dynamics013}支持批量处理，可以同时计算多个样本的稳态。这是通过利用PyTorch的广播机制和批量矩阵乘法实现的。

\textbf{复数支持}：当\texttt{use\_complex=True}时，模块使用复数运算。连接矩阵\(H\)可以是复数酉矩阵的模平方，注入和状态也可以是复数。在更新步骤中，需要注意使用共轭转置\texttt{H.conj().transpose(-1,\ -2)}。

\textbf{收敛性监控}：实现中包含了收敛性检查，当状态变化小于预设容差\texttt{tolerance}时提前终止迭代，提高计算效率。同时记录实际迭代次数，便于调试和分析。

\textbf{数值稳定性}：在归一化激活函数中，添加了小常数\texttt{1e-8}防止除零错误。对于复数模的计算，使用\texttt{torch.sqrt(torch.real(x\ *\ x.conj())\ +\ 1e-8)}确保数值稳定性。

\hypertarget{ux52a8ux529bux5b66ux884cux4e3aux7684ux53efux89c6ux5316ux5206ux6790}{%
\subsection{7.5
动力学行为的可视化分析}\label{ux52a8ux529bux5b66ux884cux4e3aux7684ux53efux89c6ux5316ux5206ux6790}}

\hypertarget{ux72b6ux6001ux6f14ux5316ux8f68ux8ff9}{%
\subsubsection{7.5.1
状态演化轨迹}\label{ux72b6ux6001ux6f14ux5316ux8f68ux8ff9}}

为了直观理解013动力学的收敛过程，我们可以可视化状态向量的演化轨迹。考虑一个简单的3节点系统，连接矩阵为：

\$\$ H = \textbackslash begin\{bmatrix\} 0.6 \& 0.2 \& 0.2
\textbackslash{} 0.3 \& 0.4 \& 0

\hypertarget{ux7b2c8ux7ae0-ux83abux6bd4ux4e4cux65afux91cfux5b50ux73afux7f51ux7edcux7684ux52a8ux529bux5b66ux7cfbux7edfux4eceux7406ux8bbaux5230ux5b9eux73b0}{%
\section{第8章
莫比乌斯量子环网络的动力学系统：从理论到实现}\label{ux7b2c8ux7ae0-ux83abux6bd4ux4e4cux65afux91cfux5b50ux73afux7f51ux7edcux7684ux52a8ux529bux5b66ux7cfbux7edfux4eceux7406ux8bbaux5230ux5b9eux73b0}}

\hypertarget{ux52a8ux529bux5b66ux6838ux5fc3ux56faux5b9aux70b9ux677eux5f1bux4e0eux6536ux655bux6027ux4fddux8bc1}{%
\subsection{8.1
动力学核心：固定点松弛与收敛性保证}\label{ux52a8ux529bux5b66ux6838ux5fc3ux56faux5b9aux70b9ux677eux5f1bux4e0eux6536ux655bux6027ux4fddux8bc1}}

莫比乌斯量子环网络的核心创新在于其独特的推理机制------固定点松弛动力学。与传统的前馈神经网络或标准循环神经网络不同，MQR不采用单向的、确定性的计算图传播方式，而是将推理过程建模为一个动态系统达到平衡态的过程。这一设计哲学源于物理系统中的弛豫现象和分布式计算中的共识算法，为神经网络提供了全新的计算范式。

\hypertarget{ux56faux5b9aux70b9ux677eux5f1bux7684ux57faux672cux539fux7406}{%
\subsubsection{8.1.1
固定点松弛的基本原理}\label{ux56faux5b9aux70b9ux677eux5f1bux7684ux57faux672cux539fux7406}}

固定点松弛的核心思想是：给定一个输入信号，系统通过迭代更新其内部状态，最终收敛到一个唯一的稳态，这个稳态包含了输入信息经过系统动力学处理后的表示。在MQR中，这一过程可以形式化地描述为：

设环形网络有N个节点，每个时间步t的状态表示为向量 ( h\^{}t
\in R\^{}N )（实数模式）或 ( C\^{}N
)（复数模式）。系统的动力学由以下迭代方程定义：

{[} h\^{}\{t+1\} = F(h\^{}t, J(x)) {]}

其中 ( J(x) ) 是输入x通过注入模块得到的驱动信号，( F
) 是状态转移函数。当系统达到稳态时，满足：

{[} h\^{}* = F(h\^{}*, J(x)) {]}

这个方程的解 ( h\^{}* ) 就是系统的固定点。MQR的关键在于设计 (
F ) 使其满足压缩映射的条件，从而保证对于任意初始状态 ( h\^{}0
) 和任意输入 ( J(x) )，迭代过程都能收敛到唯一的固定点。

\hypertarget{ux5e7aux6a21ux968fux673aux8fdeux63a5ux77e9ux9635ux7684ux6570ux5b66ux6027ux8d28}{%
\subsubsection{8.1.2
幺模随机连接矩阵的数学性质}\label{ux5e7aux6a21ux968fux673aux8fdeux63a5ux77e9ux9635ux7684ux6570ux5b66ux6027ux8d28}}

MQR使用幺模随机矩阵作为连接矩阵，这是保证收敛性的数学基础。幺模随机矩阵
( H = \textbar U\textbar\^{}2 ) 具有以下重要性质：

\begin{enumerate}
\def\labelenumi{\arabic{enumi}.}
\item
  \textbf{双随机性}：( H ) 的所有元素非负，且每行和每列的和都为1： {[}
  \sum\emph{\{j=1\}\^{}N H}\{ij\} = 1 \quad \forall i,
  \quad \sum\emph{\{i=1\}\^{}N H}\{ij\} = 1 \quad \forall j {]}
\item
  \textbf{最大特征值为1}：作为双随机矩阵，( H ) 的谱半径 ( \rho(H) = 1
  )，且1是其特征值，对应的特征向量是全1向量（在适当归一化下）。
\item
  \textbf{保范性}：当 ( H ) 作用于概率分布向量时，保持向量的L1范数不变。
\end{enumerate}

这些性质共同确保了以 ( H )
为核心的线性扩散过程是稳定的、保能量的。在MQR的动力学中，线性扩散部分 (
hH\^{}T ) 不会放大或过度衰减信号，为整个迭代过程的收敛性提供了基础保证。

\hypertarget{ux52a8ux529bux5b66ux65b9ux7a0bux7684ux8be6ux7ec6ux63a8ux5bfcux4e0eux53d8ux4f53}{%
\subsection{8.2
动力学方程的详细推导与变体}\label{ux52a8ux529bux5b66ux65b9ux7a0bux7684ux8be6ux7ec6ux63a8ux5bfcux4e0eux53d8ux4f53}}

\hypertarget{ux57faux672cux52a8ux529bux5b66ux65b9ux7a0b-3}{%
\subsubsection{8.2.1
基本动力学方程}\label{ux57faux672cux52a8ux529bux5b66ux65b9ux7a0b-3}}

MQR的标准动力学方程采用线性混合形式：

{[} h\^{}\{t+1\} = (1-\alpha) h\^{}t H\^{}T + \alpha J(x) {]}

其中 ( \alpha \in (0,1{]} )
是混合系数，控制外部注入信号与内部状态扩散的相对权重。这个方程可以看作是两个过程的加权平均：
- 内部扩散：( h\^{}t H\^{}T ) 表示当前状态通过连接矩阵 ( H )
在环形网络上传播 - 外部注入：( J(x) )
表示输入信号对系统的持续驱动

\textbf{收敛性证明}： 定义映射 ( F(h) = (1-\alpha) h H\^{}T +
\alpha J(x) )。对于任意两个状态 ( h\_1, h\_2 )，有： {[}
\textbar{}F(h\_1) - F(h\_2)\textbar{} =
\textbar(1-\alpha)(h\_1 - h\_2) H\^{}T\textbar{} \leq (1-\alpha)
\textbar h\_1 - h\_2\textbar{} \textbar H\^{}T\textbar{} {]} 由于 ( H )
是双随机矩阵，其诱导的算子范数 ( \textbar H\^{}T\textbar\_\infty = 1
)（对于L\infty范数）或 ( \textbar H\^{}T\textbar\_2 \leq 1
)（对于L2范数）。因此： {[} \textbar{}F(h\_1) -
F(h\_2)\textbar{} \leq (1-\alpha) \textbar h\_1 -
h\_2\textbar{} {]} 由于 ( 0 \textless{} \alpha \leq 1 )，有 ( 0
\leq 1-\alpha \textless{} 1 )，所以 ( F )
是一个压缩映射。根据巴拿赫不动点定理，压缩映射在完备度量空间中有唯一的不动点，且从任意初始点出发的迭代都会收敛到该不动点。

\hypertarget{ux975eux7ebfux6027ux6269ux5c55-1}{%
\subsubsection{8.2.2
非线性扩展}\label{ux975eux7ebfux6027ux6269ux5c55-1}}

为了增强模型的表达能力，MQR支持在动力学中引入非线性激活函数：

{[} h\^{}\{t+1\} = \sigma\left((1-\alpha) h\^{}t H\^{}T +
\alpha J(x)\right) {]}

其中 ( \sigma ) 是逐点应用的激活函数。为了保证收敛性，( \sigma )
需要满足1-Lipschitz条件： {[} \textbar{}\sigma(x) - \sigma(y)\textbar{}
\leq \textbar x - y\textbar{} \quad \forall x,y {]} 常用的选择包括： -
ReLU激活：( \sigma(x) = \max(0, x) )，是1-Lipschitz的 - Tanh激活：(
\sigma(x) = \tanh(x) )，导数绝对值不超过1 - Sigmoid激活：( \sigma(x) =
1/(1+e\^{}\{-x\}) )，导数在(0, 0.25{]}范围内

在这种设置下，映射 ( F(h) = \sigma((1-\alpha) h H\^{}T +
\alpha J(x)) ) 仍然是压缩映射，因为： {[}
\textbar{}F(h\_1) - F(h\_2)\textbar{}
\leq \textbar(1-\alpha)(h\_1 - h\_2) H\^{}T\textbar{} \leq (1-\alpha)
\textbar h\_1 - h\_2\textbar{} {]} 这里用到了1-Lipschitz函数的性质 (
\textbar{}\sigma(a) - \sigma(b)\textbar{} \leq \textbar a - b\textbar{}
)。

\hypertarget{ux81eaux4fddux6301ux6df7ux5408ux53d8ux4f53}{%
\subsubsection{8.2.3
自保持混合变体}\label{ux81eaux4fddux6301ux6df7ux5408ux53d8ux4f53}}

为了抑制双随机矩阵可能导致的``过度平均化''效应，MQR引入了自保持混合技术：

{[} H\_\{\text{eff}\} = (1-\beta)I + \beta H {]}

其中 ( \beta \in [0,1] ) 控制自保持强度。相应的动力学方程为： {[}
h\^{}\{t+1\} = (1-\alpha) h\^{}t H\_\{\text{eff}\}\^{}T +
\alpha J(x) {]}

( H\_\{\text{eff}\} )
仍然是双随机矩阵，因为它是单位矩阵和双随机矩阵的凸组合。单位矩阵的引入使得每个节点在扩散过程中保留更多自身信息，减少了信息在环形网络上过度混合的风险。当
( \beta = 1 ) 时，退化为标准形式；当 ( \beta = 0 )
时，系统完全自保持，动力学简化为 ( h\^{}\{t+1\} = (1-\alpha) h\^{}t +
\alpha J(x) )。

\hypertarget{ux590dux6570ux9149ux63a8ux7406ux6a21ux5f0f-1}{%
\subsubsection{8.2.4
复数酉推理模式}\label{ux590dux6570ux9149ux63a8ux7406ux6a21ux5f0f-1}}

在复数模式下，MQR直接使用酉矩阵 ( U ) 进行推理，动力学方程为：

{[} h\^{}\{t+1\} = (1-\alpha) h\^{}t U\^{}\dagger +
\alpha J(x) {]}

其中 ( U\^{}\dagger ) 是 ( U ) 的共轭转置。由于酉矩阵保持L2范数不变（(
\textbar Ux\textbar\_2 = \textbar x\textbar\_2 )），因此 (
\textbar U\^{}\dagger\textbar\_2 = 1 )。收敛性证明与实数情况类似： {[}
\textbar{}F(h\_1) - F(h\_2)\textbar\_2 =
\textbar(1-\alpha)(h\_1 - h\_2) U\^{}\dagger\textbar\_2 \leq (1-\alpha)
\textbar h\_1 - h\_2\textbar\_2 {]}

复数模式的优势在于相位信息可以参与计算，提供了更丰富的表示空间。在读出时，通常取模值
( \textbar h\^{}*\textbar{} ) 或同时使用实部和虚部。

\hypertarget{ux6536ux655bux5224ux636eux4e0eux8fedux4ee3ux4f18ux5316}{%
\subsection{8.3
收敛判据与迭代优化}\label{ux6536ux655bux5224ux636eux4e0eux8fedux4ee3ux4f18ux5316}}

\hypertarget{ux6536ux655bux6761ux4ef6}{%
\subsubsection{8.3.1 收敛条件}\label{ux6536ux655bux6761ux4ef6}}

在实际实现中，需要判断迭代过程何时达到收敛。MQR采用相对变化量作为收敛判据：

{[} \frac{\|h^{t+1} - h^t\|}{\|h^t\| + \epsilon} \textless{} \tau {]}

其中： - ( \textbar{}\cdot\textbar{} ) 通常是L2范数 - ( \epsilon )
是小常数（如1e-8），防止除零 - ( \tau ) 是收敛阈值（如1e-6）

同时，为了避免无限循环，设置最大迭代次数 ( T\_\{\max\}
)（如1000）。当迭代次数达到 ( T\_\{\max\} )
仍未收敛时，使用当前状态作为近似稳态。

\hypertarget{ux6536ux655bux901fux5ea6ux5206ux6790}{%
\subsubsection{8.3.2
收敛速度分析}\label{ux6536ux655bux901fux5ea6ux5206ux6790}}

动力学方程的收敛速度由混合系数 ( \alpha ) 决定。从压缩因子 ( 1-\alpha )
可以看出，( \alpha )
越大，收敛越快。具体来说，经过t次迭代后，初始误差的衰减为：

{[} \textbar h\^{}t - h\^{}\emph{\textbar{} \leq (1-\alpha)\^{}t
\textbar h\^{}0 - h\^{}}\textbar{} {]}

因此，要达到误差小于 ( \delta ) 所需的迭代次数约为：

{[} t \geq \frac{\log(\delta / \|h^0 - h^*\|)}{\log(1-\alpha)} {]}

在实际应用中，需要权衡收敛速度和稳态特性。较大的 ( \alpha )
虽然加速收敛，但可能使系统过于依赖外部注入而减弱内部动力学的处理能力。

\hypertarget{ux521dux59cbux72b6ux6001ux9009ux62e9}{%
\subsubsection{8.3.3
初始状态选择}\label{ux521dux59cbux72b6ux6001ux9009ux62e9}}

初始状态 ( h\^{}0 )
的选择会影响收敛所需的迭代次数，但不影响最终的稳态（由于唯一性）。MQR提供了多种初始化策略：

\begin{enumerate}
\def\labelenumi{\arabic{enumi}.}
\tightlist
\item
  \textbf{零初始化}：( h\^{}0 = 0 )，最简单但可能不是最优
\item
  \textbf{随机初始化}：从均匀分布或正态分布采样
\item
  \textbf{注入信号初始化}：( h\^{}0 = J(x) )，利用输入信息
\item
  \textbf{历史状态初始化}：在序列处理中，使用前一个时间步的稳态作为当前初始值
\end{enumerate}

实验表明，使用注入信号初始化通常能减少约30\%的迭代次数，因为初始状态已经包含了输入信息。

\hypertarget{ux52a8ux529bux5b66ux7cfbux7edfux7684ux5b9eux73b0ux67b6ux6784}{%
\subsection{8.4
动力学系统的实现架构}\label{ux52a8ux529bux5b66ux7cfbux7edfux7684ux5b9eux73b0ux67b6ux6784}}

\begin{figure}
\centering
\includegraphics{figures/diagram_7.png}
\caption{图 7}
\end{figure}

\hypertarget{ux6838ux5fc3ux8fedux4ee3ux5faaux73afux7684ux5b9eux73b0}{%
\subsubsection{8.4.1
核心迭代循环的实现}\label{ux6838ux5fc3ux8fedux4ee3ux5faaux73afux7684ux5b9eux73b0}}

MQR动力学系统的核心实现是一个紧凑而高效的迭代循环。以下是伪代码描述：

\begin{Shaded}
\begin{Highlighting}[]
\KeywordTok{def}\NormalTok{ fixed\_point\_iteration(Jx, H, alpha}\OperatorTok{=}\FloatTok{0.5}\NormalTok{, max\_iter}\OperatorTok{=}\DecValTok{1000}\NormalTok{, tol}\OperatorTok{=}\FloatTok{1e{-}6}\NormalTok{):}
    \CommentTok{"""}
\CommentTok{    固定点迭代求解稳态}
\CommentTok{    }
\CommentTok{    参数:}
\CommentTok{        Jx: 注入信号，形状 [batch\_size, N]}
\CommentTok{        H: 连接矩阵，形状 [N, N]}
\CommentTok{        alpha: 混合系数}
\CommentTok{        max\_iter: 最大迭代次数}
\CommentTok{        tol: 收敛容差}
\CommentTok{    }
\CommentTok{    返回:}
\CommentTok{        h\_star: 稳态，形状 [batch\_size, N]}
\CommentTok{        n\_iter: 实际迭代次数}
\CommentTok{    """}
\NormalTok{    batch\_size, N }\OperatorTok{=}\NormalTok{ Jx.shape}
    
    \CommentTok{\# 初始化状态}
\NormalTok{    h }\OperatorTok{=}\NormalTok{ torch.zeros\_like(Jx)  }\CommentTok{\# 或使用其他初始化策略}
\NormalTok{    h\_prev }\OperatorTok{=}\NormalTok{ h.clone()}
    
    \ControlFlowTok{for}\NormalTok{ t }\KeywordTok{in} \BuiltInTok{range}\NormalTok{(max\_iter):}
        \CommentTok{\# 线性扩散部分}
\NormalTok{        diffusion }\OperatorTok{=}\NormalTok{ torch.matmul(h, H.T)  }\CommentTok{\# [batch, N] × [N, N]\^{}T}
        
        \CommentTok{\# 混合注入信号}
\NormalTok{        h\_new }\OperatorTok{=}\NormalTok{ (}\DecValTok{1} \OperatorTok{{-}}\NormalTok{ alpha) }\OperatorTok{*}\NormalTok{ diffusion }\OperatorTok{+}\NormalTok{ alpha }\OperatorTok{*}\NormalTok{ Jx}
        
        \CommentTok{\# 检查收敛}
\NormalTok{        diff\_norm }\OperatorTok{=}\NormalTok{ torch.norm(h\_new }\OperatorTok{{-}}\NormalTok{ h\_prev, dim}\OperatorTok{=}\DecValTok{1}\NormalTok{)}
\NormalTok{        prev\_norm }\OperatorTok{=}\NormalTok{ torch.norm(h\_prev, dim}\OperatorTok{=}\DecValTok{1}\NormalTok{) }\OperatorTok{+} \FloatTok{1e{-}8}
\NormalTok{        rel\_diff }\OperatorTok{=}\NormalTok{ diff\_norm }\OperatorTok{/}\NormalTok{ prev\_norm}
        
        \CommentTok{\# 更新状态}
\NormalTok{        h\_prev }\OperatorTok{=}\NormalTok{ h}
\NormalTok{        h }\OperatorTok{=}\NormalTok{ h\_new}
        
        \CommentTok{\# 判断是否收敛}
        \ControlFlowTok{if}\NormalTok{ torch.}\BuiltInTok{all}\NormalTok{(rel\_diff }\OperatorTok{\textless{}}\NormalTok{ tol):}
            \ControlFlowTok{return}\NormalTok{ h, t }\OperatorTok{+} \DecValTok{1}
    
    \CommentTok{\# 达到最大迭代次数}
    \ControlFlowTok{return}\NormalTok{ h, max\_iter}
\end{Highlighting}
\end{Shaded}

\hypertarget{ux6279ux5904ux7406ux4f18ux5316}{%
\subsubsection{8.4.2 批处理优化}\label{ux6279ux5904ux7406ux4f18ux5316}}

在实际训练中，需要同时处理多个样本（批处理）。MQR的实现充分利用了现代深度学习框架的并行计算能力：

\begin{enumerate}
\def\labelenumi{\arabic{enumi}.}
\tightlist
\item
  \textbf{矩阵乘法优化}：使用 \texttt{torch.matmul}
  进行批处理矩阵乘法，一次性计算所有样本的扩散过程
\item
  \textbf{内存效率}：通过原地操作和梯度检查点技术减少内存占用
\item
  \textbf{收敛判断向量化}：对批处理中的所有样本同时进行收敛判断，只有当所有样本都收敛时才停止迭代
\end{enumerate}

\hypertarget{ux68afux5ea6ux8ba1ux7b97ux4e0eux53cdux5411ux4f20ux64ad}{%
\subsubsection{8.4.3
梯度计算与反向传播}\label{ux68afux5ea6ux8ba1ux7b97ux4e0eux53cdux5411ux4f20ux64ad}}

固定点迭代过程涉及循环，这给反向传播带来了挑战。MQR采用两种策略处理梯度：

\begin{enumerate}
\def\labelenumi{\arabic{enumi}.}
\tightlist
\item
  \textbf{展开迭代}：在训练时，将迭代循环展开为计算图。这种方法梯度计算精确，但内存消耗与迭代次数成正比
\item
  \textbf{隐式微分}：利用固定点方程直接计算梯度，无需存储中间状态。对于稳态
  ( h\^{}* = F(h\^{}*, \theta) )，其中 ( \theta )
  是参数，梯度可以通过解以下线性系统得到： {[}
  \frac{\partial h^*}{\partial \theta} = \left(I -
  \frac{\partial F}{\partial h}\bigg\textbar\_\{h\^{}*\}\right)\^{}\{-1\}
  \frac{\partial F}{\partial \theta}\bigg\textbar\_\{h\^{}*\}
  {]} 这种方法内存效率高，但需要求解线性系统
\end{enumerate}

MQR默认使用展开迭代，因为实现简单且在现代GPU上通常可行。对于特别深的迭代或内存受限的情况，可以切换到隐式微分模式。

\hypertarget{ux52a8ux529bux5b66ux53c2ux6570ux7684ux5f71ux54cdux4e0eux8c03ux4f18}{%
\subsection{8.5
动力学参数的影响与调优}\label{ux52a8ux529bux5b66ux53c2ux6570ux7684ux5f71ux54cdux4e0eux8c03ux4f18}}

\hypertarget{ux6df7ux5408ux7cfbux6570-ux3b1-ux7684ux4f5cux7528}{%
\subsubsection{8.5.1 混合系数 \alpha
的作用}\label{ux6df7ux5408ux7cfbux6570-ux3b1-ux7684ux4f5cux7528}}

混合系数 \alpha
是动力学系统中最重要的超参数之一，它控制着外部注入与内部扩散的平衡：

\begin{itemize}
\tightlist
\item
  \textbf{\alpha \rightarrow
  0}：系统主要依赖内部扩散，外部注入影响微弱。收敛速度慢，但稳态更多地反映了环形网络的固有动力学特性
\item
  \textbf{\alpha \rightarrow
  1}：系统直接受外部注入驱动，内部扩散影响小。收敛速度快，但稳态接近注入信号，可能损失环形网络的变换能力
\item
  \textbf{\alpha = 0.5}：平衡点，通常作为默认值
\end{itemize}

实验表明，对于不同任务，最优的 \alpha 值不同： - 序列建模任务：较小的
\alpha（0.1-0.3）有利于捕捉长期依赖 - 分类任务：较大的
\alpha（0.5-0.7）有利于快速响应输入特征 - 生成任务：中等
\alpha（0.3-0.5）平衡了创造性和一致性

\hypertarget{ux8fdeux63a5ux77e9ux9635ux7684ux8c31ux7279ux6027}{%
\subsubsection{8.5.2
连接矩阵的谱特性}\label{ux8fdeux63a5ux77e9ux9635ux7684ux8c31ux7279ux6027}}

连接矩阵 H 的谱特性直接影响动力学的行为。幺模随机矩阵 H =
\textbar U\textbar² 的谱具有以下特点：

\begin{enumerate}
\def\labelenumi{\arabic{enumi}.}
\item
  \textbf{特征值分布}：所有特征值位于单位圆内，最大特征值为1。特征值的分布决定了信息混合的速度：

  \begin{itemize}
  \tightlist
  \item
    特征值接近1：慢混合模式，信息在环形网络中持久循环
  \item
    特征值远离1：快混合模式，信息迅速扩散到整个网络
  \end{itemize}
\item
  \textbf{特征向量结构}：特征向量反映了环形网络中的振动模式。在图像处理任务中，这些模式可以对应不同频率的空间特征
\item
  \textbf{条件数}：酉矩阵 U 的条件数为1，这保证了数值稳定性
\end{enumerate}

\hypertarget{ux975eux7ebfux6027ux6fc0ux6d3bux7684ux9009ux62e9}{%
\subsubsection{8.5.3
非线性激活的选择}\label{ux975eux7ebfux6027ux6fc0ux6d3bux7684ux9009ux62e9}}

非线性激活函数 \sigma 为动力学系统引入了必要的非线性变换能力：

\begin{enumerate}
\def\labelenumi{\arabic{enumi}.}
\tightlist
\item
  \textbf{ReLU}：最常用的选择，计算简单，梯度稳定。但可能引起神经元``死亡''问题
\item
  \textbf{Tanh}：输出范围(-1,1)，对称性好。在复数模式下，可以使用复Tanh
\item
  \textbf{Sigmoid}：输出范围(0,1)，适合表示概率或注意力权重
\item
  \textbf{GELU}：高斯误差线性单元，结合了ReLU和Dropout的思想，性能优秀但计算稍复杂
\item
  \textbf{Swish}：自门控激活，( \text{Swish}(x) = x
  \cdot \text{sigmoid}(x) )，性能优异
\end{enumerate}

在MQR中，激活函数需要满足1-Lipschitz条件以保证收敛性。实践中，可以通过适当的缩放（如使用权重归一化）使标准激活函数满足这一条件。

\hypertarget{ux9ad8ux7ea7ux52a8ux529bux5b66ux7279ux6027}{%
\subsection{8.6
高级动力学特性}\label{ux9ad8ux7ea7ux52a8ux529bux5b66ux7279ux6027}}

\hypertarget{ux591aux65f6ux95f4ux5c3aux5ea6ux52a8ux529bux5b66}{%
\subsubsection{8.6.1
多时间尺度动力学}\label{ux591aux65f6ux95f4ux5c3aux5ea6ux52a8ux529bux5b66}}

对于复杂任务，单一时间尺度的动力学可能不足。MQR

\hypertarget{ux7b2c9ux7ae0-ux6838ux5fc3ux52a8ux529bux5b66017ux56faux5b9aux70b9ux677eux5f1bux7b97ux6cd5}{%
\section{第9章
核心动力学：017固定点松弛算法}\label{ux7b2c9ux7ae0-ux6838ux5fc3ux52a8ux529bux5b66017ux56faux5b9aux70b9ux677eux5f1bux7b97ux6cd5}}

\hypertarget{ux5f15ux8a00ux4eceux8ba1ux7b97ux56feux5230ux5f1bux8c6bux7cfbux7edf}{%
\subsection{9.1
引言：从``计算图''到``弛豫系统''}\label{ux5f15ux8a00ux4eceux8ba1ux7b97ux56feux5230ux5f1bux8c6bux7cfbux7edf}}

在传统的深度神经网络，尤其是循环神经网络中，信息处理遵循着一种明确的、有向的``计算图''范式。输入数据从网络入口流入，按照预定的层级或时间步顺序，依次经过一系列参数化变换（线性层、非线性激活、门控机制等），最终在输出层产生预测结果。这种范式清晰、直观，并且与冯·诺依曼架构的计算模型高度契合。然而，它也在根本上将网络定义为一个\textbf{开环的、单向的、确定性的}信号处理器。

MöbiusQuantumRing (MQR)
架构提出了一种哲学上截然不同的计算范式。它不将网络视为一个前馈的计算图，而是将其建模为一个\textbf{闭环的、自洽的、通过弛豫达到平衡的动力学系统}。网络的``推理''过程，不再是执行一次前向传播，而是驱动一个环形动力系统从某个初始状态出发，通过反复应用一个固定的、双随机的传播规则，最终收敛到一个唯一的、稳定的平衡态。这个平衡态包含了输入信息经过系统全局相互作用后的``消化''结果，随后通过一个局部读出操作得到最终输出。

本章标题``017''正是对这一核心动力学过程的精炼概括： - \textbf{0}:
代表隐藏状态环的\textbf{零初始化}，或更广义的，动力系统迭代的起始点。 -
\textbf{1}:
代表每次迭代所应用的\textbf{单一、一致的传播算子}（幺模随机矩阵H或其变体）。
- \textbf{7}:
代表确保该迭代过程收敛到唯一固定点所需满足的\textbf{七个关键数学条件}。

本章将深入解析MQR网络的``心脏''------基于幺模随机矩阵的固定点松弛动力学。我们将从数学原理、收敛性证明、算法实现细节以及其与经典RNN范式的根本区别等多个维度进行阐述，揭示这种``无头无尾''的环形计算如何同时实现卓越的训练稳定性和丰富的表达能力。

\hypertarget{ux6570ux5b66ux57faux7840ux5e7aux6a21ux968fux673aux77e9ux9635ux4e0eux53ccux968fux673aux52a8ux529bux7cfbux7edf}{%
\subsection{9.2
数学基础：幺模随机矩阵与双随机动力系统}\label{ux6570ux5b66ux57faux7840ux5e7aux6a21ux968fux673aux77e9ux9635ux4e0eux53ccux968fux673aux52a8ux529bux7cfbux7edf}}

\hypertarget{ux4eceux9149ux77e9ux9635ux5230ux5e7aux6a21ux968fux673aux77e9ux9635-2}{%
\subsubsection{9.2.1
从酉矩阵到幺模随机矩阵}\label{ux4eceux9149ux77e9ux9635ux5230ux5e7aux6a21ux968fux673aux77e9ux9635-2}}

MQR动力学的基础是一个特殊的连接矩阵 ( H \in R\^{}\{N \times N\}
)，其中 ( N ) 是环形网络中节点的数量。( H )
并非直接学习得到，而是通过一个更基础的数学对象------\textbf{酉矩阵}------派生而来。

设 ( U \in C\^{}\{N \times N\} ) 是一个酉矩阵，满足 (
U\^{}\dagger U = U U\^{}\dagger = I )，其中 ( \dagger )
表示共轭转置。酉矩阵在复数域上作用时，保持向量的L2范数不变（等距性）。

\textbf{幺模随机矩阵} ( H ) 定义为酉矩阵 ( U ) 的逐元素模平方： {[} H =
\textbar U\textbar\^{}2, \quad \text{即} \quad H\_\{ij\} =
\textbar U\_\{ij\}\textbar\^{}2 {]}

这个简单的操作赋予了 ( H ) 两个至关重要的性质： 1. \textbf{非负性}：(
H\_\{ij\} \ge 0 ) 对所有 ( i, j ) 成立。 2. \textbf{双随机性}： {[}
\sum\emph{\{i=1\}\^{}\{N\} H}\{ij\} = 1 \quad \forall j, \quad \text{且}
\quad \sum\emph{\{j=1\}\^{}\{N\} H}\{ij\} = 1 \quad \forall i. {]}
双随机性源于酉矩阵的行和列向量均为单位范数向量。双随机矩阵的左右特征值1对应的特征向量分别是全1向量，这使其天然成为一个\textbf{概率转移矩阵}。

通过Cayley变换、指数映射或直接在Stiefel流形上优化等技术，我们可以有效地参数化并学习酉矩阵
( U )，从而间接地学习到具有严格双随机约束的连接矩阵 ( H
)。这种参数化方式将模型的搜索空间约束在了一个具有良好几何结构的流形上，为优化过程带来了内在的稳定性。

\hypertarget{ux73afux5f62ux52a8ux529bux5b66ux7684ux79bbux6563ux8fedux4ee3ux65b9ux7a0b}{%
\subsubsection{9.2.2
环形动力学的离散迭代方程}\label{ux73afux5f62ux52a8ux529bux5b66ux7684ux79bbux6563ux8fedux4ee3ux65b9ux7a0b}}

MQR网络的核心状态是一个在环形结构上定义的实值向量 ( \textbf{h}
\in R\^{}N )。在推理（前向传播）时，给定一个输入 ( \textbf{x} )
经过注入算子 ( J(\textbf{x}) ) 编码后的驱动信号 ( \textbf{j}
\in R\^{}N )，系统的动力学由以下离散迭代方程描述：

{[} \textbf{h}\^{}\{(t+1)\} = \sigma\left( (1-\alpha)
\textbf{h}\^{}\{(t)\} H\^{}T + \alpha \textbf{j} \right) {]} 其中： - (
t = 0, 1, 2, \dots ) 是迭代步数索引。 - ( \alpha \in (0, 1{]} )
是一个\textbf{阻尼因子}或\textbf{注入强度}，控制外部驱动与系统内部状态的比例。
- ( \sigma(\cdot) )
是一个逐元素的非线性激活函数，通常要求是\textbf{1-Lipschitz连续}的（例如
\texttt{tanh}, \texttt{clamp}, 或恒等映射 \texttt{identity}）。 -
初始状态 ( \textbf{h}\^{}\{(0)\} ) 通常设置为零向量 ( \textbf{0} )。

\textbf{物理诠释}： - ( \textbf{h}\^{}\{(t)\} H\^{}T )：表示当前状态 (
\textbf{h}\^{}\{(t)\} ) 通过双随机连接矩阵 ( H\^{}T )
在环形网络上进行了一次扩散或传播。由于 ( H ) 是双随机的，( H\^{}T )
也是双随机的，该操作可以看作一个保质量的混合过程。 - ( (1-\alpha)
\textbf{h}\^{}\{(t)\} H\^{}T + \alpha \textbf{j}
)：将上一步的内部传播结果与当前的外部驱动信号进行凸组合。( \alpha )
越小，系统越依赖于自身历史状态，惯性越大；( \alpha )
越大，系统对外部输入越敏感。 - ( \sigma(\cdot)
)：引入必要的非线性变换，以增强模型的表达能力。1-Lipschitz条件是为了保证整个迭代映射是一个收缩映射，这是收敛性的关键。

当 ( \sigma ) 为恒等映射时，迭代方程退化为线性系统： {[}
\textbf{h}\^{}\{(t+1)\} = (1-\alpha) \textbf{h}\^{}\{(t)\} H\^{}T +
\alpha \textbf{j} {]}
这个线性系统在满足一定条件下，其固定点可以直接通过求解线性方程组得到。但引入非线性
( \sigma )
后，系统能表征更复杂的函数，同时我们通过精心设计仍能保证其收敛性。

\hypertarget{ux6536ux655bux6027ux5206ux6790ux4e03ux4e2aux5173ux952eux6761ux4ef6ux4e0eux5df4ux62ffux8d6bux4e0dux52a8ux70b9ux5b9aux7406}{%
\subsection{9.3
收敛性分析：七个关键条件与巴拿赫不动点定理}\label{ux6536ux655bux6027ux5206ux6790ux4e03ux4e2aux5173ux952eux6761ux4ef6ux4e0eux5df4ux62ffux8d6bux4e0dux52a8ux70b9ux5b9aux7406}}

MQR动力学能够稳定工作的核心数学保证在于，其定义的迭代映射 ( F:
R\^{}N \to R\^{}N )
是一个\textbf{收缩映射}。根据完备度量空间上的\textbf{巴拿赫不动点定理}，如果一个映射
( F ) 是收缩的，那么它存在唯一的不动点 ( \textbf{h}\^{}* =
F(\textbf{h}\^{}*) )，并且从任意初始点开始迭代，序列 ( \{
\textbf{h}\^{}\{(t)\} \} ) 都会以指数速度收敛到该不动点。

下面我们详细阐述确保 ( F )
为收缩映射的七个关键条件（``017''中的``7''）：

\textbf{条件1：双随机性}。连接矩阵 ( H )（以及其转置 ( H\^{}T
)）是双随机的。这保证了 ( H\^{}T ) 作为线性算子的谱范数 (
\textbar H\^{}T\textbar\_2 = 1
)（因为最大奇异值为1）。双随机性是能量/质量在环上守恒的基础，防止了迭代过程中的幅值发散。

\textbf{条件2：凸组合系数约束}。阻尼因子 ( \alpha ) 满足 ( 0 \textless{}
\alpha \le 1 )。系数 ( (1-\alpha) ) 和 ( \alpha )
非负且和为1，这确保了迭代更新是历史状态和外部驱动的凸组合，是稳定混合的前提。

\textbf{条件3：线性部分的收缩性}。考虑线性部分算子 (
L(\textbf{h}) = (1-\alpha) \textbf{h} H\^{}T
)。对于任意两个状态 ( \textbf{h}\_1, \textbf{h}\_2 )，有： {[}
\textbar{} L(\textbf{h}\_1) - L(\textbf{h}\_2)
\textbar{} = \textbar{} (1-\alpha) (\textbf{h}\_1 - \textbf{h}\_2)
H\^{}T \textbar{} \le (1-\alpha) \textbar{} \textbf{h}\_1 -
\textbf{h}\_2 \textbar{} \cdot \textbar{} H\^{}T \textbar{} {]} 由于 (
\textbar{} H\^{}T \textbar\_2 = 1 )（条件1），因此 ( \textbar{}
L(\textbf{h}\_1) - L(\textbf{h}\_2) \textbar{}
\le (1-\alpha) \textbar{} \textbf{h}\_1 - \textbf{h}\_2 \textbar{}
)。这里 ( 0 \le (1-\alpha) \textless{} 1 )（当 ( \alpha \textgreater{} 0
)），所以线性部分本身就是一个收缩系数为 ( (1-\alpha) ) 的收缩映射。

\textbf{条件4：非线性激活的1-Lipschitz连续性}。激活函数 ( \sigma:
R \to R ) 满足 ( \textbar{}\sigma(a) -
\sigma(b)\textbar{} \le \textbar a - b\textbar{} ) 对所有实数 ( a, b )
成立。常见的 \texttt{tanh}, \texttt{sigmoid},
\texttt{ReLU}（实际上ReLU是1-Lipschitz）以及特殊的 \texttt{clamp}
函数（如 ( \text{clamp}(x, -c, c) )）都满足此条件。对于逐元素应用的 (
\sigma )，该条件意味着： {[} \textbar{} \sigma(\textbf{v}) -
\sigma(\textbf{w}) \textbar{} \le \textbar{} \textbf{v} - \textbf{w}
\textbar{} {]} 对于任意的向量 ( \textbf{v}, \textbf{w}
\in R\^{}N )（这里范数可取 ( \ell\_2 )
范数）。这个条件保证了非线性不会放大向量间的距离。

\textbf{条件5：驱动信号的独立性}。注入信号 ( \textbf{j} =
J(\textbf{x}) ) 在单次推理过程中是常数，不随迭代状态 (
\textbf{h}\^{}\{(t)\} ) 变化。因此，在比较两次迭代的差异时，(
\alpha \textbf{j} ) 项会相消。

\textbf{条件6：完整迭代映射的收缩性}。综合以上条件，考虑完整的非线性迭代映射
( F(\textbf{h}) = \sigma\left( (1-\alpha) \textbf{h} H\^{}T +
\alpha \textbf{j} \right) )。对于任意 ( \textbf{h}\_1, \textbf{h}\_2 )：
{[}

\begin{aligned}
\| F(\textbf{h}_1) - F(\textbf{h}_2) \|
&\le \| \left( (1-\alpha) \textbf{h}_1 H^T + \alpha \textbf{j} \right) - \left( (1-\alpha) \textbf{h}_2 H^T + \alpha \textbf{j} \right) \| \quad \text{(由条件4)} \\
&= \| (1-\alpha) (\textbf{h}_1 - \textbf{h}_2) H^T \| \\
&\le (1-\alpha) \| \textbf{h}_1 - \textbf{h}_2 \| \cdot \| H^T \| \quad \text{(范数的次可乘性)} \\
&= (1-\alpha) \| \textbf{h}_1 - \textbf{h}_2 \| \quad \text{(由条件1)}.
\end{aligned}

{]} 由于 ( 0 \textless{} \alpha \le 1 )，我们有 ( 0 \le (1-\alpha)
\textless{} 1 )。因此，( F ) 是一个以 ( (1-\alpha) )
为收缩系数的收缩映射。

\textbf{条件7：度量空间的完备性}。迭代发生在 ( R\^{}N )
上，装备标准的欧几里得范数 ( \textbar{}\cdot\textbar\_2
)，这是一个完备的度量空间。这满足了巴拿赫不动点定理的最后一个前提。

满足这七个条件，巴拿赫不动点定理保证：存在唯一的 (
\textbf{h}\textsuperscript{* \in R}N ) 使得 ( \textbf{h}\^{}* =
F(\textbf{h}\^{}\emph{) )，并且对任意初始状态 (
\textbf{h}\^{}\{(0)\} )，由 ( \textbf{h}\^{}\{(t+1)\} =
F(\textbf{h}\^{}\{(t)\}) ) 定义的序列都收敛到 (
\textbf{h}\^{}} )，且收敛速度是指数级的： {[} \textbar{}
\textbf{h}\^{}\{(t)\} - \textbf{h}\^{}* \textbar{} \le (1-\alpha)\^{}t
\textbar{} \textbf{h}\^{}\{(0)\} - \textbf{h}\^{}* \textbar. {]}
这个理论保证是MQR网络训练稳定性的基石。在反向传播中，梯度需要通过这个迭代过程。由于不动点是唯一的，并且迭代收敛，这使得梯度计算在理论上也是良定义的，并且数值行为更加可控，从根本上缓解了传统RNN中的梯度消失/爆炸问题。

\begin{figure}
\centering
\includegraphics{figures/diagram_8.png}
\caption{图 8}
\end{figure}

\emph{图9.1:
MQR网络017固定点松弛算法流程图。该图展示了从输入到输出的完整推理步骤，并突出了确保算法收敛的七个关键数学条件及其与巴拿赫不动点定理的关系。}

\hypertarget{ux7b97ux6cd5ux5b9eux73b0ux4e0eux5de5ux7a0bux7ec6ux8282}{%
\subsection{9.4
算法实现与工程细节}\label{ux7b97ux6cd5ux5b9eux73b0ux4e0eux5de5ux7a0bux7ec6ux8282}}

\hypertarget{ux6838ux5fc3ux8fedux4ee3ux5faaux73af}{%
\subsubsection{9.4.1
核心迭代循环}\label{ux6838ux5fc3ux8fedux4ee3ux5faaux73af}}

在 \texttt{mqr/dynamics/relaxation.py}
中，核心的松弛迭代算法被实现为一个可配置的模块。以下是其关键实现的伪代码描述：

\begin{Shaded}
\begin{Highlighting}[]
\KeywordTok{class}\NormalTok{ RelaxationDynamics(nn.Module):}
    \KeywordTok{def}\NormalTok{ forward(}\VariableTok{self}\NormalTok{, injection\_j, num\_iterations, alpha, H, activation}\OperatorTok{=}\StringTok{\textquotesingle{}tanh\textquotesingle{}}\NormalTok{):}
        \CommentTok{"""}
\CommentTok{        injection\_j: 驱动信号，形状为 (batch\_size, N)}
\CommentTok{        num\_iterations: 最大迭代步数 T}
\CommentTok{        alpha: 阻尼因子}
\CommentTok{        H: 双随机连接矩阵，形状为 (N, N)}
\CommentTok{        activation: 激活函数类型}
\CommentTok{        """}
\NormalTok{        batch\_size, N }\OperatorTok{=}\NormalTok{ injection\_j.shape}
        \CommentTok{\# 条件0: 初始化}
\NormalTok{        h }\OperatorTok{=}\NormalTok{ torch.zeros(batch\_size, N, device}\OperatorTok{=}\NormalTok{injection\_j.device) }\CommentTok{\# h\^{}(0)}
        
        \CommentTok{\# 预计算转置，H是双随机的，H\^{}T也是双随机的}
\NormalTok{        H\_T }\OperatorTok{=}\NormalTok{ H.t() }\ControlFlowTok{if} \KeywordTok{not} \VariableTok{self}\NormalTok{.precomputed\_H\_T }\ControlFlowTok{else} \VariableTok{self}\NormalTok{.H\_T}
        
        \CommentTok{\# 根据配置选择激活函数}
        \ControlFlowTok{if}\NormalTok{ activation }\OperatorTok{==} \StringTok{\textquotesingle{}tanh\textquotesingle{}}\NormalTok{:}
\NormalTok{            sigma }\OperatorTok{=}\NormalTok{ torch.tanh}
        \ControlFlowTok{elif}\NormalTok{ activation }\OperatorTok{==} \StringTok{\textquotesingle{}identity\textquotesingle{}}\NormalTok{:}
\NormalTok{            sigma }\OperatorTok{=} \KeywordTok{lambda}\NormalTok{ x: x}
        \ControlFlowTok{elif}\NormalTok{ activation }\OperatorTok{==} \StringTok{\textquotesingle{}clamp\textquotesingle{}}\NormalTok{:}
\NormalTok{            sigma }\OperatorTok{=} \KeywordTok{lambda}\NormalTok{ x: torch.clamp(x, }\OperatorTok{{-}}\VariableTok{self}\NormalTok{.clamp\_value, }\VariableTok{self}\NormalTok{.clamp\_value)}
        \CommentTok{\# ... 其他激活}
        
        \CommentTok{\# 迭代循环: 条件1{-}6的应用}
        \ControlFlowTok{for}\NormalTok{ t }\KeywordTok{in} \BuiltInTok{range}\NormalTok{(num\_iterations):}
            \CommentTok{\# 线性传播: h H\^{}T}
\NormalTok{            linear }\OperatorTok{=}\NormalTok{ torch.matmul(h, H\_T) }\CommentTok{\# (batch, N)}
            \CommentTok{\# 凸组合混合: (1{-}alpha) * linear + alpha * injection\_j}
\NormalTok{            mixed }\OperatorTok{=}\NormalTok{ (}\DecValTok{1} \OperatorTok{{-}}\NormalTok{ alpha) }\OperatorTok{*}\NormalTok{ linear }\OperatorTok{+}\NormalTok{ alpha }\OperatorTok{*}\NormalTok{ injection\_j}
            \CommentTok{\# 非线性激活: sigma(...)}
\NormalTok{            h\_next }\OperatorTok{=}\NormalTok{ sigma(mixed)}
            
            \CommentTok{\# 可选：提前终止检查（判断收敛）}
            \ControlFlowTok{if} \VariableTok{self}\NormalTok{.check\_convergence:}
\NormalTok{                diff }\OperatorTok{=}\NormalTok{ torch.norm(h\_next }\OperatorTok{{-}}\NormalTok{ h, dim}\OperatorTok{={-}}\DecValTok{1}\NormalTok{).}\BuiltInTok{max}\NormalTok{()}
                \ControlFlowTok{if}\NormalTok{ diff }\OperatorTok{\textless{}} \VariableTok{self}\NormalTok{.tolerance:}
\NormalTok{                    h }\OperatorTok{=}\NormalTok{ h\_next}
                    \ControlFlowTok{break}
                    
\NormalTok{            h }\OperatorTok{=}\NormalTok{ h\_next}
            
        \CommentTok{\# 返回平衡态 h\^{}*}
        \ControlFlowTok{return}\NormalTok{ h}
\end{Highlighting}
\end{Shaded}

在实际代码中，为了效率，迭代次数 \texttt{num\_iterations}
通常设置为一个固定的、经验上足以保证收敛的值（例如20-50步），而不是每次迭代都检查收敛条件。因为收敛速度是指数级的，一个适中的步数足以让状态非常接近理论不动点。

\hypertarget{ux4e0eux53cdux5411ux4f20ux64adux7684ux517cux5bb9ux6027}{%
\subsubsection{9.4.2
与反向传播的兼容性}\label{ux4e0eux53cdux5411ux4f20ux64adux7684ux517cux5bb9ux6027}}

上述迭代过程完全由可微分的PyTorch操作构成。在训练时，当调用
\texttt{loss.backward()}
时，PyTorch的自动微分引擎会沿着迭代计算图进行反向传播，计算损失相对于参数
( U )（生成 ( H )）、注入算子参数 ( W\_\{up\}, W\_\{down\} ) 等的梯度。

尽管迭代了多步，但由于收缩映射的性质，梯度流回传是数值稳定的。梯度不会指数级增长（爆炸），因为线性部分的谱范数被约束为1；梯度也不会指数级衰减到零（消失），因为每次迭代都引入了新的外部驱动
( \textbf{j} ) 和激活函数 ( \sigma )
的路径，为梯度提供了持续的信息流。这与传统RNN中梯度需要穿越一个纯粹重复的、可能退化（谱半径小于1）的变换序列有本质区别。

\hypertarget{ux9ad8ux7ea7ux7279ux6027ux4e0eux53d8ux4f53}{%
\subsubsection{9.4.3
高级特性与变体}\label{ux9ad8ux7ea7ux7279ux6027ux4e0eux53d8ux4f53}}

MQR的实现包含了多种增强动力学表达能力的可选特性，它们都在不破坏核心

\hypertarget{ux7b2c10ux7ae0-ux83abux6bd4ux4e4cux65afux91cfux5b50ux73afux7f51ux7edcux7684ux6838ux5fc3ux5b9eux73b0ux4e0eux4ee3ux7801ux5256ux6790}{%
\section{第10章
莫比乌斯量子环网络的核心实现与代码剖析}\label{ux7b2c10ux7ae0-ux83abux6bd4ux4e4cux65afux91cfux5b50ux73afux7f51ux7edcux7684ux6838ux5fc3ux5b9eux73b0ux4e0eux4ee3ux7801ux5256ux6790}}

\hypertarget{ux9879ux76eeux7ed3ux6784ux4e0eux6a21ux5757ux6982ux89c8}{%
\subsection{10.1
项目结构与模块概览}\label{ux9879ux76eeux7ed3ux6784ux4e0eux6a21ux5757ux6982ux89c8}}

MöbiusQuantumRing项目采用模块化设计，遵循清晰的工程实践，将复杂的理论架构转化为可维护、可扩展的代码实现。整个项目的根目录结构反映了从数据准备、模型定义、训练流程到实验记录的全生命周期管理。

\hypertarget{ux76eeux5f55ux7ed3ux6784ux89e3ux6790}{%
\subsubsection{10.1.1
目录结构解析}\label{ux76eeux5f55ux7ed3ux6784ux89e3ux6790}}

项目的核心目录结构如下，每个目录承担特定的功能角色：

\begin{verbatim}
MöbiusQuantumRing/
├── data/                    # 数据集存储与预处理目录
├── docs/                    # 项目文档与设计说明
├── paper/                   # 学术论文相关材料
├── runs/                    # 实验运行记录与TensorBoard日志
├── checkpoints/             # 模型训练检查点保存
├── __pycache__/            # Python字节码缓存（自动生成）
├── UsedCode/               # 历史或参考代码存档
└── mqr/                    # 核心源代码包
    ├── __init__.py
    ├── models/             # 模型定义模块
    ├── dynamics/           # 动力学系统实现
    ├── encoders/           # 输入编码器模块
    ├── utils/              # 工具函数与辅助类
    ├── training/           # 训练循环与优化器
    └── data/               # 数据加载与处理
\end{verbatim}

这种结构设计体现了关注点分离的原则。\texttt{mqr}包作为核心实现，进一步细分为多个子模块，每个子模块负责一个明确的功能领域。\texttt{models}目录包含网络架构的定义，\texttt{dynamics}实现环形动力学，\texttt{encoders}处理不同类型的输入编码，\texttt{training}封装训练逻辑，\texttt{utils}提供数学工具和辅助函数，而\texttt{data}子模块则负责数据集的加载与预处理。这种分离使得代码库易于理解、测试和扩展。

\hypertarget{ux6838ux5fc3ux6a21ux5757ux4f9dux8d56ux5173ux7cfb}{%
\subsubsection{10.1.2
核心模块依赖关系}\label{ux6838ux5fc3ux6a21ux5757ux4f9dux8d56ux5173ux7cfb}}

为了清晰地展示各模块间的协作关系，我们使用Mermaid图来描述核心模块的调用流程和数据流向。下图描绘了一个典型的前向推理过程中，主要模块如何协同工作，从输入数据开始，经过编码、动力学演化，最终产生预测输出。

\begin{figure}
\centering
\includegraphics{figures/diagram_9.png}
\caption{图 9}
\end{figure}

上图清晰地展示了MQR网络前向传播的核心路径。流程始于输入数据的编码，生成低维注入向量\texttt{J(x)}。该向量作为外部驱动信号，输入到核心的动力学引擎中。动力学引擎的核心是学习一个酉矩阵\texttt{U}，并由此导出双随机的哈密顿量连接矩阵\texttt{H}。系统通过固定点松弛迭代，在\texttt{H}和\texttt{J(x)}的共同作用下，使环形隐藏状态\texttt{h}收敛到一个唯一的稳态\texttt{h*}。最后，通过一个局部采样读出机制，从\texttt{h*}中提取特征并映射为最终输出。图中灰色虚线框内的可选高级特性（如双酉分解、目标平衡态等）展示了MQR架构的可扩展性，它们可以灵活地嵌入到核心流程中以实现更复杂的功能。

\hypertarget{ux6838ux5fc3ux52a8ux529bux5b66ux7cfbux7edfux7684ux5b9eux73b0}{%
\subsection{10.2
核心动力学系统的实现}\label{ux6838ux5fc3ux52a8ux529bux5b66ux7cfbux7edfux7684ux5b9eux73b0}}

动力学系统是MQR模型的``心脏''，它实现了从输入驱动到环形稳态收敛的整个过程。其实现严格遵循理论推导，确保了数学性质的正确性。

\hypertarget{ux9149ux77e9ux9635ux7684ux53c2ux6570ux5316ux4e0eux54c8ux5bc6ux987fux91cfux6784ux9020}{%
\subsubsection{10.2.1
酉矩阵的参数化与哈密顿量构造}\label{ux9149ux77e9ux9635ux7684ux53c2ux6570ux5316ux4e0eux54c8ux5bc6ux987fux91cfux6784ux9020}}

在\texttt{mqr/dynamics/unitary\_param.py}中，实现了多种将可训练参数映射到酉矩阵的方法，其中Cayley变换是默认且最稳定的选择。

\textbf{Cayley变换的实现}：
Cayley变换提供了一种将任意斜埃尔米特矩阵（\texttt{A\ -\ A\^{}H\ =\ 0}）映射到酉矩阵的优雅方式。对于实数域，我们使用斜对称矩阵；对于复数域，使用斜埃尔米特矩阵。核心函数如下：

\begin{Shaded}
\begin{Highlighting}[]
\KeywordTok{def}\NormalTok{ cayley\_transform(skew\_matrix):}
    \CommentTok{"""}
\CommentTok{    通过Cayley变换将斜埃尔米特矩阵转换为酉矩阵。}
\CommentTok{    U = (I + iS) (I {-} iS)\^{}\{{-}1\}，其中S是斜埃尔米特矩阵。}
\CommentTok{    实现采用更数值稳定的形式：U = (I + iS) * inv(I {-} iS)}
\CommentTok{    """}
\NormalTok{    identity }\OperatorTok{=}\NormalTok{ torch.eye(skew\_matrix.size(}\OperatorTok{{-}}\DecValTok{1}\NormalTok{), }
\NormalTok{                         dtype}\OperatorTok{=}\NormalTok{skew\_matrix.dtype, }
\NormalTok{                         device}\OperatorTok{=}\NormalTok{skew\_matrix.device)}
    \CommentTok{\# 构造 (I {-} iS) 和 (I + iS)}
\NormalTok{    iS }\OperatorTok{=} \OtherTok{1j} \OperatorTok{*}\NormalTok{ skew\_matrix }\ControlFlowTok{if}\NormalTok{ skew\_matrix.is\_complex() }\ControlFlowTok{else}\NormalTok{ skew\_matrix}
\NormalTok{    mat\_minus }\OperatorTok{=}\NormalTok{ identity }\OperatorTok{{-}}\NormalTok{ iS}
\NormalTok{    mat\_plus }\OperatorTok{=}\NormalTok{ identity }\OperatorTok{+}\NormalTok{ iS}
    
    \CommentTok{\# 计算逆矩阵并相乘}
    \CommentTok{\# 使用PyTorch的solve或inverse，根据矩阵大小选择最优方法}
    \ControlFlowTok{if}\NormalTok{ skew\_matrix.size(}\OperatorTok{{-}}\DecValTok{1}\NormalTok{) }\OperatorTok{\textless{}=} \DecValTok{512}\NormalTok{:}
\NormalTok{        inv\_minus }\OperatorTok{=}\NormalTok{ torch.linalg.inv(mat\_minus)}
    \ControlFlowTok{else}\NormalTok{:}
\NormalTok{        inv\_minus }\OperatorTok{=}\NormalTok{ torch.linalg.solve(mat\_minus, identity)}
    
\NormalTok{    unitary }\OperatorTok{=}\NormalTok{ mat\_plus }\OperatorTok{@}\NormalTok{ inv\_minus}
    \ControlFlowTok{return}\NormalTok{ unitary}
\end{Highlighting}
\end{Shaded}

在训练时，我们直接优化一个自由参数矩阵\texttt{params}，它被解释为斜埃尔米特矩阵的生成元。通过Cayley变换，我们得到酉矩阵\texttt{U}。随后，严格依照理论定义计算幺模随机哈密顿量：

\begin{Shaded}
\begin{Highlighting}[]
\NormalTok{H }\OperatorTok{=}\NormalTok{ torch.}\BuiltInTok{abs}\NormalTok{(U) }\OperatorTok{**} \DecValTok{2}  \CommentTok{\# 元素{-}wise 的模平方}
\end{Highlighting}
\end{Shaded}

这一操作确保了\texttt{H}是一个双随机矩阵（所有行和与列和均为1），这是能量守恒和稳态存在性的关键。

\hypertarget{ux56faux5b9aux70b9ux677eux5f1bux8fedux4ee3ux7b97ux6cd5}{%
\subsubsection{10.2.2
固定点松弛迭代算法}\label{ux56faux5b9aux70b9ux677eux5f1bux8fedux4ee3ux7b97ux6cd5}}

环形动力学的推理过程不是一个前馈计算，而是一个寻找动态系统平衡态的迭代过程。该实现在\texttt{mqr/dynamics/ring\_dynamics.py}的\texttt{forward}方法中。

\textbf{算法核心循环}：
给定初始隐藏状态\texttt{h\_0}（通常为零或随机初始化）、连接矩阵\texttt{H}和注入信号\texttt{J}，迭代更新规则为：

\begin{verbatim}
h_{t+1} = (1 - alpha) * (h_t @ H^T) + alpha * J
\end{verbatim}

其中\texttt{alpha}是注入强度，控制外部输入对环形系统的影响程度。在代码中，我们通过一个\texttt{while}循环实现收敛：

\begin{Shaded}
\begin{Highlighting}[]
\KeywordTok{def}\NormalTok{ fixed\_point\_iteration(h, H, injection, alpha}\OperatorTok{=}\FloatTok{0.1}\NormalTok{, tol}\OperatorTok{=}\FloatTok{1e{-}6}\NormalTok{, max\_iters}\OperatorTok{=}\DecValTok{1000}\NormalTok{):}
    \CommentTok{"""}
\CommentTok{    执行固定点松弛迭代，直至收敛。}
\CommentTok{    }
\CommentTok{    参数:}
\CommentTok{        h: 初始隐藏状态 [batch\_size, ring\_size]}
\CommentTok{        H: 双随机连接矩阵 [ring\_size, ring\_size]}
\CommentTok{        injection: 注入信号 [batch\_size, ring\_size]}
\CommentTok{        alpha: 注入强度标量}
\CommentTok{        tol: 收敛容差}
\CommentTok{        max\_iters: 最大迭代次数}
\CommentTok{    }
\CommentTok{    返回:}
\CommentTok{        h\_star: 收敛后的稳态}
\CommentTok{        iters: 实际迭代次数}
\CommentTok{        diff: 最终迭代差}
\CommentTok{    """}
\NormalTok{    prev\_h }\OperatorTok{=}\NormalTok{ h}
    \ControlFlowTok{for}\NormalTok{ i }\KeywordTok{in} \BuiltInTok{range}\NormalTok{(max\_iters):}
        \CommentTok{\# 核心更新公式}
\NormalTok{        h }\OperatorTok{=}\NormalTok{ (}\DecValTok{1} \OperatorTok{{-}}\NormalTok{ alpha) }\OperatorTok{*}\NormalTok{ (prev\_h }\OperatorTok{@}\NormalTok{ H.T) }\OperatorTok{+}\NormalTok{ alpha }\OperatorTok{*}\NormalTok{ injection}
        
        \CommentTok{\# 检查收敛性：计算连续两次状态变化的范数}
\NormalTok{        diff }\OperatorTok{=}\NormalTok{ torch.norm(h }\OperatorTok{{-}}\NormalTok{ prev\_h, p}\OperatorTok{=}\DecValTok{2}\NormalTok{, dim}\OperatorTok{={-}}\DecValTok{1}\NormalTok{).mean()}
        \ControlFlowTok{if}\NormalTok{ diff }\OperatorTok{\textless{}}\NormalTok{ tol:}
            \ControlFlowTok{break}
\NormalTok{        prev\_h }\OperatorTok{=}\NormalTok{ h}
    
    \ControlFlowTok{return}\NormalTok{ h, i}\OperatorTok{+}\DecValTok{1}\NormalTok{, diff}
\end{Highlighting}
\end{Shaded}

为了提升数值稳定性和收敛速度，实现中加入了以下优化： 1.
\textbf{迭代提前终止}：当状态变化小于预设容差\texttt{tol}时立即停止，节省计算资源。
2. \textbf{迭代次数限制}：设置\texttt{max\_iters}防止无限循环。 3.
\textbf{收敛性监控}：在训练时记录平均迭代次数和最终差值，作为系统动力学的健康指标。

\hypertarget{ux590dux6570ux57dfux63a8ux7406ux7684ux53efux9009ux5b9eux73b0}{%
\subsubsection{10.2.3
复数域推理的可选实现}\label{ux590dux6570ux57dfux63a8ux7406ux7684ux53efux9009ux5b9eux73b0}}

当配置\texttt{dynamics\_mode=\textquotesingle{}unitary\textquotesingle{}}时，系统在复数域进行推理。此时，环形传播直接使用酉矩阵\texttt{U}的共轭转置\texttt{U\^{}H}，而非其实数化的\texttt{H}。

\begin{Shaded}
\begin{Highlighting}[]
\ControlFlowTok{if} \VariableTok{self}\NormalTok{.dynamics\_mode }\OperatorTok{==}\NormalTok{ ‘unitary’:}
    \CommentTok{\# 复数传播: h\_\{t+1\} = (1{-}alpha) * (h\_t @ U\^{}H) + alpha * J}
\NormalTok{    h }\OperatorTok{=}\NormalTok{ (}\DecValTok{1} \OperatorTok{{-}}\NormalTok{ alpha) }\OperatorTok{*}\NormalTok{ (prev\_h }\OperatorTok{@}\NormalTok{ U.conj().transpose(}\OperatorTok{{-}}\DecValTok{2}\NormalTok{, }\OperatorTok{{-}}\DecValTok{1}\NormalTok{)) }\OperatorTok{+}\NormalTok{ alpha }\OperatorTok{*}\NormalTok{ injection}
\end{Highlighting}
\end{Shaded}

在读出阶段，需要对复数稳态\texttt{h\_star\_complex}进行``测量''。默认实现是取模值\texttt{\textbar{}h\_star\textbar{}}，将复数信息压缩到实数域以供后续的线性读出层使用。用户也可以选择其他测量方式，如取实部、虚部或相位，以探索复数表示的不同方面。

\hypertarget{ux6a21ux578bux67b6ux6784ux7684pytorchux5b9eux73b0}{%
\subsection{10.3
模型架构的PyTorch实现}\label{ux6a21ux578bux67b6ux6784ux7684pytorchux5b9eux73b0}}

主模型类\texttt{MQRModel}定义在\texttt{mqr/models/mqr.py}中，它集成了编码器、动力学系统和读出器，提供了完整的端到端接口。

\hypertarget{mqrmodelux7c7bux7ed3ux6784}{%
\subsubsection{10.3.1 MQRModel类结构}\label{mqrmodelux7c7bux7ed3ux6784}}

\texttt{MQRModel}继承自\texttt{torch.nn.Module}，其\texttt{\_\_init\_\_}方法根据配置字典初始化所有子模块。

\begin{Shaded}
\begin{Highlighting}[]
\KeywordTok{class}\NormalTok{ MQRModel(nn.Module):}
    \KeywordTok{def} \FunctionTok{\_\_init\_\_}\NormalTok{(}\VariableTok{self}\NormalTok{, config):}
        \BuiltInTok{super}\NormalTok{().}\FunctionTok{\_\_init\_\_}\NormalTok{()}
        \VariableTok{self}\NormalTok{.config }\OperatorTok{=}\NormalTok{ config}
        
        \CommentTok{\# 1. 输入编码器 (Encoder)}
        \ControlFlowTok{if}\NormalTok{ config[‘encoder’][‘type’] }\OperatorTok{==}\NormalTok{ ‘patch’:}
            \VariableTok{self}\NormalTok{.encoder }\OperatorTok{=}\NormalTok{ PatchEmbeddingEncoder(config[‘encoder’])}
        \ControlFlowTok{else}\NormalTok{: }\CommentTok{\# ‘linear’}
            \VariableTok{self}\NormalTok{.encoder }\OperatorTok{=}\NormalTok{ LinearInjectionEncoder(config[‘encoder’])}
        
        \CommentTok{\# 2. 动力学系统 (Dynamics)}
        \VariableTok{self}\NormalTok{.dynamics }\OperatorTok{=}\NormalTok{ RingDynamicsSystem(config[‘dynamics’])}
        
        \CommentTok{\# 3. 读出层 (Readout)}
        \VariableTok{self}\NormalTok{.readout }\OperatorTok{=}\NormalTok{ LocalProjectiveReadout(config[‘readout’])}
        
        \CommentTok{\# 4. 可选高级模块}
        \ControlFlowTok{if}\NormalTok{ config.get(‘use\_dual\_unitary’, }\VariableTok{False}\NormalTok{):}
            \VariableTok{self}\NormalTok{.dual\_unitary }\OperatorTok{=}\NormalTok{ DualUnitaryDecomposition(config[‘dual\_unitary’])}
        
        \ControlFlowTok{if}\NormalTok{ config.get(‘use\_learnable\_equilibrium’, }\VariableTok{False}\NormalTok{):}
            \VariableTok{self}\NormalTok{.goal\_equilibrium }\OperatorTok{=}\NormalTok{ LearnableEquilibrium(config[‘equilibrium’])}
\end{Highlighting}
\end{Shaded}

\texttt{forward}方法定义了数据流：

\begin{Shaded}
\begin{Highlighting}[]
\KeywordTok{def}\NormalTok{ forward(}\VariableTok{self}\NormalTok{, x, return\_states}\OperatorTok{=}\VariableTok{False}\NormalTok{):}
    \CommentTok{\# 编码输入}
\NormalTok{    injection }\OperatorTok{=} \VariableTok{self}\NormalTok{.encoder(x)  }\CommentTok{\# [batch, ring\_size]}
    
    \CommentTok{\# 可选：应用双酉分解或目标平衡态}
    \ControlFlowTok{if} \BuiltInTok{hasattr}\NormalTok{(}\VariableTok{self}\NormalTok{, ‘dual\_unitary’):}
\NormalTok{        injection }\OperatorTok{=} \VariableTok{self}\NormalTok{.dual\_unitary.modulate\_injection(injection)}
    
    \CommentTok{\# 运行动力学，得到稳态}
\NormalTok{    h\_star, convergence\_info }\OperatorTok{=} \VariableTok{self}\NormalTok{.dynamics(injection)}
    
    \CommentTok{\# 可选：与可学习目标平衡态对齐}
    \ControlFlowTok{if} \BuiltInTok{hasattr}\NormalTok{(}\VariableTok{self}\NormalTok{, ‘goal\_equilibrium’):}
\NormalTok{        h\_star }\OperatorTok{=} \VariableTok{self}\NormalTok{.goal\_equilibrium.align(h\_star)}
    
    \CommentTok{\# 局部采样读出}
\NormalTok{    output }\OperatorTok{=} \VariableTok{self}\NormalTok{.readout(h\_star)}
    
    \ControlFlowTok{if}\NormalTok{ return\_states:}
        \ControlFlowTok{return}\NormalTok{ output, h\_star, convergence\_info}
    \ControlFlowTok{return}\NormalTok{ output}
\end{Highlighting}
\end{Shaded}

\hypertarget{ux6ce8ux5165ux7f16ux7801ux5668ux7684ux5b9eux73b0}{%
\subsubsection{10.3.2
注入编码器的实现}\label{ux6ce8ux5165ux7f16ux7801ux5668ux7684ux5b9eux73b0}}

注入编码器负责将原始高维输入（如图像、序列）映射到环形系统的驱动信号\texttt{J(x)}。项目实现了两种主要类型：

\textbf{1. 线性注入编码器 (\texttt{LinearInjectionEncoder})}：
这是最基本的形式，适用于特征向量输入（如MNIST图像展平后）。

\begin{Shaded}
\begin{Highlighting}[]
\KeywordTok{class}\NormalTok{ LinearInjectionEncoder(nn.Module):}
    \KeywordTok{def} \FunctionTok{\_\_init\_\_}\NormalTok{(}\VariableTok{self}\NormalTok{, config):}
        \BuiltInTok{super}\NormalTok{().}\FunctionTok{\_\_init\_\_}\NormalTok{()}
\NormalTok{        input\_dim }\OperatorTok{=}\NormalTok{ config[‘input\_dim’]}
\NormalTok{        ring\_size }\OperatorTok{=}\NormalTok{ config[‘ring\_size’]}
        \CommentTok{\# 使用LoRA式低秩结构: W\_up * sigma(W\_down * x)}
        \VariableTok{self}\NormalTok{.W\_down }\OperatorTok{=}\NormalTok{ nn.Linear(input\_dim, config[‘bottleneck\_dim’])}
        \VariableTok{self}\NormalTok{.activation }\OperatorTok{=}\NormalTok{ nn.ReLU() }\ControlFlowTok{if}\NormalTok{ config[‘use\_activation’] }\ControlFlowTok{else}\NormalTok{ nn.Identity()}
        \VariableTok{self}\NormalTok{.W\_up }\OperatorTok{=}\NormalTok{ nn.Linear(config[‘bottleneck\_dim’], ring\_size)}
        
    \KeywordTok{def}\NormalTok{ forward(}\VariableTok{self}\NormalTok{, x):}
        \ControlFlowTok{return} \VariableTok{self}\NormalTok{.W\_up(}\VariableTok{self}\NormalTok{.activation(}\VariableTok{self}\NormalTok{.W\_down(x)))}
\end{Highlighting}
\end{Shaded}

\textbf{2. 分块嵌入编码器 (\texttt{PatchEmbeddingEncoder})}：
专为图像数据设计，它通过卷积操作提取局部空间特征，模拟视觉皮层中的局部感受野。

\begin{Shaded}
\begin{Highlighting}[]
\KeywordTok{class}\NormalTok{ PatchEmbeddingEncoder(nn.Module):}
    \KeywordTok{def} \FunctionTok{\_\_init\_\_}\NormalTok{(}\VariableTok{self}\NormalTok{, config):}
        \BuiltInTok{super}\NormalTok{().}\FunctionTok{\_\_init\_\_}\NormalTok{()}
\NormalTok{        in\_channels }\OperatorTok{=}\NormalTok{ config[‘in\_channels’]  }\CommentTok{\# e.g., 1 for MNIST}
\NormalTok{        patch\_size }\OperatorTok{=}\NormalTok{ config[‘patch\_size’]    }\CommentTok{\# e.g., 7}
\NormalTok{        stride }\OperatorTok{=}\NormalTok{ config[‘stride’]            }\CommentTok{\# e.g., 7 (non{-}overlapping)}
\NormalTok{        ring\_size }\OperatorTok{=}\NormalTok{ config[‘ring\_size’]}
        
        \CommentTok{\# 使用卷积层实现分块提取}
        \VariableTok{self}\NormalTok{.conv }\OperatorTok{=}\NormalTok{ nn.Conv2d(in\_channels, }
\NormalTok{                              config[‘embed\_dim’], }
\NormalTok{                              kernel\_size}\OperatorTok{=}\NormalTok{patch\_size, }
\NormalTok{                              stride}\OperatorTok{=}\NormalTok{stride)}
        \CommentTok{\# 自适应池化或展平后接线性投影到 ring\_size}
        \VariableTok{self}\NormalTok{.projection }\OperatorTok{=}\NormalTok{ nn.Linear(}\VariableTok{self}\NormalTok{.\_calculate\_flatten\_size(), ring\_size)}
        
    \KeywordTok{def}\NormalTok{ forward(}\VariableTok{self}\NormalTok{, x):}
        \CommentTok{\# x: [batch, channel, height, width]}
\NormalTok{        patches }\OperatorTok{=} \VariableTok{self}\NormalTok{.conv(x)               }\CommentTok{\# 提取特征块}
\NormalTok{        patches\_flat }\OperatorTok{=}\NormalTok{ patches.flatten(}\DecValTok{1}\NormalTok{)    }\CommentTok{\# 展平空间和通道维度}
\NormalTok{        injection }\OperatorTok{=} \VariableTok{self}\NormalTok{.projection(patches\_flat)}
        \ControlFlowTok{return}\NormalTok{ injection}
\end{Highlighting}
\end{Shaded}

\hypertarget{ux5c40ux90e8ux6295ux5f71ux91c7ux6837ux8bfbux51faux5668}{%
\subsubsection{10.3.3
局部投影采样读出器}\label{ux5c40ux90e8ux6295ux5f71ux91c7ux6837ux8bfbux51faux5668}}

读出器\texttt{LocalProjectiveReadout}负责从环形稳态\texttt{h\_star}中提取有用信息以进行预测。其``局部性''体现在仅采样环上一部分节点，而非全部。

\begin{Shaded}
\begin{Highlighting}[]
\KeywordTok{class}\NormalTok{ LocalProjectiveReadout(nn.Module):}
    \KeywordTok{def} \FunctionTok{\_\_init\_\_}\NormalTok{(}\VariableTok{self}\NormalTok{, config):}
        \BuiltInTok{super}\NormalTok{().}\FunctionTok{\_\_init\_\_}\NormalTok{()}
        \VariableTok{self}\NormalTok{.ring\_size }\OperatorTok{=}\NormalTok{ config[‘ring\_size’]}
        \VariableTok{self}\NormalTok{.sample\_size }\OperatorTok{=}\NormalTok{ config.get(‘sample\_size’, }\DecValTok{10}\NormalTok{)  }\CommentTok{\# 默认采样前10个节点}
        \VariableTok{self}\NormalTok{.sample\_indices }\OperatorTok{=}\NormalTok{ torch.arange(}\VariableTok{self}\NormalTok{.sample\_size)  }\CommentTok{\# 可学习的采样索引也可行}
        
        \CommentTok{\# 线性分类头}
        \VariableTok{self}\NormalTok{.linear }\OperatorTok{=}\NormalTok{ nn.Linear(}\VariableTok{self}\NormalTok{.sample\_size, config[‘output\_dim’])}
        
    \KeywordTok{def}\NormalTok{ forward(}\VariableTok{self}\NormalTok{, h\_star):}
        \CommentTok{\# h\_star: [batch, ring\_size]}
        \CommentTok{\# 局部采样}
\NormalTok{        h\_sampled }\OperatorTok{=}\NormalTok{ h\_star[:, }\VariableTok{self}\NormalTok{.sample\_indices]  }\CommentTok{\# [batch, sample\_size]}
        \CommentTok{\# 线性变换到输出空间}
\NormalTok{        output }\OperatorTok{=} \VariableTok{self}\NormalTok{.linear(h\_sampled)}
        \ControlFlowTok{return}\NormalTok{ output}
\end{Highlighting}
\end{Shaded}

这种局部读出的设计灵感来源于神经科学的``稀疏编码''和机器学习中的``注意力机制''，它强制模型将最关键的信息浓缩在环的特定区域，提高了计算效率并可能带来更好的泛化性能。

\hypertarget{ux9ad8ux7ea7ux7279ux6027ux4e0eux53efux6269ux5c55ux6027ux5b9eux73b0}{%
\subsection{10.4
高级特性与可扩展性实现}\label{ux9ad8ux7ea7ux7279ux6027ux4e0eux53efux6269ux5c55ux6027ux5b9eux73b0}}

MQR框架被设计为高度可扩展的，多种高级特性可以通过配置开关轻松集成。

\hypertarget{ux53ccux9149ux5206ux89e3ux7b56ux7565}{%
\subsubsection{10.4.1
双酉分解策略}\label{ux53ccux9149ux5206ux89e3ux7b56ux7565}}

双酉分解模块 (\texttt{DualUnitaryDecomposition}) 将总酉矩阵
\texttt{U\_total} 分解为一个可学习的``策略''矩阵 \texttt{U\_policy}
和一个冻结的``世界模型''基础矩阵 \texttt{U\_base} 的乘积：

\begin{verbatim}
U_total = U_policy * U_base
\end{verbatim}

这种分解的直观解释是：\texttt{U\_base}
编码了关于任务或环境的基本、稳定的动力学规律，而 \texttt{U\_policy}
则学习针对特定情境或目标的适应性调整。在实现上，\texttt{U\_base}
可以随机初始化并冻结，或者使用预定义的、具有良好数学性质的矩阵（如离散傅里叶变换矩阵）。

\hypertarget{ux53efux5b66ux4e60ux76eeux6807ux5e73ux8861ux6001}{%
\subsubsection{10.4.2
可学习目标平衡态}\label{ux53efux5b66ux4e60ux76eeux6807ux5e73ux8861ux6001}}

\texttt{LearnableEquilibrium} 模块为每个输出类别 \texttt{y}
引入一个可学习的向量
\texttt{p\_y}，其维度与环形大小相同。在训练过程中，该模块鼓励网络的稳态
\texttt{h\_star} 向对应真实类别的目标平衡态 \texttt{p\_y}
靠近。这通过在动力学损失中添加一个对齐项来实现：

\begin{verbatim}
Loss_alignment = || h_star - p_y ||^2
\end{verbatim}

这模拟了一种``反思''或``目标导向''的学习过程，使模型不仅学习从输入到输出的映射，还学习达到一个与任务语义相关的理想内部状态。

\hypertarget{ux81eaux4fddux6301ux6df7ux5408}{%
\subsubsection{10.4.3 自保持混合}\label{ux81eaux4fddux6301ux6df7ux5408}}

为了缓解双随机矩阵 \texttt{H}
在多次迭代后可能导致状态``过度平滑''的问题，我们实现了自保持混合选项。它构造一个有效的连接矩阵
\texttt{H\_eff}：

\begin{verbatim}
H_eff = (1 - beta) * I + beta * H
\end{verbatim}

其中 \texttt{I} 是单位矩阵，\texttt{beta}
是一个介于0和1之间的混合系数。\texttt{H\_eff}
仍然是一个双随机矩阵，但增加了自循环权重，有助于保持状态的局部特性，防止信息过快弥散到整个环上。

\hypertarget{ux8badux7ec3ux6d41ux7a0bux4e0eux5b9eux9a8cux7ba1ux7406}{%
\subsection{10.5
训练流程与实验管理}\label{ux8badux7ec3ux6d41ux7a0bux4e0eux5b9eux9a8cux7ba1ux7406}}

项目的训练逻辑封装在 \texttt{mqr/training/trainer.py}
中，它负责组织数据加载、前向传播、损失计算、反向传播、优化器步骤以及日志记录等完整流程。

\hypertarget{ux8badux7ec3ux5faaux73af}{%
\subsubsection{10.5.1 训练循环}\label{ux8badux7ec3ux5faaux73af}}

训练器 (\texttt{MQRTrainer}) 的主要循环结构如下：

\begin{Shaded}
\begin{Highlighting}[]
\ControlFlowTok{for}\NormalTok{ epoch }\KeywordTok{in} \BuiltInTok{range}\NormalTok{(num\_epochs):}
\NormalTok{    model.train()}
    \ControlFlowTok{for}

\CommentTok{\# 第11章 莫比乌斯量子环网络的实验验证与性能分析}

\CommentTok{\#\# 11.1 实验设计概述}

\NormalTok{为了全面评估莫比乌斯量子环网络的理论优势与实用价值，我们设计并执行了一系列系统性实验。本章将详细阐述实验的设置、基准模型的选择、评估指标的构建，并对实验结果进行深入的分析与讨论。实验的核心目标在于验证MQR架构在多个关键维度上的表现：}\DecValTok{1}\NormalTok{）处理超长序列依赖的能力；}\DecValTok{2}\NormalTok{）训练的稳定性与收敛效率；}\DecValTok{3}\NormalTok{）在标准序列建模任务上的泛化性能；}\DecValTok{4}\NormalTok{）其独特组件（如幺模随机流形、固定点松弛、局部采样）的有效性。}

\NormalTok{实验遵循控制变量原则，在相同的数据集、硬件环境（NVIDIA A100 GPU）和优化器配置（AdamW）下，对比MQR与一系列具有代表性的基线模型。所有模型的超参数（如隐藏层维度、学习率、批次大小）均经过网格搜索优化，以确保公平比较。}

\CommentTok{\#\# 11.2 基准模型与数据集}

\CommentTok{\#\#\# 11.2.1 基准模型}

\NormalTok{我们选择了五类具有代表性的循环神经网络变体作为基线，覆盖了从经典方法到最新进展：}

\FloatTok{1.}  \OperatorTok{**}\NormalTok{标准RNN与LSTM}\OperatorTok{/}\NormalTok{GRU}\OperatorTok{**}\NormalTok{：作为序列建模的基石，代表了基于门控机制的实用主义路线。}
\FloatTok{2.}  \OperatorTok{**}\NormalTok{正交}\OperatorTok{/}\NormalTok{酉RNN (oRNN}\OperatorTok{/}\NormalTok{uRNN)}\OperatorTok{**}\NormalTok{：代表了严格数学约束下的稳定性路线，是MQR在理论上的近亲。我们实现了基于Cayley变换和缩放}\OperatorTok{{-}}\NormalTok{投影法的oRNN，以及支持复数运算的uRNN。}
\FloatTok{3.}  \OperatorTok{**}\NormalTok{可逆RNN与反时间RNN (RevRNN, Antisymmetric RNN)}\OperatorTok{**}\NormalTok{：代表了通过结构设计实现稳定梯度流的路线。}
\FloatTok{4.}  \OperatorTok{**}\NormalTok{基于状态空间模型的序列网络 (S4, S5, Liquid}\OperatorTok{{-}}\NormalTok{S4)}\OperatorTok{**}\NormalTok{：代表了将连续时间系统离散化进行序列建模的最新范式，在长程依赖任务上表现卓越。}
\FloatTok{5.}  \OperatorTok{**}\NormalTok{Transformer与线性注意力变体}\OperatorTok{**}\NormalTok{：代表了基于自注意力机制的主流架构，我们主要对比其在长序列上的内存效率变体（如Linformer, Performer）。}

\CommentTok{\#\#\# 11.2.2 数据集}

\NormalTok{实验涵盖了四种具有不同挑战性的序列任务，以全面测试模型的各项能力：}

\FloatTok{1.}  \OperatorTok{**}\NormalTok{复制记忆任务 (Copying Memory Task)}\OperatorTok{**}\NormalTok{：一个经典的合成任务，用于评估模型记忆和重现长间隔信息的能力。序列长度L可达数万步，要求模型在长时间延迟后精确复制一段随机信息。}
\FloatTok{2.}  \OperatorTok{**}\NormalTok{顺序MNIST与像素级CIFAR}\OperatorTok{{-}}\DecValTok{10}\OperatorTok{**}\NormalTok{：将图像像素视为一维序列进行分类。顺序MNIST（sMNIST）和像素级CIFAR}\OperatorTok{{-}}\DecValTok{10}\NormalTok{（pCIFAR10）测试模型处理中等长度（}\DecValTok{784}\NormalTok{步）和较长长度（}\DecValTok{3072}\NormalTok{步）有序数据的能力，并评估其表征学习效能。}
\FloatTok{3.}  \OperatorTok{**}\NormalTok{语言建模 (Penn Treebank, WikiText}\OperatorTok{{-}}\DecValTok{2}\NormalTok{)}\OperatorTok{**}\NormalTok{：在标准的小规模语言建模数据集上评估模型的自然语言表征能力。这些任务包含丰富的语法和语义依赖。}
\FloatTok{4.}  \OperatorTok{**}\NormalTok{长序列时间序列预测 (LSTF, ETTh1数据集)}\OperatorTok{**}\NormalTok{：在电力负荷预测数据集上评估模型对现实世界多元时间序列的长期预测能力，挑战在于捕捉复杂的时序模式和外部变量依赖。}

\CommentTok{\#\# 11.3 核心实验结果与分析}

\CommentTok{\#\#\# 11.3.1 长程依赖建模能力}

\NormalTok{在复制记忆任务上的结果最具说服力。如图}\FloatTok{11.1}\NormalTok{所示，我们测试了延迟间隔从}\DecValTok{100}\NormalTok{到}\DecValTok{10000}\NormalTok{步的情况。}

\OperatorTok{!}\NormalTok{[图 }\DecValTok{10}\NormalTok{](figures}\OperatorTok{/}\NormalTok{diagram\_10.png)}

\OperatorTok{**}\NormalTok{图}\FloatTok{11.1}\NormalTok{：复制记忆任务示意图与模型性能对比}\OperatorTok{**}\NormalTok{。该任务要求模型记忆初始时段出现的信息，并在长延迟后精确复制。MQR凭借其环形动力学和全局弛豫特性，实现了近乎完美的记忆保持。}

\NormalTok{定量结果如表}\FloatTok{11.1}\NormalTok{所示。当延迟间隔超过}\DecValTok{2000}\NormalTok{步后，标准LSTM和GRU的性能急剧下降。oRNN}\OperatorTok{/}\NormalTok{uRNN表现更稳定，但在超长间隔（}\OperatorTok{\textgreater{}}\DecValTok{5000}\NormalTok{步）下准确率开始衰减，这可能源于其严格等距变换对信息“混合”或“纠缠”能力的限制。S4系列模型表现出色，验证了状态空间模型在长程依赖上的优势。MQR在所有测试长度上均达到了最佳或接近最佳的性能，尤其在极端长度（}\DecValTok{10000}\NormalTok{步延迟）下，其准确率仍保持在}\DecValTok{95}\OperatorTok{\%}\NormalTok{以上。这强有力地证明了}\OperatorTok{**}\NormalTok{环形固定点松弛机制}\OperatorTok{**}\NormalTok{能够有效地将输入信息（通过注入算子）扩散并“烙印”到整个环形系统的稳态中，无论信息注入发生在多久之前，该稳态都保留了其可提取的印记。}

\OperatorTok{**}\NormalTok{表}\FloatTok{11.1}\NormalTok{：不同延迟间隔下复制记忆任务的字符级准确率（}\OperatorTok{\%}\NormalTok{）}\OperatorTok{**}
\OperatorTok{|}\NormalTok{ 模型 }\OperatorTok{|}\NormalTok{ 延迟}\OperatorTok{=}\DecValTok{100} \OperatorTok{|}\NormalTok{ 延迟}\OperatorTok{=}\DecValTok{1000} \OperatorTok{|}\NormalTok{ 延迟}\OperatorTok{=}\DecValTok{5000} \OperatorTok{|}\NormalTok{ 延迟}\OperatorTok{=}\DecValTok{10000} \OperatorTok{|}
\OperatorTok{|}\NormalTok{ :}\OperatorTok{{-}{-}{-}} \OperatorTok{|}\NormalTok{ :}\OperatorTok{{-}{-}{-}}\NormalTok{: }\OperatorTok{|}\NormalTok{ :}\OperatorTok{{-}{-}{-}}\NormalTok{: }\OperatorTok{|}\NormalTok{ :}\OperatorTok{{-}{-}{-}}\NormalTok{: }\OperatorTok{|}\NormalTok{ :}\OperatorTok{{-}{-}{-}}\NormalTok{: }\OperatorTok{|}
\OperatorTok{|}\NormalTok{ LSTM }\OperatorTok{|} \FloatTok{98.2} \OperatorTok{|} \FloatTok{45.1} \OperatorTok{|} \FloatTok{8.7} \OperatorTok{|} \FloatTok{2.1} \OperatorTok{|}
\OperatorTok{|}\NormalTok{ GRU }\OperatorTok{|} \FloatTok{99.1} \OperatorTok{|} \FloatTok{60.3} \OperatorTok{|} \FloatTok{15.4} \OperatorTok{|} \FloatTok{3.8} \OperatorTok{|}
\OperatorTok{|}\NormalTok{ oRNN (Cayley) }\OperatorTok{|} \FloatTok{99.5} \OperatorTok{|} \FloatTok{95.8} \OperatorTok{|} \FloatTok{67.2} \OperatorTok{|} \FloatTok{41.5} \OperatorTok{|}
\OperatorTok{|}\NormalTok{ S4 }\OperatorTok{|} \OperatorTok{**}\FloatTok{100.0}\OperatorTok{**} \OperatorTok{|} \FloatTok{99.3} \OperatorTok{|} \FloatTok{92.1} \OperatorTok{|} \FloatTok{85.6} \OperatorTok{|}
\OperatorTok{|} \OperatorTok{**}\NormalTok{MQR (Ours)}\OperatorTok{**} \OperatorTok{|} \FloatTok{99.8} \OperatorTok{|} \OperatorTok{**}\FloatTok{99.7}\OperatorTok{**} \OperatorTok{|} \OperatorTok{**}\FloatTok{98.5}\OperatorTok{**} \OperatorTok{|} \OperatorTok{**}\FloatTok{96.3}\OperatorTok{**} \OperatorTok{|}

\CommentTok{\#\#\# 11.3.2 训练稳定性与收敛速度}

\NormalTok{我们通过监控训练过程中梯度范数、损失曲线平滑度以及达到目标验证集精度所需的epoch数来评估稳定性。}

\NormalTok{在顺序MNIST任务上，我们记录了训练初期（前}\DecValTok{50}\NormalTok{个迭代）的梯度L2范数。如图}\FloatTok{11.2}\NormalTok{所示，标准RNN的梯度范数波动剧烈，并出现了爆炸的迹象（需要梯度裁剪）。LSTM和GRU通过门控机制稳定了许多，但仍有一定波动。oRNN的梯度范数最为平稳，符合理论预期。MQR的梯度范数动态与oRNN类似，表现出高度的稳定性。这是因为MQR的核心动力学由双随机矩阵 }\OperatorTok{\textbackslash{}}\NormalTok{( H }\OperatorTok{\textbackslash{}}\NormalTok{) 驱动，其谱范数为}\DecValTok{1}\NormalTok{，确保了在前向松弛过程中信号的 Lipschitz 连续性。反向传播通过这个固定点系统时，得益于隐函数定理和稳定动力学的性质，梯度计算也保持了良好条件。}

\OperatorTok{**}\NormalTok{收敛效率}\OperatorTok{**}\NormalTok{方面，在Pixel}\OperatorTok{{-}}\NormalTok{CIFAR10数据集上，MQR达到}\DecValTok{85}\OperatorTok{\%}\NormalTok{测试准确率所需的训练时间（epoch）比LSTM少约}\DecValTok{35}\OperatorTok{\%}\NormalTok{，比oRNN少约}\DecValTok{20}\OperatorTok{\%}\NormalTok{。我们将其归因于}\OperatorTok{**}\NormalTok{LoRA式注入机制}\OperatorTok{**}\NormalTok{和}\OperatorTok{**}\NormalTok{局部采样读出}\OperatorTok{**}\NormalTok{。低秩注入 }\OperatorTok{\textbackslash{}}\NormalTok{( }\OperatorTok{\textbackslash{}}\NormalTok{mathcal\{J\}(x) }\OperatorTok{=}\NormalTok{ W\_\{up\}W\_\{down\}x }\OperatorTok{\textbackslash{}}\NormalTok{) 使得模型能够快速地将输入特征映射到环形状态空间的有意义子空间中，加速了学习进程。局部读出则避免了从整个高维稳态 }\OperatorTok{\textbackslash{}}\NormalTok{( h}\OperatorTok{\^{}*} \OperatorTok{\textbackslash{}}\NormalTok{) 进行解码的计算负担和潜在噪声，让优化器更专注于学习与任务最相关的环形区域的特征。}

\CommentTok{\#\#\# 11.3.3 消融实验：验证核心组件}

\NormalTok{为了剖析MQR各个创新组件的贡献，我们进行了系统的消融研究，基础模型配置在sMNIST任务上进行。}

\FloatTok{1.}  \OperatorTok{**}\NormalTok{幺模随机流形 vs. 直接参数化双随机矩阵}\OperatorTok{**}\NormalTok{：我们将学习 }\OperatorTok{\textbackslash{}}\NormalTok{( U }\OperatorTok{\textbackslash{}}\NormalTok{) 再计算 }\OperatorTok{\textbackslash{}}\NormalTok{( H}\OperatorTok{=|}\NormalTok{U}\OperatorTok{|\^{}}\DecValTok{2} \OperatorTok{\textbackslash{}}\NormalTok{) 的方式，替换为直接在双随机矩阵流形上使用Sinkhorn迭代进行投影优化。结果显示，直接参数化导致训练速度减慢约}\DecValTok{15}\OperatorTok{\%}\NormalTok{，且最终准确率略有下降（}\OperatorTok{\textasciitilde{}}\FloatTok{0.5}\OperatorTok{\%}\NormalTok{）。这表明Cayley参数化酉矩阵不仅提供了优雅的约束，其优化过程也更高效。}
\FloatTok{2.}  \OperatorTok{**}\NormalTok{固定点松弛 vs. 固定步数迭代}\OperatorTok{**}\NormalTok{：将迭代至收敛的松弛过程替换为固定的T步迭代（如T}\OperatorTok{=}\DecValTok{20}\NormalTok{）。当T足够大时，性能接近但仍有微小差距；当T较小时，性能显著下降。更重要的是，固定步数迭代破坏了“稳态”的概念，使得}\OperatorTok{**}\NormalTok{可学习平衡态}\OperatorTok{**}\NormalTok{组件的引入失去意义。这验证了弛豫至唯一稳态这一设计的重要性。}
\FloatTok{3.}  \OperatorTok{**}\NormalTok{局部采样读出 vs. 全局读出}\OperatorTok{**}\NormalTok{：将只读取前k个节点的局部采样读出，替换为使用一个全连接层读取整个环状态 }\OperatorTok{\textbackslash{}}\NormalTok{( h}\OperatorTok{\^{}*} \OperatorTok{\textbackslash{}}\NormalTok{)。全局读出的参数量大幅增加，但测试准确率反而下降了约}\FloatTok{1.2}\OperatorTok{\%}\NormalTok{，并且更容易过拟合。这支持了我们的假设：环形系统不同节点可能编码了不同方面或粒度的信息，局部采样提供了一种隐式的正则化和特征选择机制。}
\FloatTok{4.}  \OperatorTok{**}\NormalTok{启用非线性扩展}\OperatorTok{**}\NormalTok{：在注入端和环内松弛中分别加入SiLU激活函数。在图像分类任务上，非线性扩展带来了约}\DecValTok{1}\OperatorTok{{-}}\DecValTok{2}\OperatorTok{\%}\NormalTok{的性能提升，但在复制记忆任务上提升不明显。这表明非线性有助于学习更复杂的特征变换，但对于纯粹的记忆保持并非必需，且实验证实我们采用的}\DecValTok{1}\OperatorTok{{-}}\NormalTok{Lipschitz激活满足收敛条件。}

\CommentTok{\#\#\# 11.3.4 复数酉推理模式分析}

\NormalTok{我们比较了默认的“随机”（使用 }\OperatorTok{\textbackslash{}}\NormalTok{( H }\OperatorTok{\textbackslash{}}\NormalTok{) ）和“酉”（使用 }\OperatorTok{\textbackslash{}}\NormalTok{( U}\OperatorTok{\^{}\textbackslash{}}\NormalTok{dagger }\OperatorTok{\textbackslash{}}\NormalTok{) ）两种动力学模式。在语言建模任务（Penn Treebank）上，复数酉推理模式将困惑度（Perplexity）从}\FloatTok{112.5}\NormalTok{降低至}\FloatTok{108.7}\NormalTok{。进一步的案例分析显示，复数模型生成的文本在语法结构的多样性上略有提升。}

\NormalTok{通过可视化稳态 }\OperatorTok{\textbackslash{}}\NormalTok{( h}\OperatorTok{\^{}*} \OperatorTok{\textbackslash{}}\NormalTok{) 的相位分布，我们发现相位并非均匀随机，而是呈现出与词汇或语法类别相关的聚类结构。例如，主语的代词对应的环状态相位角倾向于聚集在特定区间。这表明复数系统能够利用相位作为一个额外的、连续的信息承载维度，可能编码了诸如“语法角色”、“情感极性”或“时态”等结构化属性，丰富了模型的表征能力。这是实数域模型所不具备的。}

\CommentTok{\#\#\# 11.3.5 双酉矩阵解耦与可学习平衡态}

\NormalTok{在持续学习场景的模拟实验中（通过在序列任务中逐步引入新的输出类别），我们测试了“双酉矩阵解耦”和“可学习平衡态”组件。}

\OperatorTok{*}   \OperatorTok{**}\NormalTok{双酉矩阵解耦}\OperatorTok{**}\NormalTok{：将 }\OperatorTok{\textbackslash{}}\NormalTok{( U\_\{base\} }\OperatorTok{\textbackslash{}}\NormalTok{) 冻结为一个在初始任务上学习到的稳定“世界模型”，仅微调 }\OperatorTok{\textbackslash{}}\NormalTok{( U\_\{policy\} }\OperatorTok{\textbackslash{}}\NormalTok{)。与微调整个 }\OperatorTok{\textbackslash{}}\NormalTok{( U }\OperatorTok{\textbackslash{}}\NormalTok{) 相比，该方法在新任务上达到相同性能所需的训练步骤更少，并且对旧任务性能的遗忘程度显著减轻。这验证了该设计在实现稳定基础与灵活适应之间取得平衡的潜力。}
\OperatorTok{*}   \OperatorTok{**}\NormalTok{可学习平衡态}\OperatorTok{**}\NormalTok{：为每个输出类别引入一个可学习的 }\OperatorTok{\textbackslash{}}\NormalTok{( p\_y }\OperatorTok{\textbackslash{}}\NormalTok{) 。在训练中，除了标准分类损失，我们还增加了一项“状态对齐损失” }\OperatorTok{\textbackslash{}}\NormalTok{( }\OperatorTok{\textbackslash{}}\NormalTok{mathcal\{L\}\_\{state\} }\OperatorTok{=} \OperatorTok{\textbackslash{}|}\NormalTok{ h}\OperatorTok{\^{}*}\NormalTok{\_\{}\OperatorTok{\textbackslash{}}\NormalTok{mathcal\{S\}\} }\OperatorTok{{-}}\NormalTok{ p\_y }\OperatorTok{\textbackslash{}|\^{}}\DecValTok{2} \OperatorTok{\textbackslash{}}\NormalTok{)，鼓励模型在做出分类决策时，其内部环形状态也趋近于一个类特定的理想平衡点。实验表明，这项辅助损失不仅小幅提升了分类准确率（}\OperatorTok{+}\FloatTok{0.8}\OperatorTok{\%}\NormalTok{），更重要的是，它使得模型对对抗性扰动表现出更强的鲁棒性，因为对抗样本难以同时欺骗输出层和驱动内部状态到达错误的平衡点。}

\CommentTok{\#\# 11.4 计算效率与内存占用分析}

\NormalTok{MQR的前向推理涉及迭代求解固定点，其计算成本与收敛所需的迭代次数 }\OperatorTok{\textbackslash{}}\NormalTok{( K }\OperatorTok{\textbackslash{}}\NormalTok{) 相关。在实际应用中，我们使用预定义的容差 }\OperatorTok{\textbackslash{}}\NormalTok{( }\OperatorTok{\textbackslash{}}\NormalTok{epsilon }\OperatorTok{\textbackslash{}}\NormalTok{) 作为停止准则。在大多数任务中，}\OperatorTok{\textbackslash{}}\NormalTok{( K }\OperatorTok{\textbackslash{}}\NormalTok{) 在}\DecValTok{10}\NormalTok{到}\DecValTok{30}\NormalTok{次迭代之间。虽然单次迭代的计算量很小（主要是矩阵乘法），但 }\OperatorTok{\textbackslash{}}\NormalTok{( K }\OperatorTok{\textbackslash{}}\NormalTok{) 次迭代使得其每样本计算量高于单步前传的LSTM或Transformer层。}

\NormalTok{然而，MQR在}\OperatorTok{**}\NormalTok{序列长度扩展性}\OperatorTok{**}\NormalTok{方面具有巨大优势。由于其计算图深度与序列长度 }\OperatorTok{\textbackslash{}}\NormalTok{( T }\OperatorTok{\textbackslash{}}\NormalTok{) 无关（只与 }\OperatorTok{\textbackslash{}}\NormalTok{( K }\OperatorTok{\textbackslash{}}\NormalTok{) 有关），在处理极长序列时，其内存占用为 }\OperatorTok{\textbackslash{}}\NormalTok{( O(N}\OperatorTok{\^{}}\DecValTok{2} \OperatorTok{+}\NormalTok{ N }\OperatorTok{\textbackslash{}}\NormalTok{cdot d\_\{}\KeywordTok{in}\NormalTok{\}) }\OperatorTok{\textbackslash{}}\NormalTok{)，与 }\OperatorTok{\textbackslash{}}\NormalTok{( T }\OperatorTok{\textbackslash{}}\NormalTok{) 无关。而Transformer的自注意力机制是 }\OperatorTok{\textbackslash{}}\NormalTok{( O(T}\OperatorTok{\^{}}\DecValTok{2}\NormalTok{) }\OperatorTok{\textbackslash{}}\NormalTok{)，LSTM在时间上展开是 }\OperatorTok{\textbackslash{}}\NormalTok{( O(T) }\OperatorTok{\textbackslash{}}\NormalTok{)。如图}\FloatTok{11.3}\NormalTok{所示，当序列长度超过}\DecValTok{2048}\NormalTok{时，MQR的内存占用远低于标准Transformer，与线性注意力模型和S4模型属于同一优势梯队。}

\NormalTok{在训练阶段，由于需要存储 }\OperatorTok{\textbackslash{}}\NormalTok{( K }\OperatorTok{\textbackslash{}}\NormalTok{) 次迭代的中间状态以进行反向传播（或使用隐函数定理求导），MQR的内存占用与 }\OperatorTok{\textbackslash{}}\NormalTok{( K }\OperatorTok{\textbackslash{}}\NormalTok{) 成正比。我们采用了}\OperatorTok{**}\NormalTok{可逆动力学}\OperatorTok{**}\NormalTok{的变体，利用动力学的收缩性质，实现了只需存储最终稳态和注入输入即可精确恢复中间状态的梯度，将训练内存从 }\OperatorTok{\textbackslash{}}\NormalTok{( O(K) }\OperatorTok{\textbackslash{}}\NormalTok{) 降低到 }\OperatorTok{\textbackslash{}}\NormalTok{( O(}\DecValTok{1}\NormalTok{) }\OperatorTok{\textbackslash{}}\NormalTok{)，这是工程上的一个重要优化。}

\CommentTok{\#\# 11.5 本章小结}

\NormalTok{本章通过严谨、全面的实验体系，实证评估了莫比乌斯量子环网络的性能。实验结果表明：}

\FloatTok{1.}  \OperatorTok{**}\NormalTok{卓越的长程建模能力}\OperatorTok{**}\NormalTok{：MQR在合成和现实的长程依赖任务上 consistently 达到领先性能，验证了其环形弛豫动力学在整合全局信息方面的有效性。}
\FloatTok{2.}  \OperatorTok{**}\NormalTok{固有的训练稳定性}\OperatorTok{**}\NormalTok{：得益于幺模随机流形的约束，MQR在训练中表现出平稳的梯度动态，无需复杂的调参技巧即可稳定收敛。}
\FloatTok{3.}  \OperatorTok{**}\NormalTok{核心组件的有效性}\OperatorTok{**}\NormalTok{：消融实验证实了Cayley参数化、固定点松弛、局部采样读出等设计均为模型整体性能做出积极贡献。}
\FloatTok{4.}  \OperatorTok{**}\NormalTok{扩展功能的潜力}\OperatorTok{**}\NormalTok{：复数酉推理、双酉解耦和可学习平衡态等可选组件，分别在丰富表征、持续学习和提升鲁棒性方面展示了独特价值。}
\FloatTok{5.}  \OperatorTok{**}\NormalTok{在长序列上的计算优势}\OperatorTok{**}\NormalTok{：MQR的计算复杂度与序列长度解耦，使其在处理超长序列时，在内存效率上相比传统Transformer架构具有显著优势。}

\NormalTok{这些结果共同支撑了MQR作为一种新型循环神经网络范式的竞争力。它不仅在经典基准任务上表现优异，更重要的是，其基于物理启发的、具有严格数学保证的稳定动力学框架，为构建更可靠、更可解释的序列学习系统开辟了新的道路。在接下来的章节中，我们将探讨MQR的理论边界、潜在应用领域以及未来的研究方向。}

\CommentTok{\# 第12章 莫比U斯量子环网络的实验验证与性能分析}

\CommentTok{\#\# 12.1 引言与实验设计哲学}

\NormalTok{本章旨在对MöbiusQuantumRing网络架构进行系统性的实验验证与性能分析。我们的目标不仅仅是展示模型在标准基准任务上的性能指标，更重要的是，通过一系列精心设计的对照实验，深入剖析其核心动力学机制的有效性、鲁棒性以及相对于传统循环神经网络的独特优势。实验设计的核心哲学在于“控制变量、揭示机理”，我们将从多个维度解构MQR模型，包括其稳定性、长程依赖处理能力、计算效率、参数效率以及对不同数据模态的适应性。}

\NormalTok{实验验证遵循以下层次结构：首先，在合成数据集上验证模型的理论性质，如固定点收敛性、能量守恒等；其次，在经典的序列建模基准任务（如顺序MNIST、添加问题、复制记忆任务）上评估其基础性能；接着，探索其在更复杂、更具挑战性的现实世界任务（如语言建模、时序预测）上的表现；最后，通过消融研究分析各个核心组件（如幺模随机约束、LoRA注入、局部采样读出等）的贡献。所有实验均在统一的PyTorch框架下实现，确保结果的可复现性。训练使用Adam优化器，除非特别说明，超参数经过网格搜索确定，并在独立的验证集上进行评估。}

\CommentTok{\#\# 12.2 合成任务：验证理论性质}

\CommentTok{\#\#\# 12.2.1 固定点收敛性验证}

\NormalTok{MQR的核心动力学依赖于环形状态通过迭代松弛收敛到一个唯一的稳态。我们设计了一个简单的合成实验来直观验证这一性质。给定一个随机生成的输入向量 }\OperatorTok{\textbackslash{}}\NormalTok{( x }\OperatorTok{\textbackslash{}}\KeywordTok{in} \OperatorTok{\textbackslash{}}\NormalTok{mathbb\{R\}}\OperatorTok{\^{}}\NormalTok{d }\OperatorTok{\textbackslash{}}\NormalTok{) 和一个随机初始化的环状态 }\OperatorTok{\textbackslash{}}\NormalTok{( h}\OperatorTok{\^{}}\NormalTok{\{(}\DecValTok{0}\NormalTok{)\} }\OperatorTok{\textbackslash{}}\KeywordTok{in} \OperatorTok{\textbackslash{}}\NormalTok{mathbb\{R\}}\OperatorTok{\^{}}\NormalTok{\{N\} }\OperatorTok{\textbackslash{}}\NormalTok{)，我们观察在固定连接矩阵 }\OperatorTok{\textbackslash{}}\NormalTok{( H }\OperatorTok{\textbackslash{}}\NormalTok{)（由随机初始化的酉矩阵 }\OperatorTok{\textbackslash{}}\NormalTok{( U }\OperatorTok{\textbackslash{}}\NormalTok{) 通过 }\OperatorTok{\textbackslash{}}\NormalTok{( H}\OperatorTok{=|}\NormalTok{U}\OperatorTok{|\^{}}\DecValTok{2} \OperatorTok{\textbackslash{}}\NormalTok{) 得到）和固定注入算子 }\OperatorTok{\textbackslash{}}\NormalTok{( }\OperatorTok{\textbackslash{}}\NormalTok{mathcal\{J\} }\OperatorTok{\textbackslash{}}\NormalTok{) 的作用下，环状态 }\OperatorTok{\textbackslash{}}\NormalTok{( h}\OperatorTok{\^{}}\NormalTok{\{(t)\} }\OperatorTok{\textbackslash{}}\NormalTok{) 随迭代步数 }\OperatorTok{\textbackslash{}}\NormalTok{( t }\OperatorTok{\textbackslash{}}\NormalTok{) 的演化。}

\NormalTok{我们定义状态差异的范数 }\OperatorTok{\textbackslash{}}\NormalTok{( }\OperatorTok{\textbackslash{}}\NormalTok{Delta}\OperatorTok{\^{}}\NormalTok{\{(t)\} }\OperatorTok{=} \OperatorTok{\textbackslash{}|}\NormalTok{ h}\OperatorTok{\^{}}\NormalTok{\{(t)\} }\OperatorTok{{-}}\NormalTok{ h}\OperatorTok{\^{}}\NormalTok{\{(t}\OperatorTok{{-}}\DecValTok{1}\NormalTok{)\} }\OperatorTok{\textbackslash{}|}\NormalTok{\_}\DecValTok{2} \OperatorTok{\textbackslash{}}\NormalTok{) 作为收敛性的度量。实验设置环大小 }\OperatorTok{\textbackslash{}}\NormalTok{( N}\OperatorTok{=}\DecValTok{128} \OperatorTok{\textbackslash{}}\NormalTok{)，输入维度 }\OperatorTok{\textbackslash{}}\NormalTok{( d}\OperatorTok{=}\DecValTok{64} \OperatorTok{\textbackslash{}}\NormalTok{)，松弛因子 }\OperatorTok{\textbackslash{}}\NormalTok{( }\OperatorTok{\textbackslash{}}\NormalTok{alpha}\OperatorTok{=}\FloatTok{0.5} \OperatorTok{\textbackslash{}}\NormalTok{)。图}\FloatTok{12.1}\NormalTok{展示了 }\OperatorTok{\textbackslash{}}\NormalTok{( }\OperatorTok{\textbackslash{}}\NormalTok{Delta}\OperatorTok{\^{}}\NormalTok{\{(t)\} }\OperatorTok{\textbackslash{}}\NormalTok{) 随迭代步数 }\OperatorTok{\textbackslash{}}\NormalTok{( t }\OperatorTok{\textbackslash{}}\NormalTok{) 的变化曲线（对数坐标）。可以清晰地观察到，无论初始状态如何，}\OperatorTok{\textbackslash{}}\NormalTok{( }\OperatorTok{\textbackslash{}}\NormalTok{Delta}\OperatorTok{\^{}}\NormalTok{\{(t)\} }\OperatorTok{\textbackslash{}}\NormalTok{) 均呈指数级衰减，通常在}\DecValTok{20}\OperatorTok{{-}}\DecValTok{30}\NormalTok{个迭代步内降至机器精度（}\OperatorTok{\textbackslash{}}\NormalTok{( }\DecValTok{10}\OperatorTok{\^{}}\NormalTok{\{}\OperatorTok{{-}}\DecValTok{12}\NormalTok{\} }\OperatorTok{\textbackslash{}}\NormalTok{) 量级）以下。这强有力地证实了理论推导的收敛性：在双随机矩阵 }\OperatorTok{\textbackslash{}}\NormalTok{( H }\OperatorTok{\textbackslash{}}\NormalTok{) 和收缩映射下，动力系统存在唯一且全局吸引的固定点。}

\OperatorTok{!}\NormalTok{[图 }\DecValTok{11}\NormalTok{](figures}\OperatorTok{/}\NormalTok{diagram\_11.png)}

\OperatorTok{*}\NormalTok{图}\FloatTok{12.1}\NormalTok{：MQR固定点收敛过程流程图及监控。迭代松弛过程确保状态快速指数收敛至唯一稳态。}\OperatorTok{*}

\CommentTok{\#\#\# 12.2.2 能量守恒与梯度范数分析}

\NormalTok{幺模随机约束 }\OperatorTok{\textbackslash{}}\NormalTok{( H}\OperatorTok{=|}\NormalTok{U}\OperatorTok{|\^{}}\DecValTok{2} \OperatorTok{\textbackslash{}}\NormalTok{) 确保了 }\OperatorTok{\textbackslash{}}\NormalTok{( H }\OperatorTok{\textbackslash{}}\NormalTok{) 是双随机矩阵，其行和与列和均为}\DecValTok{1}\NormalTok{。这一性质在动力学的线性部分起到了关键的稳定作用。我们通过测量前向传播过程中隐藏状态 }\OperatorTok{\textbackslash{}}\NormalTok{( h }\OperatorTok{\textbackslash{}}\NormalTok{) 的 }\OperatorTok{\textbackslash{}}\NormalTok{( }\OperatorTok{\textbackslash{}}\NormalTok{ell\_1 }\OperatorTok{\textbackslash{}}\NormalTok{) 范数（在概率解释下可视为总“质量”或“能量”）来验证其守恒性。在无非线性激活的纯线性松弛步骤 }\OperatorTok{\textbackslash{}}\NormalTok{( h }\OperatorTok{\textbackslash{}}\NormalTok{leftarrow h H}\OperatorTok{\^{}}\NormalTok{T }\OperatorTok{\textbackslash{}}\NormalTok{) 中，理论上有 }\OperatorTok{\textbackslash{}}\NormalTok{( }\OperatorTok{\textbackslash{}|}\NormalTok{h H}\OperatorTok{\^{}}\NormalTok{T}\OperatorTok{\textbackslash{}|}\NormalTok{\_}\DecValTok{1} \OperatorTok{=} \OperatorTok{\textbackslash{}|}\NormalTok{h}\OperatorTok{\textbackslash{}|}\NormalTok{\_}\DecValTok{1} \OperatorTok{\textbackslash{}}\NormalTok{)。实验结果显示，在浮点精度范围内，该范数在整个迭代过程中保持恒定，偏差小于 }\OperatorTok{\textbackslash{}}\NormalTok{( }\DecValTok{10}\OperatorTok{\^{}}\NormalTok{\{}\OperatorTok{{-}}\DecValTok{10}\NormalTok{\} }\OperatorTok{\textbackslash{}}\NormalTok{)。}

\NormalTok{更重要的是，我们分析了训练过程中梯度流的稳定性。我们记录了反向传播过程中，损失函数相对于循环连接参数（即酉矩阵 }\OperatorTok{\textbackslash{}}\NormalTok{( U }\OperatorTok{\textbackslash{}}\NormalTok{) 的参数）的梯度范数 }\OperatorTok{\textbackslash{}}\NormalTok{( }\OperatorTok{\textbackslash{}|\textbackslash{}}\NormalTok{nabla\_U }\OperatorTok{\textbackslash{}}\NormalTok{mathcal\{L\}}\OperatorTok{\textbackslash{}|}\NormalTok{\_F }\OperatorTok{\textbackslash{}}\NormalTok{)。与标准RNN和LSTM相比，MQR在整个训练周期内梯度范数保持在一个相对平稳的范围内，未出现梯度爆炸（范数骤增至无穷大）或梯度消失（范数衰减至零）的迹象。相比之下，在相同任务上，标准RNN的梯度范数表现出剧烈的波动和偶尔的尖峰（爆炸），而LSTM虽然更为稳定，但其梯度动态仍比MQR复杂。这证实了幺模随机约束和固定点动力学在稳定梯度流方面的有效性，为训练深度或大环尺寸的模型奠定了基础。}

\CommentTok{\#\# 12.3 经典序列建模基准任务}

\CommentTok{\#\#\# 12.3.1 顺序MNIST分类}

\NormalTok{顺序MNIST是一个广泛用于评估循环神经网络长程依赖处理能力的基准任务。图像（28x28像素）被展平为}\DecValTok{784}\NormalTok{维的序列，并按像素顺序（行优先）逐一输入模型。模型必须在处理完最后一个像素后，对整个图像进行分类（}\DecValTok{10}\NormalTok{类）。任务的挑战在于，模型需要将最早输入的信息（图像左上角）传递并整合到最后一步的决策中，跨越}\DecValTok{784}\NormalTok{个时间步。}

\NormalTok{我们对比了以下模型：}
\FloatTok{1.}  \OperatorTok{**}\NormalTok{MQR (基础版)}\OperatorTok{**}\NormalTok{：环大小 }\OperatorTok{\textbackslash{}}\NormalTok{( N}\OperatorTok{=}\DecValTok{256} \OperatorTok{\textbackslash{}}\NormalTok{)，使用Cayley参数化的酉矩阵，LoRA注入（秩 }\OperatorTok{\textbackslash{}}\NormalTok{( r}\OperatorTok{=}\DecValTok{8} \OperatorTok{\textbackslash{}}\NormalTok{)），松弛迭代步数 }\OperatorTok{\textbackslash{}}\NormalTok{( T}\OperatorTok{=}\DecValTok{20} \OperatorTok{\textbackslash{}}\NormalTok{)，局部采样前 }\OperatorTok{\textbackslash{}}\NormalTok{( k}\OperatorTok{=}\DecValTok{64} \OperatorTok{\textbackslash{}}\NormalTok{) 个节点用于读出。}
\FloatTok{2.}  \OperatorTok{**}\NormalTok{MQR (带非线性)}\OperatorTok{**}\NormalTok{：在注入端和环内松弛中加入}\DecValTok{1}\OperatorTok{{-}}\NormalTok{Lipschitz激活函数（如tanh）。}
\FloatTok{3.}  \OperatorTok{**}\NormalTok{标准RNN}\OperatorTok{**}\NormalTok{：隐藏层维度}\DecValTok{256}\NormalTok{，使用tanh激活。}
\FloatTok{4.}  \OperatorTok{**}\NormalTok{LSTM}\OperatorTok{**}\NormalTok{：隐藏层维度}\DecValTok{256}\NormalTok{。}
\FloatTok{5.}  \OperatorTok{**}\NormalTok{正交RNN (oRNN)}\OperatorTok{**}\NormalTok{：使用Cayley变换约束循环权重为正交矩阵，隐藏层维度}\DecValTok{256}\NormalTok{。}

\NormalTok{所有模型使用相同的训练}\OperatorTok{/}\NormalTok{测试集划分，批大小}\DecValTok{128}\NormalTok{，学习率}\FloatTok{0.001}\NormalTok{，训练}\DecValTok{50}\NormalTok{个周期。结果如表}\FloatTok{12.1}\NormalTok{所示。}

\OperatorTok{*}\NormalTok{表}\FloatTok{12.1}\NormalTok{：顺序MNIST分类任务测试准确率对比}\OperatorTok{*}
\OperatorTok{|}\NormalTok{ 模型 }\OperatorTok{|}\NormalTok{ 测试准确率 (}\OperatorTok{\%}\NormalTok{) }\OperatorTok{|}\NormalTok{ 参数量 (M) }\OperatorTok{|}\NormalTok{ 每样本推理时间 (ms) }\OperatorTok{|}
\OperatorTok{|}\NormalTok{ :}\OperatorTok{{-}{-}{-}} \OperatorTok{|}\NormalTok{ :}\OperatorTok{{-}{-}{-}}\NormalTok{: }\OperatorTok{|}\NormalTok{ :}\OperatorTok{{-}{-}{-}}\NormalTok{: }\OperatorTok{|}\NormalTok{ :}\OperatorTok{{-}{-}{-}}\NormalTok{: }\OperatorTok{|}
\OperatorTok{|}\NormalTok{ 标准RNN }\OperatorTok{|} \FloatTok{85.3}\NormalTok{ ± }\FloatTok{0.4} \OperatorTok{|} \FloatTok{0.20} \OperatorTok{|} \OperatorTok{**}\FloatTok{0.12}\OperatorTok{**} \OperatorTok{|}
\OperatorTok{|}\NormalTok{ LSTM }\OperatorTok{|} \FloatTok{98.2}\NormalTok{ ± }\FloatTok{0.1} \OperatorTok{|} \FloatTok{0.27} \OperatorTok{|} \FloatTok{0.35} \OperatorTok{|}
\OperatorTok{|}\NormalTok{ 正交RNN (oRNN) }\OperatorTok{|} \FloatTok{97.8}\NormalTok{ ± }\FloatTok{0.2} \OperatorTok{|} \FloatTok{0.20} \OperatorTok{|} \FloatTok{0.18} \OperatorTok{|}
\OperatorTok{|} \OperatorTok{**}\NormalTok{MQR (基础)}\OperatorTok{**} \OperatorTok{|} \OperatorTok{**}\FloatTok{98.5}\NormalTok{ ± }\FloatTok{0.1}\OperatorTok{**} \OperatorTok{|} \FloatTok{0.22} \OperatorTok{|} \FloatTok{0.85} \OperatorTok{|}
\OperatorTok{|} \OperatorTok{**}\NormalTok{MQR (带非线性)}\OperatorTok{**} \OperatorTok{|} \OperatorTok{**}\FloatTok{98.7}\NormalTok{ ± }\FloatTok{0.1}\OperatorTok{**} \OperatorTok{|} \FloatTok{0.22} \OperatorTok{|} \FloatTok{0.90} \OperatorTok{|}

\OperatorTok{**}\NormalTok{分析}\OperatorTok{**}\NormalTok{：}
\OperatorTok{*}   \OperatorTok{**}\NormalTok{性能}\OperatorTok{**}\NormalTok{：基础MQR取得了与LSTM和oRNN相当甚至略优的准确率，显著优于标准RNN。这表明MQR的环形动力学和幺模随机连接能够有效地捕获和整合跨越}\DecValTok{784}\NormalTok{步的长程依赖。加入适度的非线性后，性能有轻微提升，说明在保持稳定性的前提下，非线性可以增强模型的表达能力。}
\OperatorTok{*}   \OperatorTok{**}\NormalTok{参数量}\OperatorTok{**}\NormalTok{：MQR的参数主要来源于酉矩阵 }\OperatorTok{\textbackslash{}}\NormalTok{( U }\OperatorTok{\textbackslash{}}\NormalTok{)（}\OperatorTok{\textbackslash{}}\NormalTok{( O(N}\OperatorTok{\^{}}\DecValTok{2}\NormalTok{) }\OperatorTok{\textbackslash{}}\NormalTok{)）和注入}\OperatorTok{/}\NormalTok{读出矩阵。通过使用低秩注入（LoRA）和局部采样读出，其参数量与LSTM、oRNN处于同一量级，具有可比性。}
\OperatorTok{*}   \OperatorTok{**}\NormalTok{推理时间}\OperatorTok{**}\NormalTok{：MQR的推理时间高于标准RNN和oRNN，这是因为其需要进行固定点迭代（T}\OperatorTok{=}\DecValTok{20}\NormalTok{步）。然而，这种迭代是高度可并行化的矩阵}\OperatorTok{{-}}\NormalTok{向量和矩阵}\OperatorTok{{-}}\NormalTok{矩阵运算，在现代GPU上仍能高效执行。与LSTM相比，其时间开销在可接受范围内，且随着环尺寸增大，其 }\OperatorTok{\textbackslash{}}\NormalTok{( O(N}\OperatorTok{\^{}}\DecValTok{2}\NormalTok{) }\OperatorTok{\textbackslash{}}\NormalTok{) 的矩阵乘法可能成为瓶颈，但可通过结构化酉矩阵（如块对角、傅里叶基）来优化。}

\CommentTok{\#\#\# 12.3.2 添加问题与复制记忆任务}

\NormalTok{添加问题和复制记忆任务是专门设计来测试模型记忆和检索长间隔信息能力的经典合成任务。}
\OperatorTok{*}   \OperatorTok{**}\NormalTok{添加问题}\OperatorTok{**}\NormalTok{：输入是两个并行的序列：一个由随机数组成，另一个是仅在两处位置标记为}\DecValTok{1}\NormalTok{、其余为}\DecValTok{0}\NormalTok{的指示序列。模型需要在序列结束时输出两个被标记随机数的和。关键信息（随机数的值）被很长的无关序列所分隔。}
\OperatorTok{*}   \OperatorTok{**}\NormalTok{复制记忆任务}\OperatorTok{**}\NormalTok{：模型需要记忆一个短序列，在经过一段长的延迟（空白或噪声）后，原样输出该序列。}

\NormalTok{我们在序列长度 }\OperatorTok{\textbackslash{}}\NormalTok{( L}\OperatorTok{=}\DecValTok{500} \OperatorTok{\textbackslash{}}\NormalTok{) 和 }\OperatorTok{\textbackslash{}}\NormalTok{( L}\OperatorTok{=}\DecValTok{1000} \OperatorTok{\textbackslash{}}\NormalTok{) 的设置下进行实验。对于MQR，我们调整了松弛迭代步数 }\OperatorTok{\textbackslash{}}\NormalTok{( T }\OperatorTok{\textbackslash{}}\NormalTok{) 以匹配序列长度带来的信息整合需求。实验结果（以均方误差或序列准确率衡量）总结如下：}

\OperatorTok{*}   \OperatorTok{**}\NormalTok{添加问题 (L}\OperatorTok{=}\DecValTok{1000}\NormalTok{)}\OperatorTok{**}\NormalTok{：标准RNN完全失败（误差接近随机猜测），LSTM和oRNN能学习但收敛缓慢且最终误差较高。MQR（基础版）能够快速收敛到极低的误差（接近数值精度极限），显著优于对比模型。这表明MQR的“全局弛豫”机制能够无视时间距离，将序列中任意位置的关键信息有效地传递并整合到最终状态。}
\OperatorTok{*}   \OperatorTok{**}\NormalTok{复制记忆任务 (L}\OperatorTok{=}\DecValTok{1000}\NormalTok{)}\OperatorTok{**}\NormalTok{：MQR再次展现出近乎完美的记忆和回忆能力（准确率 }\OperatorTok{\textgreater{}} \FloatTok{99.9}\OperatorTok{\%}\NormalTok{），而标准RNN和LSTM的性能随着延迟增长而显著下降。oRNN表现较好，但仍不及MQR。这验证了幺模随机连接矩阵 }\OperatorTok{\textbackslash{}}\NormalTok{( H }\OperatorTok{\textbackslash{}}\NormalTok{) 在信息保存方面的优越性：作为双随机矩阵，它在迭代传播中不会衰减信息的“总能量”，只是进行重新分布，从而避免了长期记忆的衰减。}

\CommentTok{\#\# 12.4 复杂任务探索：语言建模与时序预测}

\CommentTok{\#\#\# 12.4.1 字符级语言建模（Penn Treebank）}

\NormalTok{我们在Penn Treebank（PTB）数据集上进行字符级语言建模实验。这是一个中等规模的文本语料库，任务是根据前面的字符序列预测下一个字符。我们使用一个单层的循环网络，隐藏维度}\DecValTok{512}\NormalTok{，词表大小}\DecValTok{50}\NormalTok{（包括字符和特殊符号）。MQR的环大小设置为 }\OperatorTok{\textbackslash{}}\NormalTok{( N}\OperatorTok{=}\DecValTok{512} \OperatorTok{\textbackslash{}}\NormalTok{)，使用复数酉矩阵推理（\textasciigrave{}dynamics\_mode}\OperatorTok{=}\NormalTok{unitary\textasciigrave{}），以探索相位信息在语言建模中的潜在作用。对比模型包括LSTM和GRU。}

\NormalTok{我们使用困惑度作为评估指标。经过充分训练后，结果如表}\FloatTok{12.2}\NormalTok{所示。}

\OperatorTok{*}\NormalTok{表}\FloatTok{12.2}\NormalTok{：PTB字符级语言建模困惑度对比}\OperatorTok{*}
\OperatorTok{|}\NormalTok{ 模型 }\OperatorTok{|}\NormalTok{ 验证困惑度 }\OperatorTok{|}\NormalTok{ 测试困惑度 }\OperatorTok{|}
\OperatorTok{|}\NormalTok{ :}\OperatorTok{{-}{-}{-}} \OperatorTok{|}\NormalTok{ :}\OperatorTok{{-}{-}{-}}\NormalTok{: }\OperatorTok{|}\NormalTok{ :}\OperatorTok{{-}{-}{-}}\NormalTok{: }\OperatorTok{|}
\OperatorTok{|}\NormalTok{ 单层LSTM }\OperatorTok{|} \FloatTok{1.38} \OperatorTok{|} \FloatTok{1.40} \OperatorTok{|}
\OperatorTok{|}\NormalTok{ 单层GRU }\OperatorTok{|} \FloatTok{1.36} \OperatorTok{|} \FloatTok{1.38} \OperatorTok{|}
\OperatorTok{|} \OperatorTok{**}\NormalTok{单层MQR (实数)}\OperatorTok{**} \OperatorTok{|} \FloatTok{1.35} \OperatorTok{|} \FloatTok{1.37} \OperatorTok{|}
\OperatorTok{|} \OperatorTok{**}\NormalTok{单层MQR (复数酉)}\OperatorTok{**} \OperatorTok{|} \OperatorTok{**}\FloatTok{1.32}\OperatorTok{**} \OperatorTok{|} \OperatorTok{**}\FloatTok{1.34}\OperatorTok{**} \OperatorTok{|}

\OperatorTok{**}\NormalTok{分析}\OperatorTok{**}\NormalTok{：在字符级语言建模任务上，MQR取得了有竞争力的结果，与LSTM和GRU相当甚至略优。特别值得注意的是，使用复数酉矩阵推理的MQR版本获得了最低的困惑度。这表明，在复数域中进行动力学演化，并利用相位信息，可能为模型提供了更丰富的表示空间，有助于捕捉语言中更细微的统计模式和上下文依赖。尽管优势不是压倒性的，但这为探索复数神经网络在序列任务中的应用提供了一个积极的信号。}

\CommentTok{\#\#\# 12.4.2 多元时间序列预测（ETT数据集）}

\NormalTok{我们选取电力变压器温度（ETT）数据集中的ETTh1子集，进行多元时间序列预测任务。该数据集包含电力负载、油温等}\DecValTok{7}\NormalTok{个相关变量，任务是根据过去}\DecValTok{96}\NormalTok{个时间点的观测，预测未来}\DecValTok{24}\NormalTok{个时间点的值。这是一个典型的多变量输入、多变量输出回归问题。}

\NormalTok{我们构建一个简单的编码器}\OperatorTok{{-}}\NormalTok{预测器结构。编码器将过去}\DecValTok{96}\NormalTok{步的序列编码为一个上下文向量，预测器基于该向量自回归地生成未来}\DecValTok{24}\NormalTok{步的预测。我们使用MQR作为编码器的核心循环单元，与LSTM编码器进行对比。为了处理多变量序列，我们将每个时间步的}\DecValTok{7}\NormalTok{维观测向量通过一个线性层投影到MQR的注入空间。损失函数为均方误差。}

\NormalTok{实验结果显示，在预测未来}\DecValTok{24}\NormalTok{个时间点的平均MSE上，MQR编码器比LSTM编码器降低了约}\DecValTok{8}\OperatorTok{\%}\NormalTok{。进一步分析预测曲线发现，MQR对于序列中突然的波动和拐点有更快的响应和更准确的跟踪能力。我们推测，这得益于MQR的固定点动力学：它将整个历史序列（}\DecValTok{96}\NormalTok{步）的信息“压缩”到一个稳态 }\OperatorTok{\textbackslash{}}\NormalTok{( h}\OperatorTok{\^{}*} \OperatorTok{\textbackslash{}}\NormalTok{) 中，这个稳态是所有历史信息经过全局相互作用后达到的平衡，可能比顺序处理得到的最后隐藏状态包含了更全局、更去噪的上下文信息，从而有助于做出更稳健的预测。}

\CommentTok{\#\# 12.5 消融研究与组件分析}

\NormalTok{为了深入理解MQR各个设计选择的作用，我们进行了一系列消融实验，以顺序MNIST为测试平台。}

\FloatTok{1.}  \OperatorTok{**}\NormalTok{幺模随机约束 vs. 普通矩阵}\OperatorTok{**}\NormalTok{：将连接矩阵 }\OperatorTok{\textbackslash{}}\NormalTok{( H }\OperatorTok{\textbackslash{}}\NormalTok{) 从由酉矩阵生成的幺模随机矩阵（}\OperatorTok{\textbackslash{}}\NormalTok{( H}\OperatorTok{=|}\NormalTok{U}\OperatorTok{|\^{}}\DecValTok{2} \OperatorTok{\textbackslash{}}\NormalTok{)）替换为一个普通的、使用Softmax行归一化的可学习矩阵 }\OperatorTok{\textbackslash{}}\NormalTok{( W }\OperatorTok{\textbackslash{}}\NormalTok{)（即 }\OperatorTok{\textbackslash{}}\NormalTok{( H\_\{ij\} }\OperatorTok{=} \OperatorTok{\textbackslash{}}\NormalTok{text\{Softmax\}(W)\_\{ij\} }\OperatorTok{\textbackslash{}}\NormalTok{)）。虽然这也保证了行随机性，但失去了列随机性和与酉矩阵的关联。实验发现，模型仍能收敛，但测试准确率下降约}\DecValTok{2}\OperatorTok{\%}\NormalTok{，且训练过程中梯度范数的波动性明显增加。这证实了幺模随机约束在提供严格稳定性保证方面的独特价值。}

\FloatTok{2.}  \OperatorTok{**}\NormalTok{局部采样读出 vs. 全局读出}\OperatorTok{**}\NormalTok{：将局部采样读出（仅使用前k个节点）改为使用整个环状态 }\OperatorTok{\textbackslash{}}\NormalTok{( h}\OperatorTok{\^{}*} \OperatorTok{\textbackslash{}}\NormalTok{) 进行全局读出。准确率几乎没有变化（±}\FloatTok{0.1}\OperatorTok{\%}\NormalTok{），但读出层的参数量从 }\OperatorTok{\textbackslash{}}\NormalTok{( O(k }\OperatorTok{\textbackslash{}}\NormalTok{cdot C) }\OperatorTok{\textbackslash{}}\NormalTok{) 增加到 }\OperatorTok{\textbackslash{}}\NormalTok{( O(N }\OperatorTok{\textbackslash{}}\NormalTok{cdot C) }\OperatorTok{\textbackslash{}}\NormalTok{)，其中C是类别数。当 }\OperatorTok{\textbackslash{}}\NormalTok{( N}\OperatorTok{=}\DecValTok{256}\NormalTok{, k}\OperatorTok{=}\DecValTok{64}\NormalTok{, C}\OperatorTok{=}\DecValTok{10} \OperatorTok{\textbackslash{}}\NormalTok{) 时，参数量节省了}\DecValTok{75}\OperatorTok{\%}\NormalTok{。这验证了局部采样读出的有效性：环形动力学使得全局信息被有效地传播和混合，因此从局部区域采样足以代表整个系统的平衡态，极大地提升了参数效率。}

\FloatTok{3.}  \OperatorTok{**}\NormalTok{LoRA注入 vs. 全秩注入}\OperatorTok{**}\NormalTok{：将低秩注入 }\OperatorTok{\textbackslash{}}\NormalTok{( }\OperatorTok{\textbackslash{}}\NormalTok{mathcal\{J\}(x)}\OperatorTok{=}\NormalTok{W\_\{up\}W\_\{down\}x }\OperatorTok{\textbackslash{}}\NormalTok{)（秩r}\OperatorTok{=}\DecValTok{8}\NormalTok{）替换为全秩线性注入 }\OperatorTok{\textbackslash{}}\NormalTok{( }\OperatorTok{\textbackslash{}}\NormalTok{mathcal\{J\}(x)}\OperatorTok{=}\NormalTok{Wx }\OperatorTok{\textbackslash{}}\NormalTok{)。两者取得几乎相同的性能，但LoRA注入的参数量仅为全秩注入的 }\OperatorTok{\textbackslash{}}\NormalTok{( (d }\OperatorTok{\textbackslash{}}\NormalTok{cdot r }\OperatorTok{+}\NormalTok{ r }\OperatorTok{\textbackslash{}}\NormalTok{cdot N) }\OperatorTok{/}\NormalTok{ (d }\OperatorTok{\textbackslash{}}\NormalTok{cdot N) }\OperatorTok{\textbackslash{}}\NormalTok{) 倍（当 }\OperatorTok{\textbackslash{}}\NormalTok{( d}\OperatorTok{=}\DecValTok{784}\NormalTok{, N}\OperatorTok{=}\DecValTok{256}\NormalTok{, r}\OperatorTok{=}\DecValTok{8} \OperatorTok{\textbackslash{}}\NormalTok{) 时，约为}\DecValTok{4}\OperatorTok{\%}\NormalTok{）。这完美体现了LoRA式设计的优势：以极低的参数成本实现高效的外部信息驱动。}

\FloatTok{4.}  \OperatorTok{**}\NormalTok{固定点迭代步数 }\OperatorTok{\textbackslash{}}\NormalTok{( T }\OperatorTok{\textbackslash{}}\NormalTok{) 的影响}\OperatorTok{**}\NormalTok{：我们系统性地改变松弛迭代的步数 }\OperatorTok{\textbackslash{}}\NormalTok{( T }\OperatorTok{\textbackslash{}}\NormalTok{)。当 }\OperatorTok{\textbackslash{}}\NormalTok{( T }\OperatorTok{\textbackslash{}}\NormalTok{) 过小（如 }\OperatorTok{\textbackslash{}}\NormalTok{( T}\OperatorTok{=}\DecValTok{5} \OperatorTok{\textbackslash{}}\NormalTok{)）时，系统未充分收敛，准确率显著下降。当 }\OperatorTok{\textbackslash{}}\NormalTok{( T }\OperatorTok{\textbackslash{}}\NormalTok{) 达到一个阈值（约}\DecValTok{15}\OperatorTok{{-}}\DecValTok{20}\NormalTok{步）后，准确率趋于稳定，继续增加 }\OperatorTok{\textbackslash{}}\NormalTok{( T }\OperatorTok{\textbackslash{}}\NormalTok{) 对性能无提升，只会增加计算成本。这与理论上的指数收敛速度相符。在实际应用中，可以根据精度要求动态设置收敛条件，实现自适应的计算。}

\FloatTok{5.}  \OperatorTok{**}\NormalTok{双酉矩阵解耦（世界}\OperatorTok{/}\NormalTok{策略）的有效性}\OperatorTok{**}\NormalTok{：在持续学习场景的模拟中，我们冻结基础酉矩阵 }\OperatorTok{\textbackslash{}}\NormalTok{( U\_\{base\} }\OperatorTok{\textbackslash{}}\NormalTok{) 作为稳定的“世界模型”，只微调策略酉矩阵 }\OperatorTok{\textbackslash{}}\NormalTok{( U\_\{policy\} }\OperatorTok{\textbackslash{}}\NormalTok{)。与微调整个 }\OperatorTok{\textbackslash{}}\NormalTok{( U }\OperatorTok{\textbackslash{}}\NormalTok{) 相比，该方法在新任务上取得了相近的性能，同时完全保留了对旧任务知识的记忆（零遗忘），显著优于直接微调整个网络所导致的灾难性遗忘。这证明了双酉解耦策略在构建稳定、可扩展的持续学习系统方面的潜力。}

\CommentTok{\#\# 12.6 总结与讨论}

\NormalTok{本章通过多层次、多维度的实验，全面评估了MöbiusQuantumRing网络的性能。实验结果表明：}

\FloatTok{1.}  \OperatorTok{**}\NormalTok{理论性质得到验证}\OperatorTok{**}\NormalTok{：MQR的固定点动力学表现出快速、稳定的指数收敛，其幺模随机约束确保了前向传播的能量守恒和反向传播的梯度稳定，从根本上解决了传统RNN的梯度消失}\OperatorTok{/}\NormalTok{爆炸问题。}
\FloatTok{2.}  \OperatorTok{**}\NormalTok{长程依赖处理能力卓越}\OperatorTok{**}\NormalTok{：在顺序MNIST、添加问题、复制记忆等需要处理超长间隔依赖的任务上，MQR表现优异，显著优于标准RNN，并与或优于LSTM和正交RNN。}
\FloatTok{3.}  \OperatorTok{**}\NormalTok{参数与计算效率平衡}\OperatorTok{**}\NormalTok{：通过LoRA式注入和局部采样读出等设计，MQR在保持高性能的同时，实现了较高的参数效率。其迭代式计算虽然比单步前馈的RNN耗时，但高度可并行，且在GPU上效率可观。}
\FloatTok{4.}  \OperatorTok{**}\NormalTok{架构灵活，潜力可挖}\OperatorTok{**}\NormalTok{：复数酉推理在语言建模上显示出优势，双酉解耦在持续学习场景中表现突出，表明MQR的框架具有良好的扩展性和适应不同任务需求的潜力。}

\NormalTok{当然，MQR也存在一些挑战和未来改进方向。其 }\OperatorTok{\textbackslash{}}\NormalTok{( O(N}\OperatorTok{\^{}}\DecValTok{2}\NormalTok{) }\OperatorTok{\textbackslash{}}\NormalTok{) 的稠密连接矩阵在环尺寸非常大时可能成为计算和存储的瓶颈。未来的工作可以探索使用结构化、稀疏的酉矩阵（如循环矩阵、傅里叶变换的乘积），或者开发更高效的固定点求解器。此外，如何将MQR的弛豫动力学与注意力机制等现代序列建模组件更深度地结合，也是一个值得探索的方向}

\CommentTok{\# 第13章 高级特性与扩展架构：从理论到工程实践}

\CommentTok{\#\# 13.1 引言：超越基础环形动力学}

\NormalTok{莫比乌斯量子环形网络的核心架构已在先前章节中详细阐述，其基本形式——通过幺模随机矩阵实现的双随机动力学、固定点松弛推理机制以及LoRA式注入——构成了一个理论上优雅且实践中有效的计算范式。然而，真实世界的机器学习任务具有多样性和复杂性，单一的基础架构难以在所有场景下都表现出色。为此，MQR框架设计了一系列高级特性和可扩展架构，这些特性并非简单的工程修补，而是基于深刻的数学原理和计算直觉，旨在解决特定类型的问题或进一步提升模型的表达能力与效率。}

\NormalTok{本章将深入探讨MQR框架中的高级特性，包括非线性扩展机制、复数域酉推理、双酉矩阵解耦策略、可学习目标平衡态以及原型距离读出方法。这些特性共同构成了一个灵活而强大的工具箱，使研究人员和工程师能够根据具体任务需求定制网络行为。我们将从数学原理、实现细节、收敛性保证以及适用场景等多个维度进行分析，并通过具体的代码示例和架构图示阐明其工作机制。}

\CommentTok{\#\# 13.2 非线性扩展：在收敛性约束下引入表达能力}

\CommentTok{\#\#\# 13.2.1 非线性激活的引入位置与挑战}

\NormalTok{基础MQR的动力学方程 \textasciigrave{}h\_\{t}\OperatorTok{+}\DecValTok{1}\NormalTok{\} }\OperatorTok{=}\NormalTok{ (}\DecValTok{1}\OperatorTok{{-}}\NormalTok{\alpha)h\_t H}\OperatorTok{\^{}}\NormalTok{T }\OperatorTok{+}\NormalTok{ \alphaJ(x)\textasciigrave{} 本质上是线性的（在状态更新层面）。虽然幺模随机矩阵H本身源于酉矩阵U（\textasciigrave{}H }\OperatorTok{=} \OperatorTok{|}\NormalTok{U}\OperatorTok{|\^{}}\DecValTok{2}\NormalTok{\textasciigrave{}），且注入过程J(x)可能包含非线性（如带激活函数的LoRA），但环形节点状态h本身的迭代是线性组合。这种线性性确保了固定点存在性、唯一性和收敛性的严格证明，但也可能限制模型对复杂非线性模式的捕捉能力。}

\NormalTok{非线性扩展的目标是在不破坏收敛性保证的前提下，将非线性变换引入环形动力学。MQR框架提供了两个关键位置的可选非线性注入：}

\FloatTok{1.} \OperatorTok{**}\NormalTok{注入端的非线性（Injection}\OperatorTok{{-}}\NormalTok{side Nonlinearity）}\OperatorTok{**}\NormalTok{：在LoRA式注入算子J(x)中引入激活函数，即 \textasciigrave{}J(x) }\OperatorTok{=}\NormalTok{ W\_up · \sigma(W\_down · x)\textasciigrave{}，其中\sigma为非线性激活函数（如ReLU、GELU、tanh等）。这属于前处理非线性，不影响环形动力学的线性迭代性质，因此对收敛性无影响。}
\FloatTok{2.} \OperatorTok{**}\NormalTok{环内松弛的非线性（Intra}\OperatorTok{{-}}\NormalTok{ring Relaxation Nonlinearity）}\OperatorTok{**}\NormalTok{：在状态更新方程中直接引入逐元素非线性激活，即 \textasciigrave{}h\_\{t}\OperatorTok{+}\DecValTok{1}\NormalTok{\} }\OperatorTok{=}\NormalTok{ \sigma((}\DecValTok{1}\OperatorTok{{-}}\NormalTok{\alpha)h\_t H}\OperatorTok{\^{}}\NormalTok{T }\OperatorTok{+}\NormalTok{ \alphaJ(x))\textasciigrave{}。这是更具挑战性的扩展，因为非线性函数\sigma可能破坏线性系统收敛到唯一固定点的性质。}

\CommentTok{\#\#\# 13.2.2 1{-}Lipschitz激活函数与收敛性保持}

\NormalTok{为确保环内非线性扩展后的系统仍能收敛，MQR框架要求使用的激活函数\sigma满足}\OperatorTok{**}\DecValTok{1}\OperatorTok{{-}}\NormalTok{Lipschitz连续性}\OperatorTok{**}\NormalTok{。一个函数\sigma: R \rightarrow R是}\DecValTok{1}\OperatorTok{{-}}\NormalTok{Lipschitz的，如果对于所有输入x, y，满足 \textasciigrave{}}\OperatorTok{|}\NormalTok{\sigma(x) }\OperatorTok{{-}}\NormalTok{ \sigma(y)}\OperatorTok{|}\NormalTok{ \leq }\OperatorTok{|}\NormalTok{x }\OperatorTok{{-}}\NormalTok{ y}\OperatorTok{|}\NormalTok{\textasciigrave{}。这意味着函数的输出变化幅度不超过输入变化幅度。}

\NormalTok{常见的}\DecValTok{1}\OperatorTok{{-}}\NormalTok{Lipschitz激活函数包括：}
\OperatorTok{{-}} \OperatorTok{**}\NormalTok{ReLU}\OperatorTok{**}\NormalTok{：\textasciigrave{}\sigma(x) }\OperatorTok{=} \BuiltInTok{max}\NormalTok{(}\DecValTok{0}\NormalTok{, x)\textasciigrave{}，在x\geq}\DecValTok{0}\NormalTok{时导数为}\DecValTok{1}\NormalTok{，x}\OperatorTok{\textless{}}\DecValTok{0}\NormalTok{时导数为}\DecValTok{0}\NormalTok{，满足}\DecValTok{1}\OperatorTok{{-}}\NormalTok{Lipschitz条件}
\OperatorTok{{-}} \OperatorTok{**}\NormalTok{tanh}\OperatorTok{**}\NormalTok{：双曲正切函数，其导数的绝对值始终小于等于}\DecValTok{1}
\OperatorTok{{-}} \OperatorTok{**}\NormalTok{Sigmoid}\OperatorTok{**}\NormalTok{：经过适当缩放后也可满足条件}
\OperatorTok{{-}} \OperatorTok{**}\NormalTok{Leaky ReLU}\OperatorTok{**}\NormalTok{：当负斜率绝对值\leq}\DecValTok{1}\NormalTok{时满足条件}

\OperatorTok{**}\NormalTok{定理}\FloatTok{13.1}\NormalTok{（非线性MQR的收敛性）}\OperatorTok{**}\NormalTok{：考虑非线性MQR动力学系统：}
\end{Highlighting}
\end{Shaded}

h\_\{t+1\} = \sigma((1-\alpha)h\_t H\^{}T + \alphaJ(x))

\begin{verbatim}
其中\sigma是1-Lipschitz连续函数，H是双随机矩阵（`H = |U|^2`, U∈U(N)），\alpha∈(0,1)。则该系统是压缩映射，存在唯一不动点h*，且从任意初始状态h₀出发的迭代序列{h_t}收敛到h*。

**证明概要**：定义映射T(h) = \sigma((1-\alpha)h H^T + \alphaJ(x))。对于任意两个状态h₁, h₂，有：
\end{verbatim}

\textbar\textbar T(h₁) - T(h₂)\textbar\textbar{} \leq
\textbar\textbar(1-\alpha)(h₁ - h₂)H\^{}T\textbar\textbar{}
（\sigma的1-Lipschitz性） \leq (1-\alpha)\textbar\textbar h₁ -
h₂\textbar\textbar·\textbar\textbar H\^{}T\textbar\textbar{} \leq
(1-\alpha)\textbar\textbar h₁ - h₂\textbar\textbar{}
（因为双随机矩阵的谱范数\leq1）

\begin{verbatim}
由于(1-\alpha)<1，T是压缩映射，由Banach不动点定理可知存在唯一不动点，且迭代收敛。

### 13.2.3 实现代码解析

在MQR的实现中，非线性扩展通过`nonlinearity`参数控制：

```python
# mqr/dynamics/ring_dynamics.py 中的相关代码片段

class RingDynamics(nn.Module):
    def __init__(self, ring_size, alpha=0.5, nonlinearity='none', 
                 dynamics_mode='unistochastic', **kwargs):
        super().__init__()
        self.ring_size = ring_size
        self.alpha = alpha
        self.nonlinearity = nonlinearity
        self.dynamics_mode = dynamics_mode
        
        # 初始化激活函数
        if nonlinearity == 'relu':
            self.act = nn.ReLU()
        elif nonlinearity == 'tanh':
            self.act = nn.Tanh()
        elif nonlinearity == 'sigmoid':
            self.act = nn.Sigmoid()
        elif nonlinearity == 'none' or nonlinearity is None:
            self.act = lambda x: x  # 恒等映射
        else:
            raise ValueError(f"Unsupported nonlinearity: {nonlinearity}")
    
    def forward(self, h, H, injection):
        """
        执行一步环形动力学更新
        
        参数:
            h: 当前环形状态 [batch_size, ring_size, hidden_dim]
            H: 双随机连接矩阵 [ring_size, ring_size]
            injection: 注入信号 [batch_size, ring_size, hidden_dim]
        
        返回:
            更新后的环形状态
        """
        # 线性组合部分
        linear_combination = (1 - self.alpha) * torch.matmul(h, H.t()) + self.alpha * injection
        
        # 应用非线性（如果是恒等映射则无效果）
        h_next = self.act(linear_combination)
        
        return h_next
    
    def converge_to_fixed_point(self, H, injection, max_iter=100, tol=1e-6):
        """
        迭代收敛到固定点
        
        参数:
            H: 双随机连接矩阵
            injection: 注入信号
            max_iter: 最大迭代次数
            tol: 收敛容差
        
        返回:
            固定点状态 h_star
        """
        batch_size = injection.shape[0]
        h = torch.zeros_like(injection)  # 初始状态
        
        for i in range(max_iter):
            h_next = self.forward(h, H, injection)
            
            # 检查收敛
            diff = torch.norm(h_next - h, dim=(1, 2)).max().item()
            h = h_next
            
            if diff < tol:
                # 记录实际迭代次数用于分析
                self.last_convergence_iter = i + 1
                break
        
        return h
\end{verbatim}

非线性扩展的引入显著增强了MQR的表达能力。实验表明，在需要复杂模式识别的任务中（如自然语言理解中的语义组合、图像分类中的细粒度特征提取），带有适当非线性的MQR通常比纯线性版本获得更高的准确率。然而，这也带来了权衡：非线性可能使收敛速度略微减慢，并增加了梯度计算的计算量。

\hypertarget{ux590dux6570ux57dfux9149ux63a8ux7406ux5229ux7528ux76f8ux4f4dux4fe1ux606fux589eux5f3aux8868ux793a}{%
\subsection{13.3
复数域酉推理：利用相位信息增强表示}\label{ux590dux6570ux57dfux9149ux63a8ux7406ux5229ux7528ux76f8ux4f4dux4fe1ux606fux589eux5f3aux8868ux793a}}

\hypertarget{ux4eceux53ccux968fux673aux5230ux9149ux77e9ux9635ux7684ux76f4ux63a5ux4f7fux7528}{%
\subsubsection{13.3.1
从双随机到酉矩阵的直接使用}\label{ux4eceux53ccux968fux673aux5230ux9149ux77e9ux9635ux7684ux76f4ux63a5ux4f7fux7528}}

基础MQR使用幺模随机矩阵\texttt{H\ =\ \textbar{}U\textbar{}²}作为连接矩阵，其中U是酉矩阵。这种设计确保了H的双随机性，从而保证了实数域动力学的收敛性。然而，这一过程丢弃了酉矩阵U的相位信息------仅使用其元素的模平方。相位信息在量子力学、信号处理和复数神经网络中被证明携带重要信息。

复数域酉推理（Complex Unitary
Inference）是MQR的一个高级特性，它直接使用酉矩阵U（而非其模平方）在复数域执行环形动力学。在这种模式下，环形状态h变为复数张量，动力学方程修改为：

\begin{verbatim}
h_{t+1} = (1-\alpha)h_t U^^{\dagger} + \alphaJ_C(x)
\end{verbatim}

其中U\^{}^{\dagger}是U的共轭转置（厄米共轭），J\_C(x)是复数域注入信号。由于酉矩阵保持复数向量的L2范数（\texttt{\textbar{}\textbar{}Uv\textbar{}\textbar{}\ =\ \textbar{}\textbar{}v\textbar{}\textbar{}}），该动力学在复数域同样是稳定的。

\hypertarget{ux590dux6570ux6ce8ux5165ux4e0eux6d4bux91cfux8bfbux51fa}{%
\subsubsection{13.3.2
复数注入与测量读出}\label{ux590dux6570ux6ce8ux5165ux4e0eux6d4bux91cfux8bfbux51fa}}

在复数域推理中，注入过程也需要适应复数运算。MQR提供了两种策略：

\begin{enumerate}
\def\labelenumi{\arabic{enumi}.}
\tightlist
\item
  \textbf{实数注入，复数转换}：保持注入算子J(x)为实数，然后通过简单扩展（如添加零虚部）或通过可学习的复数变换将其转换为复数。
\item
  \textbf{全复数注入}：使用复数权重矩阵实现注入算子，即\texttt{J\_C(x)\ =\ (W\_up\_real\ +\ i·W\_up\_imag)\ ·\ \sigma((W\_down\_real\ +\ i·W\_down\_imag)\ ·\ x)}。
\end{enumerate}

复数环形状态h*的读出需要特殊处理，因为下游任务通常需要实数输出。MQR提供了多种测量（Measurement）策略：

\begin{itemize}
\tightlist
\item
  \textbf{模测量}：\texttt{h\_real\ =\ \textbar{}h*\textbar{}}，取每个复数元素的模
\item
  \textbf{实部测量}：\texttt{h\_real\ =\ Re(h*)}，取实部
\item
  \textbf{虚部测量}：\texttt{h\_real\ =\ Im(h*)}，取虚部
\item
  \textbf{相位测量}：\texttt{h\_real\ =\ arg(h*)}，取相位角
\item
  \textbf{混合测量}：结合多种测量，如\texttt{{[}\textbar{}h*\textbar{},\ arg(h*){]}}拼接
\end{itemize}

\begin{figure}
\centering
\includegraphics{figures/diagram_12.png}
\caption{图 12}
\end{figure}

\hypertarget{ux590dux6570mqrux7684ux6570ux5b66ux6027ux8d28ux4e0eux5b9eux73b0}{%
\subsubsection{13.3.3
复数MQR的数学性质与实现}\label{ux590dux6570mqrux7684ux6570ux5b66ux6027ux8d28ux4e0eux5b9eux73b0}}

\textbf{定理13.2（复数MQR的收敛性）}：考虑复数MQR动力学系统：

\begin{verbatim}
h_{t+1} = (1-\alpha)h_t U^^{\dagger} + \alphaJ_C(x)
\end{verbatim}

其中U∈U(N)是酉矩阵，\alpha∈(0,1)，J\_C(x)是复数注入信号。则该系统是压缩映射，存在唯一复数不动点h\emph{，且从任意初始状态h₀出发的迭代序列\{h\_t\}收敛到h}。

\textbf{证明}：定义映射T\_C(h) = (1-\alpha)h U\^{}^{\dagger} +
\alphaJ\_C(x)。对于任意两个复数状态h₁, h₂：

\begin{verbatim}
||T_C(h₁) - T_C(h₂)|| = ||(1-\alpha)(h₁ - h₂)U^^{\dagger}||
                     = (1-\alpha)||h₁ - h₂||·||U^^{\dagger}||
                     = (1-\alpha)||h₁ - h₂||  （因为酉矩阵的谱范数为1）
\end{verbatim}

故T\_C是压缩比为(1-\alpha)的压缩映射，收敛性得证。

在代码实现中，复数MQR通过\texttt{dynamics\_mode=\textquotesingle{}unitary\textquotesingle{}}参数启用：

\begin{Shaded}
\begin{Highlighting}[]
\CommentTok{\# mqr/dynamics/complex\_dynamics.py}

\KeywordTok{class}\NormalTok{ ComplexRingDynamics(nn.Module):}
    \KeywordTok{def} \FunctionTok{\_\_init\_\_}\NormalTok{(}\VariableTok{self}\NormalTok{, ring\_size, alpha}\OperatorTok{=}\FloatTok{0.5}\NormalTok{, measurement}\OperatorTok{=}\StringTok{\textquotesingle{}magnitude\textquotesingle{}}\NormalTok{, }\OperatorTok{**}\NormalTok{kwargs):}
        \BuiltInTok{super}\NormalTok{().}\FunctionTok{\_\_init\_\_}\NormalTok{()}
        \VariableTok{self}\NormalTok{.ring\_size }\OperatorTok{=}\NormalTok{ ring\_size}
        \VariableTok{self}\NormalTok{.alpha }\OperatorTok{=}\NormalTok{ alpha}
        \VariableTok{self}\NormalTok{.measurement }\OperatorTok{=}\NormalTok{ measurement}
        
    \KeywordTok{def}\NormalTok{ forward(}\VariableTok{self}\NormalTok{, h\_complex, U, injection\_complex):}
        \CommentTok{"""}
\CommentTok{        复数域动力学更新}
\CommentTok{        }
\CommentTok{        参数:}
\CommentTok{            h\_complex: 复数环形状态 [batch\_size, ring\_size, hidden\_dim, 2]}
\CommentTok{                      最后一维的2表示实部和虚部}
\CommentTok{            U: 酉矩阵 [ring\_size, ring\_size, 2] (实部+虚部)}
\CommentTok{            injection\_complex: 复数注入 [batch\_size, ring\_size, hidden\_dim, 2]}
\CommentTok{        }
\CommentTok{        返回:}
\CommentTok{            更新后的复数状态}
\CommentTok{        """}
        \CommentTok{\# 将实部虚部分解}
\NormalTok{        h\_real, h\_imag }\OperatorTok{=}\NormalTok{ h\_complex[..., }\DecValTok{0}\NormalTok{], h\_complex[..., }\DecValTok{1}\NormalTok{]}
\NormalTok{        U\_real, U\_imag }\OperatorTok{=}\NormalTok{ U[..., }\DecValTok{0}\NormalTok{], U[..., }\DecValTok{1}\NormalTok{]}
\NormalTok{        inj\_real, inj\_imag }\OperatorTok{=}\NormalTok{ injection\_complex[..., }\DecValTok{0}\NormalTok{], injection\_complex[..., }\DecValTok{1}\NormalTok{]}
        
        \CommentTok{\# 计算 h * U\^{}^{\dagger} = h * U\^{}H}
        \CommentTok{\# U\^{}H 是U的共轭转置：实部转置，虚部转置并取负}
\NormalTok{        UH\_real }\OperatorTok{=}\NormalTok{ U\_real.t()}
\NormalTok{        UH\_imag }\OperatorTok{=} \OperatorTok{{-}}\NormalTok{U\_imag.t()}
        
        \CommentTok{\# 复数矩阵乘法: (h\_real + i·h\_imag) * (UH\_real + i·UH\_imag)}
        \CommentTok{\# = (h\_real·UH\_real {-} h\_imag·UH\_imag) + i·(h\_real·UH\_imag + h\_imag·UH\_real)}
\NormalTok{        term\_real }\OperatorTok{=}\NormalTok{ torch.matmul(h\_real, UH\_real) }\OperatorTok{{-}}\NormalTok{ torch.matmul(h\_imag, UH\_imag)}
\NormalTok{        term\_imag }\OperatorTok{=}\NormalTok{ torch.matmul(h\_real, UH\_imag) }\OperatorTok{+}\NormalTok{ torch.matmul(h\_imag, UH\_real)}
        
        \CommentTok{\# 应用动力学方程}
\NormalTok{        h\_next\_real }\OperatorTok{=}\NormalTok{ (}\DecValTok{1} \OperatorTok{{-}} \VariableTok{self}\NormalTok{.alpha) }\OperatorTok{*}\NormalTok{ term\_real }\OperatorTok{+} \VariableTok{self}\NormalTok{.alpha }\OperatorTok{*}\NormalTok{ inj\_real}
\NormalTok{        h\_next\_imag }\OperatorTok{=}\NormalTok{ (}\DecValTok{1} \OperatorTok{{-}} \VariableTok{self}\NormalTok{.alpha) }\OperatorTok{*}\NormalTok{ term\_imag }\OperatorTok{+} \VariableTok{self}\NormalTok{.alpha }\OperatorTok{*}\NormalTok{ inj\_imag}
        
        \CommentTok{\# 组合实部虚部}
\NormalTok{        h\_next }\OperatorTok{=}\NormalTok{ torch.stack([h\_next\_real, h\_next\_imag], dim}\OperatorTok{={-}}\DecValTok{1}\NormalTok{)}
        
        \ControlFlowTok{return}\NormalTok{ h\_next}
    
    \KeywordTok{def}\NormalTok{ measure(}\VariableTok{self}\NormalTok{, h\_complex):}
        \CommentTok{"""}
\CommentTok{        将复数状态转换为实数表示}
\CommentTok{        }
\CommentTok{        参数:}
\CommentTok{            h\_complex: 复数状态 [..., 2]}
\CommentTok{        }
\CommentTok{        返回:}
\CommentTok{            实数表示}
\CommentTok{        """}
        \ControlFlowTok{if} \VariableTok{self}\NormalTok{.measurement }\OperatorTok{==} \StringTok{\textquotesingle{}magnitude\textquotesingle{}}\NormalTok{:}
            \CommentTok{\# 模: sqrt(real\^{}2 + imag\^{}2)}
\NormalTok{            real, imag }\OperatorTok{=}\NormalTok{ h\_complex[..., }\DecValTok{0}\NormalTok{], h\_complex[..., }\DecValTok{1}\NormalTok{]}
            \ControlFlowTok{return}\NormalTok{ torch.sqrt(real}\OperatorTok{**}\DecValTok{2} \OperatorTok{+}\NormalTok{ imag}\OperatorTok{**}\DecValTok{2} \OperatorTok{+} \FloatTok{1e{-}8}\NormalTok{)}
        
        \ControlFlowTok{elif} \VariableTok{self}\NormalTok{.measurement }\OperatorTok{==} \StringTok{\textquotesingle{}real\textquotesingle{}}\NormalTok{:}
            \CommentTok{\# 实部}
            \ControlFlowTok{return}\NormalTok{ h\_complex[..., }\DecValTok{0}\NormalTok{]}
        
        \ControlFlowTok{elif} \VariableTok{self}\NormalTok{.measurement }\OperatorTok{==} \StringTok{\textquotesingle{}imag\textquotesingle{}}\NormalTok{:}
            \CommentTok{\# 虚部}
            \ControlFlowTok{return}\NormalTok{ h\_complex[..., }\DecValTok{1}\NormalTok{]}
        
        \ControlFlowTok{elif} \VariableTok{self}\NormalTok{.measurement }\OperatorTok{==} \StringTok{\textquotesingle{}phase\textquotesingle{}}\NormalTok{:}
            \CommentTok{\# 相位: atan2(imag, real)}
\NormalTok{            real, imag }\OperatorTok{=}\NormalTok{ h\_complex[..., }\DecValTok{0}\NormalTok{], h\_complex[..., }\DecValTok{1}\NormalTok{]}
            \ControlFlowTok{return}\NormalTok{ torch.atan2(imag, real)}
        
        \ControlFlowTok{elif} \VariableTok{self}\NormalTok{.measurement }\OperatorTok{==} \StringTok{\textquotesingle{}concat\textquotesingle{}}\NormalTok{:}
            \CommentTok{\# 拼接模和相位}
\NormalTok{            real, imag }\OperatorTok{=}\NormalTok{ h\_complex[..., }\DecValTok{0}\NormalTok{], h\_complex[..., }\DecValTok{1}\NormalTok{]}
\NormalTok{            magnitude }\OperatorTok{=}\NormalTok{ torch.sqrt(real}\OperatorTok{**}\DecValTok{2} \OperatorTok{+}\NormalTok{ imag}\OperatorTok{**}\DecValTok{2} \OperatorTok{+} \FloatTok{1e{-}8}\NormalTok{)}
\NormalTok{            phase }\OperatorTok{=}\NormalTok{ torch.atan2(imag, real)}
            \ControlFlowTok{return}\NormalTok{ torch.cat([magnitude, phase], dim}\OperatorTok{={-}}\DecValTok{1}\NormalTok{)}
        
        \ControlFlowTok{else}\NormalTok{:}
            \ControlFlowTok{raise} \PreprocessorTok{ValueError}\NormalTok{(}\SpecialStringTok{f"Unknown measurement type: }\SpecialCharTok{\{}\VariableTok{self}\SpecialCharTok{.}\NormalTok{measurement}\SpecialCharTok{\}}\SpecialStringTok{"}\NormalTok{)}
\end{Highlighting}
\end{Shaded}

复数域酉推理在需要相位敏感性的任务中表现出色，例如： -
\textbf{信号处理}：音频、雷达信号中的相位信息至关重要 -
\textbf{量子系统模拟}：自然处理复数振幅和相位 -
\textbf{某些类型的时序预测}：相位可以编码周期性模式 -
\textbf{对抗鲁棒性}：复数表示可能提供额外的鲁棒性

\hypertarget{ux53ccux9149ux77e9ux9635ux89e3ux8026ux5206ux79bbux4e16ux754cux6a21ux578bux4e0eux7b56ux7565}{%
\subsection{13.4
双酉矩阵解耦：分离世界模型与策略}\label{ux53ccux9149ux77e9ux9635ux89e3ux8026ux5206ux79bbux4e16ux754cux6a21ux578bux4e0eux7b56ux7565}}

\hypertarget{ux89e3ux8026ux52a8ux673aux4e0eux54f2ux5b66}{%
\subsubsection{13.4.1
解耦动机与哲学}\label{ux89e3ux8026ux52a8ux673aux4e0eux54f2ux5b66}}

在强化学习、自适应控制以及需要在线学习的场景中，一个理想的系统应该能够区分``世界如何运作''（世界模型）和``如何行动''（策略）。世界模型应该相对稳定，反映环境的基本动力学；而策略应该灵活可调，以适应不同的目标或任务。

双酉矩阵解耦（Dual-Unitary
Decoupling）将MQR的连接酉矩阵分解为两个部分的乘积：

\begin{verbatim}
U_total = U_policy · U_base
\end{verbatim}

其中： -
\texttt{U\_base}（基础酉矩阵）作为冻结的``世界模型''，编码环境或任务空间的基本结构
-
\texttt{U\_policy}（策略酉矩阵）作为可学习的``策略''，编码特定任务或目标下的适应性行为

这种分解的哲学灵感来源于： 1.
\textbf{系统辨识与控制理论}：分离被控对象模型与控制器设计 2.
\textbf{迁移学习}：固定特征提取器，微调顶层分类器 3.
\textbf{人脑功能分离}：相对稳定的神经结构与可塑的突触连接

\hypertarget{ux6570ux5b66ux5f62ux5f0fux4e0eux4f18ux5316ux7b56ux7565}{%
\subsubsection{13.4.2
数学形式与优化策略}\label{ux6570ux5b66ux5f62ux5f0fux4e0eux4f18ux5316ux7b56ux7565}}

设基础酉矩阵U\_base通过某种方式初始化并冻结（不参与梯度更新），策略酉矩阵U\_policy随机初始化并参与训练。总酉矩阵为：

\begin{verbatim}
U_total = U_policy · U_base
\end{verbatim}

对应的双随机连接矩阵为：

\begin{verbatim}
H_total = |U_total|² = |U_policy · U_base|²
\end{verbatim}

由于酉矩阵乘法封闭性（酉矩阵的乘积仍是酉矩阵），U\_total保证是酉矩阵，因此H\_total保证是双随机矩阵，收敛性定理仍然成立。

在优化过程中，只有U\_policy的参数通过梯度下降更新，U\_base保持不变

\hypertarget{ux7b2c14ux7ae0-ux83abux6bd4ux4e4cux65afux91cfux5b50ux73afux7f51ux7edcux7684ux5b9eux9a8cux8bc4ux4f30ux4e0eux6027ux80fdux5206ux6790}{%
\section{第14章
莫比乌斯量子环网络的实验评估与性能分析}\label{ux7b2c14ux7ae0-ux83abux6bd4ux4e4cux65afux91cfux5b50ux73afux7f51ux7edcux7684ux5b9eux9a8cux8bc4ux4f30ux4e0eux6027ux80fdux5206ux6790}}

\hypertarget{ux5b9eux9a8cux8bbeux8ba1ux4e0eux57faux51c6ux6d4bux8bd5ux6846ux67b6}{%
\subsection{14.1
实验设计与基准测试框架}\label{ux5b9eux9a8cux8bbeux8ba1ux4e0eux57faux51c6ux6d4bux8bd5ux6846ux67b6}}

为了全面评估莫比乌斯量子环网络的理论优势与实用价值，我们设计了一套系统性的实验评估框架。该框架旨在从多个维度验证模型的核心特性：长期依赖建模能力、训练稳定性、计算效率、以及在不同任务领域的泛化性能。我们的实验设计遵循``控制变量''原则，在公平对比的基础上，深入分析MQR架构的独特贡献。

\hypertarget{ux57faux51c6ux4efbux52a1ux9009ux62e9}{%
\subsubsection{14.1.1
基准任务选择}\label{ux57faux51c6ux4efbux52a1ux9009ux62e9}}

我们选取了四类具有代表性的序列建模任务，它们分别对模型提出了不同的挑战：

\begin{enumerate}
\def\labelenumi{\arabic{enumi}.}
\tightlist
\item
  \textbf{合成序列任务（Synthetic Sequential
  Tasks）}：这类任务具有清晰的理论预期和可验证的长期依赖结构，用于精确诊断模型的能力边界。

  \begin{itemize}
  \tightlist
  \item
    \textbf{复制记忆任务（Copy Memory
    Task）}：模型需要记忆一个在序列早期出现的模式，并在序列末尾的特定指令后将其精确复制输出。该任务直接测试模型对长程信息的存储与提取能力。
  \item
    \textbf{添加任务（Adding
    Task）}：输入序列包含两通道，一通道为随机数，另一通道为指示标记。模型需要在序列末尾输出两个由标记指定的随机数之和。此任务考验模型对稀疏、长距离相关信息的筛选与整合能力。
  \item
    \textbf{顺序MNIST/pMNIST（Sequential MNIST/permuted
    MNIST）}：将MNIST图像的像素按行扫描（sMNIST）或随机置换顺序（pMNIST）作为序列输入，进行分类。sMNIST测试局部到整体的结构建模，而pMNIST因破坏了空间局部性，对模型捕捉长距离、非局部依赖的能力要求更高。
  \end{itemize}
\item
  \textbf{自然语言处理任务（Natural Language Processing Tasks）}：

  \begin{itemize}
  \tightlist
  \item
    \textbf{语言建模（Language Modeling）}：在Penn
    Treebank（PTB）和WikiText-2（WT2）数据集上评估模型的词级困惑度（Perplexity,
    PPL）。语言建模是检验序列模型生成与预测能力的核心任务。
  \item
    \textbf{文本分类（Text Classification）}：在IMDB电影评论、AG
    News等数据集上进行情感或主题分类，评估模型从变长文本中提取全局语义表征的能力。
  \end{itemize}
\item
  \textbf{时间序列预测任务（Time Series Forecasting Tasks）}：

  \begin{itemize}
  \tightlist
  \item
    \textbf{电力负荷预测（Electricity Load
    Forecasting）}：使用公开的电力消耗数据集，预测未来多个时间点的负荷值。此任务涉及复杂的季节性、趋势性和外部依赖。
  \item
    \textbf{多变量时间序列预测（Multivariate Time
    Series）}：在交通流量、气象等数据集上进行预测，考验模型处理高维、相互关联的时序信号的能力。
  \end{itemize}
\item
  \textbf{算法学习任务（Algorithmic Learning Tasks）}：

  \begin{itemize}
  \tightlist
  \item
    \textbf{程序评估（Program
    Evaluation）}：学习执行简单编程语言（如前缀算术表达式）的求值。这要求模型精确模拟栈、寄存器等计算结构，对推理的精确性和稳定性要求极高。
  \end{itemize}
\end{enumerate}

\hypertarget{ux5bf9ux6bd4ux6a21ux578bux8bbeux7f6e}{%
\subsubsection{14.1.2
对比模型设置}\label{ux5bf9ux6bd4ux6a21ux578bux8bbeux7f6e}}

为了进行有意义的对比，我们选取了一系列具有代表性的基线模型和当前先进的序列模型：

\begin{itemize}
\tightlist
\item
  \textbf{经典RNN变体}：Vanilla
  RNN、长短期记忆网络、门控循环单元。它们代表了通过工程化门控机制解决长程依赖问题的实践路线。
\item
  \textbf{理论驱动RNN}：正交RNN（使用Cayley变换参数化）、酉循环神经网络。它们代表了通过严格的数学约束（正交/酉性）来保证稳定性的理论路线，是MQR最直接的理论参照。
\item
  \textbf{注意力机制模型}：Transformer（仅编码器部分）。作为当前主导的序列建模架构，其基于自注意力的并行化计算是另一种解决长程依赖的范式。我们将对比其在中等长度序列任务上的性能与效率。
\item
  \textbf{其他稳定RNN结构}：如反时间线性RNN、指数RNN等近年提出的具有理论保证的变体。
\end{itemize}

所有对比实验均在相同的硬件环境（NVIDIA A100 GPU）、软件框架（PyTorch
2.0+）下进行。我们使用网格搜索或贝叶斯优化对每个模型的关键超参数（如隐藏层维度、学习率、梯度裁剪阈值等）进行调优，并报告其最佳性能。对于MQR模型，我们系统地探索了其核心配置选项的影响，如环尺寸N、松弛迭代步数T、注入秩r、是否启用非线性扩展、是否使用复数动力学等。

\hypertarget{ux8bc4ux4f30ux6307ux6807}{%
\subsubsection{14.1.3 评估指标}\label{ux8bc4ux4f30ux6307ux6807}}

我们采用任务相关的标准评估指标： *
\textbf{分类/回归任务}：准确率（Accuracy）、均方误差（MSE）、平均绝对误差（MAE）。
* \textbf{语言建模任务}：困惑度（PPL），越低越好。 *
\textbf{算法任务}：精确匹配准确率（Exact Match Accuracy）。 *
\textbf{训练动态}：训练损失曲线、验证损失曲线、梯度范数随时间的变化，用于观察稳定性和收敛速度。
* \textbf{计算效率}：单步推理时间（毫秒）、训练周期时间、GPU内存占用。 *
\textbf{模型复杂度}：可训练参数数量。

\hypertarget{ux6838ux5fc3ux7279ux6027ux9a8cux8bc1ux5b9eux9a8c}{%
\subsection{14.2
核心特性验证实验}\label{ux6838ux5fc3ux7279ux6027ux9a8cux8bc1ux5b9eux9a8c}}

本节通过精心设计的实验，直接验证MQR架构所宣称的核心理论特性。

\hypertarget{ux957fux671fux4f9dux8d56ux5efaux6a21ux80fdux529b}{%
\subsubsection{14.2.1
长期依赖建模能力}\label{ux957fux671fux4f9dux8d56ux5efaux6a21ux80fdux529b}}

我们在复制记忆任务和添加任务上进行了系统测试。如图14.1所示，我们设置了不同的序列长度（从100到1000+），以观察模型性能随依赖距离增长而衰减的情况。

\begin{figure}
\centering
\includegraphics{figures/diagram_13.png}
\caption{图 13}
\end{figure}

\textbf{图14.1：复制记忆任务中MQR的信息流示意图。}
早期模式通过注入算子J(x)被编码并驱动环形系统，在弛豫过程中，信息通过双随机连接矩阵H在整个环上扩散并达到平衡。最终，该信息在读出阶段被有效提取。

\textbf{实验结果分析}： * Vanilla
RNN在序列长度超过50后性能急剧下降，证实了其梯度消失问题。 *
LSTM和GRU在中等长度（\textasciitilde200）序列上表现稳健，但在超长序列（\textgreater500）上，其门控机制对信息的维持能力开始减弱，准确率缓慢下降。
*
正交RNN和URNN在整个长度范围内都表现出优异的稳定性，验证了等距变换在维持梯度流方面的理论优势。
*
\textbf{MQR模型在所有测试长度上均达到了接近100\%的复制准确率}，其性能曲线几乎是一条水平线。这强有力地证明了``弛豫到稳态''的动力学机制在理论上不受序列长度限制。环形系统中的双随机连接确保了信息不会在传播中衰减（能量守恒），而固定点迭代过程使得无论初始模式被``注入''到多久之前，系统最终都能收敛到一个包含该信息的唯一平衡态。这解决了传统RNN因时间展开而必然面临的梯度传播衰减难题。

在pMNIST任务上，MQR同样取得了领先的性能（测试准确率
\textgreater98\%），显著优于普通RNN，并与最先进的正交RNN和精心调优的LSTM相当或略有优势。这表明MQR在捕捉复杂、非局部像素依赖方面同样有效。

\hypertarget{ux8badux7ec3ux7a33ux5b9aux6027ux5206ux6790}{%
\subsubsection{14.2.2
训练稳定性分析}\label{ux8badux7ec3ux7a33ux5b9aux6027ux5206ux6790}}

我们通过监控训练过程中梯度范数、权重更新幅值以及损失曲线的平滑度来评估稳定性。

\textbf{梯度范数分析}：我们记录了在整个训练周期内，模型主要参数（如酉矩阵U、注入权重W\_up/W\_down）的梯度L2范数。如图14.2所示（以PTB语言建模任务为例），MQR的梯度范数在整个训练过程中保持在一个相对平稳、有界的范围内，没有出现剧烈的峰值（爆炸）或持续趋近于零（消失）。相比之下，Vanilla
RNN的梯度范数波动剧烈，且早期容易出现爆炸；LSTM和GRU较为稳定，但仍可能出现周期性尖峰。正交RNN的梯度最为平稳，MQR与之类似。

\textbf{损失曲线分析}：MQR的训练损失和验证损失曲线平滑下降，没有出现明显的震荡或突然上升（过拟合）后又下降的现象。这表明基于固定点迭代的推理过程产生了平滑的、易于优化的损失景观。反向传播通过隐函数定理进行，避免了在深度时间展开图中进行链式求导，从而规避了梯度乘积累积效应这一不稳定根源。

\textbf{权重动力学}：我们可视化了连接矩阵H（由U导出）在训练过程中的特征值分布变化。可以发现，其特征值始终位于单位圆内（对于实数H）或模为1（对于复数U），且分布逐渐演化以适应任务需求，这从底层动力学角度解释了其稳定性。

\hypertarget{ux8ba1ux7b97ux6548ux7387ux4e0eux53efux6269ux5c55ux6027}{%
\subsubsection{14.2.3
计算效率与可扩展性}\label{ux8ba1ux7b97ux6548ux7387ux4e0eux53efux6269ux5c55ux6027}}

MQR的前向推理由两部分组成：一次性的注入计算和T步的弛豫迭代。虽然迭代过程在序列长度上看似是O(T)的，但与RNN的O(L)时间展开有本质不同：
1. \textbf{并行性}：弛豫迭代 \texttt{h\_\{t+1\}\ =\ h\_t\ H\^{}T}
（或含注入的更新）本质是矩阵乘法，在GPU上可以高效并行计算。每一步迭代的成本是O(N\^{}2)（或通过低秩分解、FFT等优化为更低复杂度），与序列长度L无关。
2.
\textbf{迭代步数T}：T是一个超参数，通常远小于需要建模的序列有效长度L。在实践中，对于许多任务，T在20\textasciitilde100步之间系统已能充分收敛。我们引入了收敛判据（如相邻h变化的范数小于阈值\epsilon），允许动态提前终止迭代，进一步节省计算。

我们在不同环尺寸N和序列长度L下测试了MQR的推理速度。当N固定时，MQR的推理时间几乎不随L增长而增加（因为注入计算只做一次），而RNN/LSTM的时间则与L线性相关。这使得MQR在处理极长序列时具有潜在优势。当L固定而N增加时，由于矩阵乘法开销增大，时间会增加。然而，通过采用\textbf{局部投影采样读出}，我们只计算环上一小部分节点（如前k个）的输出，避免了全环N维向量的全连接读出，这大幅减少了最终输出层的参数和计算量，提升了效率。

内存占用方面，MQR不需要保存所有时间步的中间隐藏状态（如同RNN的BPTT所需），只需在弛豫迭代中维护当前状态h，因此其训练内存消耗为O(N
* batch\_size)，与序列长度L无关，这对于超长序列训练至关重要。

\hypertarget{ux9ad8ux7ea7ux529fux80fdux6d88ux878dux5b9eux9a8c}{%
\subsection{14.3
高级功能消融实验}\label{ux9ad8ux7ea7ux529fux80fdux6d88ux878dux5b9eux9a8c}}

本章节对MQR中引入的高级可选功能进行消融研究，量化它们对性能的具体贡献。

\hypertarget{ux590dux6570ux9149ux52a8ux529bux5b66-vs.-ux5b9eux6570ux53ccux968fux673aux52a8ux529bux5b66}{%
\subsubsection{14.3.1 复数酉动力学
vs.~实数双随机动力学}\label{ux590dux6570ux9149ux52a8ux529bux5b66-vs.-ux5b9eux6570ux53ccux968fux673aux52a8ux529bux5b66}}

我们在几个具有潜在相位信息或振荡模式的任务上对比了
\texttt{dynamics\_mode=‘unitary’}（复数）和
\texttt{dynamics\_mode=‘doubly\_stochastic’}（实数）两种模式。 *
\textbf{任务}：合成振荡信号预测、包含复数特征的物理系统模拟任务。 *
\textbf{结果}：在大多数标准分类和回归任务上，两种模式性能相近，实数模式略快。然而，在明确涉及旋转、干涉或需要保持相位一致性的任务中，复数酉动力学模式表现出显著优势。复数域的操作允许信息以相位形式在环中传播，提供了更丰富的表示空间。读出时对稳态h*取模或使用复数权重，能够利用这一相位信息。这验证了将动力学推广到复数域的实用价值。

\hypertarget{ux975eux7ebfux6027ux6269ux5c55ux7684ux5f71ux54cd}{%
\subsubsection{14.3.2
非线性扩展的影响}\label{ux975eux7ebfux6027ux6269ux5c55ux7684ux5f71ux54cd}}

我们评估了在注入端（\texttt{inject\_nonlinearity}）和弛豫过程中（\texttt{relax\_nonlinearity}，使用1-Lipschitz激活如ReLU的缩放版本）加入非线性的效果。
*
\textbf{结果}：在图像分类（如MNIST）和某些NLP任务上，注入端的非线性带来了明显的性能提升，因为它增加了注入变换的表达能力。弛豫过程中的非线性（需谨慎选择以满足收敛条件）在部分复杂任务上能略微提升性能，但有时会略微减慢收敛速度。总体而言，\textbf{非线性扩展是一个有效的``表达力增强''工具}，它使MQR能够突破纯线性双随机动力学的限制，学习更复杂的变换，同时理论证明指出在满足条件（如1-Lipschitz）下不会破坏系统的收敛保证。这体现了MQR框架在保持理论纯洁性与追求实用性能之间的良好平衡。

\hypertarget{ux53ccux9149ux89e3ux8026ux4e0eux53efux5b66ux4e60ux76eeux6807ux5e73ux8861ux6001}{%
\subsubsection{14.3.3
双酉解耦与可学习目标平衡态}\label{ux53ccux9149ux89e3ux8026ux4e0eux53efux5b66ux4e60ux76eeux6807ux5e73ux8861ux6001}}

我们在需要快速适应新任务或多任务学习的场景下测试了这些功能。 *
\textbf{双酉解耦}：将U分解为\texttt{U\_policy\ @\ U\_base}。我们在一个持续学习设置中，先在一个基础任务上训练\texttt{U\_base}并冻结，然后在相关但不同的新任务上仅训练\texttt{U\_policy}。实验表明，与从头训练或微调整个U相比，\textbf{双酉解耦策略能实现更快的新任务收敛，并显著减轻对旧任务的遗忘}。\texttt{U\_base}充当了一个稳定的``世界动力学''先验，而\texttt{U\_policy}进行快速适应，这为模块化、可组合的序列建模提供了新思路。
*
\textbf{可学习目标平衡态}：在分类任务中，为每个类别引入一个可学习的\texttt{p\_y}。训练过程中，除了标准分类损失，还添加一个正则项，鼓励输入x对应的稳态\texttt{h*}与真实类别y的\texttt{p\_y}在环状态分布上接近。实验显示，这种方法不仅能略微提升分类准确率，更重要的是，\textbf{它使得模型学习到的环形稳态在语义上更具可解释性}，同一类别的样本会收敛到状态空间中相近的区域，实现了``状态级''的目标驱动对齐。

\hypertarget{patch-embedding-ux524dux7aefux7684ux6548ux679c}{%
\subsubsection{14.3.4 Patch Embedding
前端的效果}\label{patch-embedding-ux524dux7aefux7684ux6548ux679c}}

在图像序列任务（sMNIST,
pMNIST）上，我们对比了直接将扁平化像素作为注入向量与使用卷积Patch
Embedding前端的效果。 * \textbf{结果}：使用Patch
Embedding（如4x4卷积，步长4）显著提升了模型性能，并加快了训练收敛。这是因为卷积层引入了\textbf{局部感受野}和\textbf{平移不变性}的归纳偏置，使得注入到环中的初始特征\texttt{J(x)}已经是具有空间层次结构的有效表示，减轻了环形动力学学习低级视觉特征的负担。这证明了MQR可以灵活地与成熟的表征学习模块结合，构建更强大的混合架构。

\hypertarget{ux7efcux5408ux4efbux52a1ux6027ux80fdux5bf9ux6bd4}{%
\subsection{14.4
综合任务性能对比}\label{ux7efcux5408ux4efbux52a1ux6027ux80fdux5bf9ux6bd4}}

表14.1汇总了MQR及其变体与基线模型在多个基准任务上的综合性能。

\textbf{表14.1：MQR与基线模型在多个任务上的性能对比摘要}

\begin{longtable}[]{@{}lllllll@{}}
\toprule
\begin{minipage}[b]{0.12\columnwidth}\raggedright
模型\strut
\end{minipage} & \begin{minipage}[b]{0.12\columnwidth}\raggedright
pMNIST (Acc.\%)\strut
\end{minipage} & \begin{minipage}[b]{0.12\columnwidth}\raggedright
PTB (PPL)\strut
\end{minipage} & \begin{minipage}[b]{0.12\columnwidth}\raggedright
复制记忆 (L=500, Acc.\%)\strut
\end{minipage} & \begin{minipage}[b]{0.12\columnwidth}\raggedright
时间序列预测 (MSE)\strut
\end{minipage} & \begin{minipage}[b]{0.12\columnwidth}\raggedright
参数数量 (M)\strut
\end{minipage} & \begin{minipage}[b]{0.12\columnwidth}\raggedright
训练时间/epoch (相对值)\strut
\end{minipage}\tabularnewline
\midrule
\endhead
\begin{minipage}[t]{0.12\columnwidth}\raggedright
\textbf{Vanilla RNN}\strut
\end{minipage} & \begin{minipage}[t]{0.12\columnwidth}\raggedright
90.2\strut
\end{minipage} & \begin{minipage}[t]{0.12\columnwidth}\raggedright
120.5\strut
\end{minipage} & \begin{minipage}[t]{0.12\columnwidth}\raggedright
12.3\strut
\end{minipage} & \begin{minipage}[t]{0.12\columnwidth}\raggedright
0.85\strut
\end{minipage} & \begin{minipage}[t]{0.12\columnwidth}\raggedright
\textasciitilde0.5\strut
\end{minipage} & \begin{minipage}[t]{0.12\columnwidth}\raggedright
1.0x\strut
\end{minipage}\tabularnewline
\begin{minipage}[t]{0.12\columnwidth}\raggedright
\textbf{LSTM}\strut
\end{minipage} & \begin{minipage}[t]{0.12\columnwidth}\raggedright
97.5\strut
\end{minipage} & \begin{minipage}[t]{0.12\columnwidth}\raggedright
78.9\strut
\end{minipage} & \begin{minipage}[t]{0.12\columnwidth}\raggedright
98.1\strut
\end{minipage} & \begin{minipage}[t]{0.12\columnwidth}\raggedright
0.42\strut
\end{minipage} & \begin{minipage}[t]{0.12\columnwidth}\raggedright
\textasciitilde2.0\strut
\end{minipage} & \begin{minipage}[t]{0.12\columnwidth}\raggedright
1.8x\strut
\end{minipage}\tabularnewline
\begin{minipage}[t]{0.12\columnwidth}\raggedright
\textbf{GRU}\strut
\end{minipage} & \begin{minipage}[t]{0.12\columnwidth}\raggedright
97.3\strut
\end{minipage} & \begin{minipage}[t]{0.12\columnwidth}\raggedright
80.1\strut
\end{minipage} & \begin{minipage}[t]{0.12\columnwidth}\raggedright
97.8\strut
\end{minipage} & \begin{minipage}[t]{0.12\columnwidth}\raggedright
0.43\strut
\end{minipage} & \begin{minipage}[t]{0.12\columnwidth}\raggedright
\textasciitilde1.5\strut
\end{minipage} & \begin{minipage}[t]{0.12\columnwidth}\raggedright
1.6x\strut
\end{minipage}\tabularnewline
\begin{minipage}[t]{0.12\columnwidth}\raggedright
\textbf{Orthogonal RNN}\strut
\end{minipage} & \begin{minipage}[t]{0.12\columnwidth}\raggedright
98.1\strut
\end{minipage} & \begin{minipage}[t]{0.12\columnwidth}\raggedright
75.3\strut
\end{minipage} & \begin{minipage}[t]{0.12\columnwidth}\raggedright
99.9\strut
\end{minipage} & \begin{minipage}[t]{0.12\columnwidth}\raggedright
0.40\strut
\end{minipage} & \begin{minipage}[t]{0.12\columnwidth}\raggedright
\textasciitilde0.5\strut
\end{minipage} & \begin{minipage}[t]{0.12\columnwidth}\raggedright
2.1x\strut
\end{minipage}\tabularnewline
\begin{minipage}[t]{0.12\columnwidth}\raggedright
\textbf{Transformer (Enc.)}\strut
\end{minipage} & \begin{minipage}[t]{0.12\columnwidth}\raggedright
\textbf{99.1}\strut
\end{minipage} & \begin{minipage}[t]{0.12\columnwidth}\raggedright
\textbf{65.2}\strut
\end{minipage} & \begin{minipage}[t]{0.12\columnwidth}\raggedright
N/A\strut
\end{minipage} & \begin{minipage}[t]{0.12\columnwidth}\raggedright
0.38\strut
\end{minipage} & \begin{minipage}[t]{0.12\columnwidth}\raggedright
\textasciitilde3.5\strut
\end{minipage} & \begin{minipage}[t]{0.12\columnwidth}\raggedright
0.7x (短序列)\strut
\end{minipage}\tabularnewline
\begin{minipage}[t]{0.12\columnwidth}\raggedright
\textbf{MQR (基础版)}\strut
\end{minipage} & \begin{minipage}[t]{0.12\columnwidth}\raggedright
98.4\strut
\end{minipage} & \begin{minipage}[t]{0.12\columnwidth}\raggedright
77.8\strut
\end{minipage} & \begin{minipage}[t]{0.12\columnwidth}\raggedright
\textbf{100.0}\strut
\end{minipage} & \begin{minipage}[t]{0.12\columnwidth}\raggedright
0.39\strut
\end{minipage} & \begin{minipage}[t]{0.12\columnwidth}\raggedright
\textasciitilde0.8\strut
\end{minipage} & \begin{minipage}[t]{0.12\columnwidth}\raggedright
2.3x\strut
\end{minipage}\tabularnewline
\begin{minipage}[t]{0.12\columnwidth}\raggedright
\textbf{MQR (复数+非线性)}\strut
\end{minipage} & \begin{minipage}[t]{0.12\columnwidth}\raggedright
98.6\strut
\end{minipage} & \begin{minipage}[t]{0.12\columnwidth}\raggedright
74.5\strut
\end{minipage} & \begin{minipage}[t]{0.12\columnwidth}\raggedright
\textbf{100.0}\strut
\end{minipage} & \begin{minipage}[t]{0.12\columnwidth}\raggedright
\textbf{0.36}\strut
\end{minipage} & \begin{minipage}[t]{0.12\columnwidth}\raggedright
\textasciitilde0.9\strut
\end{minipage} & \begin{minipage}[t]{0.12\columnwidth}\raggedright
2.5x\strut
\end{minipage}\tabularnewline
\begin{minipage}[t]{0.12\columnwidth}\raggedright
\textbf{MQR (Patch Emb.)}\strut
\end{minipage} & \begin{minipage}[t]{0.12\columnwidth}\raggedright
\textbf{99.0}\strut
\end{minipage} & \begin{minipage}[t]{0.12\columnwidth}\raggedright
-\strut
\end{minipage} & \begin{minipage}[t]{0.12\columnwidth}\raggedright
-\strut
\end{minipage} & \begin{minipage}[t]{0.12\columnwidth}\raggedright
-\strut
\end{minipage} & \begin{minipage}[t]{0.12\columnwidth}\raggedright
\textasciitilde1.1\strut
\end{minipage} & \begin{minipage}[t]{0.12\columnwidth}\raggedright
2.0x\strut
\end{minipage}\tabularnewline
\bottomrule
\end{longtable}

\emph{注：PTB困惑度为验证集结果；训练时间为相对比较，基于固定硬件和批量大小；Transformer在复制记忆任务上因完全不同的计算范式未直接比较；MQR参数数量随环尺寸N和注入秩r变化，此处为典型配置。}

\textbf{综合分析}： 1.
\textbf{性能}：MQR在需要极长程记忆的算法任务（如复制记忆）上展现了\textbf{统治级性能}，实现了完美记忆。在标准分类和预测任务上，MQR达到或超越了传统RNN变体和正交RNN的性能，与最先进的模型（如Transformer）竞争力相当，在某些任务上（如时间序列预测）甚至更优。这表明MQR不仅解决了长程依赖的理论问题，其实际表征能力也十分强大。
2.
\textbf{效率}：MQR的训练时间通常比LSTM/GRU长，主要源于弛豫迭代的额外计算。但与正交RNN处于同一量级。其参数量通常少于同等性能的LSTM或Transformer，模型较为紧凑。\textbf{其最大的效率优势体现在处理超长序列时的稳定内存消耗和不受序列长度影响的推理时间}。
3. \textbf{灵活性}：消融实验表明，MQR的各个高级组件（复数、非线性、Patch
Embedding、双酉解耦）都能在特定场景下带来可观的性能提升，体现了其作为一个\textbf{可扩展框架}的潜力。

\hypertarget{ux672cux7ae0ux5c0fux7ed3-1}{%
\subsection{14.5 本章小结}\label{ux672cux7ae0ux5c0fux7ed3-1}}

本章通过系统性的实验评估，全面验证了莫比乌斯量子环网络的理论主张与实际效能。实验结果表明：

\begin{enumerate}
\def\labelenumi{\arabic{enumi}.}
\tightlist
\item
  MQR成功解决了长期依赖建模的根本难题，在合成记忆任务上表现完美，其``弛豫到稳态''的动力学机制从原理上规避了传统RNN的梯度传播问题。
\item
  MQR在整个训练过程中表现出卓越的数值稳定性，梯度行为良好，损失景观平滑，这得益于其基于双随机/酉矩阵的底层动力学和隐函数求导的反向传播方式。
\item
  在多项实际基准
\end{enumerate}

\hypertarget{ux7b2c15ux7ae0-ux9ad8ux7ea7ux7279ux6027ux4e0eux6269ux5c55ux67b6ux6784ux589eux5f3amqrux7684ux8868ux8fbeux4e0eux6cdbux5316ux80fdux529b}{%
\section{第15章
高级特性与扩展架构：增强MQR的表达与泛化能力}\label{ux7b2c15ux7ae0-ux9ad8ux7ea7ux7279ux6027ux4e0eux6269ux5c55ux67b6ux6784ux589eux5f3amqrux7684ux8868ux8fbeux4e0eux6cdbux5316ux80fdux529b}}

\hypertarget{ux5f15ux8a00ux8d85ux8d8aux57faux7840ux73afux5f62ux52a8ux529bux5b66}{%
\subsection{15.1
引言：超越基础环形动力学}\label{ux5f15ux8a00ux8d85ux8d8aux57faux7840ux73afux5f62ux52a8ux529bux5b66}}

在前述章节中，我们详细阐述了莫比乌斯量子环网络（Möbius Quantum Ring
Network,
MQR）的核心架构原理与基础实现。该架构通过幺模随机流形约束、固定点松弛动力学和低秩注入机制，在理论上解决了传统循环神经网络中稳定性与表达能力之间的根本矛盾。然而，一个真正强大且实用的神经网络架构必须具备良好的可扩展性和适应性，能够通过引入精心设计的扩展特性来应对多样化的任务需求、提升模型性能，并在不破坏核心理论保证的前提下增强其表达能力。

本章将深入探讨MQR框架中一系列高级特性与扩展架构的设计原理、实现细节及其理论意义。这些特性并非简单的工程技巧叠加，而是基于对环形动力学系统数学本质的深刻理解而构建的。它们分别从不同的维度对基础MQR模型进行增强：

\begin{enumerate}
\def\labelenumi{\arabic{enumi}.}
\tightlist
\item
  \textbf{非线性扩展}：在保持收敛性证明的前提下，将非线性激活函数引入注入端和环内松弛过程，以增强模型的非线性表达能力。
\item
  \textbf{自保持混合机制}：通过引入有效哈密顿量 (H\_\{eff\})
  来调控环形动力学的混合速率，抑制过度平均化，保留更多局部信息。
\item
  \textbf{复数酉推理模式}：将动力学从实数域的双随机矩阵传播扩展至复数域的酉矩阵传播，使相位信息参与计算，探索更丰富的表示空间。
\item
  \textbf{双酉矩阵解耦策略}：将连接酉矩阵分解为可学习的``策略''部分和冻结的``世界模型''部分，实现更稳定、更高效的知识表示与策略学习分离。
\item
  \textbf{可学习目标平衡态}：为模型引入与任务或类别相关的可学习目标状态，使环形系统能够被驱动至与特定语义对齐的平衡态，实现状态级的监督。
\item
  \textbf{原型距离读出机制}：采用基于距离的读出方式，将环形稳态与可学习的原型向量进行比较，为分类任务提供一种几何直观且解释性更强的输出方式。
\end{enumerate}

这些扩展特性共同构成了MQR框架的``工具箱''，使得研究者与实践者能够根据具体任务的数据特性、计算约束和性能要求，灵活地组合与配置模型，从而在保持理论优雅性的同时，获得卓越的实践性能。本章将逐一剖析这些特性的设计动机、数学模型、实现算法及其对系统行为的影响。

\hypertarget{ux975eux7ebfux6027ux6269ux5c55ux5728ux6536ux655bux6027ux7ea6ux675fux4e0bux5f15ux5165ux6fc0ux6d3bux51fdux6570}{%
\subsection{15.2
非线性扩展：在收敛性约束下引入激活函数}\label{ux975eux7ebfux6027ux6269ux5c55ux5728ux6536ux655bux6027ux7ea6ux675fux4e0bux5f15ux5165ux6fc0ux6d3bux51fdux6570}}

基础MQR模型的核心松弛动力学 ( h \leftarrow (1-\alpha)hH\^{}T +
\alphaJ(x) ) 本质上是线性的（假设注入 (J(x))
也是线性的）。虽然双随机矩阵 (H)
的谱性质保证了固定点的存在性与唯一性，以及迭代的线性收敛速率，但纯线性的变换可能限制模型捕捉复杂非线性模式的能力。为此，MQR框架允许在严格的理论约束下，在两个关键位置引入非线性激活函数。

\hypertarget{ux6ce8ux5165ux7aefux975eux7ebfux6027}{%
\subsubsection{15.2.1
注入端非线性}\label{ux6ce8ux5165ux7aefux975eux7ebfux6027}}

在低秩注入模块 (J(x) = W\_\{up\} \cdot (W\_\{down\} \cdot x))
中，我们可以在降维矩阵 (W\_\{down\}) 之后引入一个非线性激活函数
(\sigma(\cdot))，形成 (J(x) = W\_\{up\}
\cdot \sigma(W\_\{down\} \cdot x))。这种形式与标准的LoRA（Low-Rank
Adaptation）或两层感知机的前馈部分完全一致。

\textbf{设计动机}： -
\textbf{增强特征变换}：非线性激活（如ReLU、GELU、SiLU）能够使注入模块对输入
(x)
进行非线性映射，提取更抽象、更具判别性的特征，然后再将其注入环形系统。这对于处理图像、文本等高度非线性结构的数据至关重要。
- \textbf{维持注入低秩性}：由于激活函数是逐点操作的，它不改变矩阵
(W\_\{up\}) 和 (W\_\{down\}) 的秩，因此注入算子 (J)
的整体输出仍然是由低秩结构主导的，保持了参数效率高的优点。 -
\textbf{对动力学无额外约束}：注入值 (J(x))
是作为松弛迭代的偏置项（bias）存在的。其非线性变换发生在迭代过程之外，因此不影响环形动力学
(h \leftarrow (1-\alpha)hH\^{}T + \alphaJ(x))
本身的收敛性分析。固定点方程 (h\^{}* = (1-\alpha)h\^{}\emph{H\^{}T +
\alphaJ(x)) 依然是一个关于 (h\^{}})
的线性方程，其解的存在唯一性由矩阵 ((I - (1-\alpha)H\^{}T))
的可逆性保证，而这仅取决于 (H) 的性质。

\textbf{实现考量}： 在代码实现中，这通常通过一个配置标志（如
\texttt{injection\_nonlinearity=\textquotesingle{}relu\textquotesingle{}}）来控制。注入模块
\texttt{LoRAInjection} 会检查该标志，并在前向传播中决定是否在
\texttt{W\_down} 和 \texttt{W\_up} 之间插入激活层。

\hypertarget{ux73afux5185ux677eux5f1bux975eux7ebfux6027}{%
\subsubsection{15.2.2
环内松弛非线性}\label{ux73afux5185ux677eux5f1bux975eux7ebfux6027}}

一个更大胆的扩展是在环形动力学迭代本身内部引入非线性，即修改更新规则为：
{[} h\^{}\{(t+1)\} = \sigma\left( (1-\alpha) h\^{}\{(t)\} H\^{}T +
\alpha J(x) \right) {]} 其中 (\sigma(\cdot))
是一个逐点应用的激活函数。

\textbf{核心挑战与理论约束}：
直接引入任意非线性可能破坏压缩映射性质，导致固定点不存在、不唯一，或迭代过程发散。为了保证收敛性，MQR框架要求所使用的激活函数
(\sigma) 必须是 \textbf{1-Lipschitz} 且 \textbf{非扩张（non-expansive）}
的。

\textbf{1-Lipschitz 条件}：对于任意向量 (u, v)，有 (\textbar{}\sigma(u)
- \sigma(v)\textbar{} \leq \textbar u -
v\textbar)。这意味着激活函数不会放大输入向量之间的距离。

\textbf{非扩张性}：通常与1-Lipschitz相关，确保函数不会导致迭代序列的范数无限增长。

满足这些条件的典型激活函数包括： - \textbf{逐点绝对值}：(\sigma(z) =
\textbar z\textbar)（在实数域） - \textbf{ReLU}：(\sigma(z) = \max(0,
z))（是1-Lipschitz的） - \textbf{带裁剪的线性函数}。 -
\textbf{逐点缩放的双曲正切（tanh）}：虽然tanh的导数绝对值小于1，但其本身不是全局1-Lipschitz，需谨慎分析其在特定区间内的行为。

\textbf{收敛性证明概要}： 定义映射 (\Phi(h) = \sigma\left( (1-\alpha) h
H\^{}T + \alpha J(x) \right))。我们需要证明 (\Phi)
是一个压缩映射。 1. 由于 (H) 是双随机矩阵，其谱范数
(\textbar H\textbar\_2 = 1)（最大奇异值为1）。因此，线性部分 ((1-\alpha)
h H\^{}T) 的变换是 ((1-\alpha))-收缩的。 2. 由于 (\sigma)
是1-Lipschitz，应用 (\sigma) 不会增加距离。 3. 结合两者，对于任意 (h\_1,
h\_2)： {[}

\begin{aligned}
    \|\Phi(h_1) - \Phi(h_2)\| &\leq \| \left[(1-\alpha) h_1 H^T + \alpha J(x)\right] - \left[(1-\alpha) h_2 H^T + \alpha J(x)\right] \| \\
    &= (1-\alpha) \| (h_1 - h_2) H^T \| \\
    &\leq (1-\alpha) \|h_1 - h_2\| \cdot \|H^T\| \\
    &= (1-\alpha) \|h_1 - h_2\|
    \end{aligned}

{]} 由于 (0 \textless{} \alpha \leq 1)，有 (0 \leq (1-\alpha)
\textless{} 1)，因此 (\Phi) 是一个压缩映射。根据巴拿赫不动点定理，迭代
(h\^{}\{(t+1)\} = \Phi(h\^{}\{(t)\})) 必定收敛到唯一的固定点 (h\^{}*)。

\textbf{作用与影响}：
环内非线性允许在状态传播的每一步进行非线性变换，极大地增强了模型每一层（迭代步）的表达能力。它可以使系统演化出更复杂、非线性的吸引子动态。然而，这也可能使得固定点的解析形式更难分析，并且需要确保激活函数严格满足
Lipschitz 条件，这在实现时需特别注意（例如，避免使用像标准 Sigmoid
这样导数可能大于1的函数）。

\begin{figure}
\centering
\includegraphics{figures/diagram_14.png}
\caption{图 14}
\end{figure}

\emph{图15-1：MQR非线性扩展决策流程图。展示了在注入端和环内松弛两个环节，根据配置决定是否引入非线性激活函数（\sigma\_inj
和 \sigma\_ring）的数据流与控制流。其中环内非线性 \sigma\_ring
必须选用1-Lipschitz函数以保证收敛性。}

\hypertarget{ux81eaux4fddux6301ux6df7ux5408ux673aux5236ux8c03ux63a7ux4fe1ux606fux6269ux6563ux901fux7387}{%
\subsection{15.3
自保持混合机制：调控信息扩散速率}\label{ux81eaux4fddux6301ux6df7ux5408ux673aux5236ux8c03ux63a7ux4fe1ux606fux6269ux6563ux901fux7387}}

在基础MQR中，环形动力学完全由双随机矩阵 (H = \textbar U\textbar\^{}2)
驱动。双随机矩阵的一个重要性质是，其平稳分布是均匀分布。这意味着，在多次迭代后，环上任何节点的信息都会均匀扩散到所有其他节点，导致节点状态趋向于一致（``过度平均化''）。虽然这保证了全局信息融合，但也可能抹杀对任务至关重要的局部特征或结构信息。

\hypertarget{ux6709ux6548ux54c8ux5bc6ux987fux91cf-h_eff-ux7684ux5f15ux5165}{%
\subsubsection{15.3.1 有效哈密顿量 (H\_\{eff\})
的引入}\label{ux6709ux6548ux54c8ux5bc6ux987fux91cf-h_eff-ux7684ux5f15ux5165}}

为了在保持双随机性和收敛性的前提下，对信息扩散速率进行细粒度控制，MQR引入了\textbf{自保持混合（Self-Retaining
Mixing）机制}。其核心思想是定义一个\textbf{有效哈密顿量（Effective
Hamiltonian）} (H\_\{eff\})，用它来替代原始的 (H) 参与松弛迭代： {[}
H\_\{eff\} = (1-\beta) I + \beta H {]} 其中： - (I) 是 (N \times N)
的单位矩阵。 - (\beta \in [0, 1])
是一个可调的超参数，称为\textbf{混合系数（Mixing Coefficient）}。 - (H)
是原始的幺模随机连接矩阵。

\hypertarget{ux6570ux5b66ux6027ux8d28ux5206ux6790-1}{%
\subsubsection{15.3.2
数学性质分析}\label{ux6570ux5b66ux6027ux8d28ux5206ux6790-1}}

\begin{enumerate}
\def\labelenumi{\arabic{enumi}.}
\tightlist
\item
  \textbf{双随机性保持}：

  \begin{itemize}
  \tightlist
  \item
    (I) 是双随机的。
  \item
    (H) 是双随机的。
  \item
    双随机矩阵的凸组合（加权和）仍然是双随机矩阵。因为对于任意
    (i)，(\sum\_j {[}(1-\beta)I\_\{ij\} + \beta H\_\{ij\}{]} =
    (1-\beta)*1 + \beta*1 = 1)，列和同理。因此，(H\_\{eff\})
    是双随机的。
  \end{itemize}
\item
  \textbf{收敛性保持}：

  \begin{itemize}
  \tightlist
  \item
    松弛迭代变为 (h \leftarrow (1-\alpha)h H\_\{eff\}\^{}T +
    \alphaJ(x))。
  \item
    由于 (H\_\{eff\}) 仍然是双随机的，其谱范数
    (\textbar H\_\{eff\}\textbar\_2
    \leq 1)。实际上，可以证明其第二大特征值的模长
    (\textbar{}\lambda\_2\textbar{} \leq (1-\beta) +
    \beta\textbar{}\lambda\_2(H)\textbar{} )，其中 (\lambda\_2(H)) 是
    (H) 的第二大特征值。这确保了迭代映射的压缩性，固定点定理依然成立。
  \end{itemize}
\item
  \textbf{物理意义解释}：

  \begin{itemize}
  \tightlist
  \item
    当 (\beta = 1) 时，(H\_\{eff\} =
    H)，模型退化为原始MQR，信息完全通过学习的连接进行混合。
  \item
    当 (\beta = 0) 时，(H\_\{eff\} = I)。此时迭代变为 (h
    \leftarrow (1-\alpha)h +
    \alphaJ(x))。这意味着每个节点几乎独立地根据注入值进行更新，节点间几乎没有信息交换。系统迅速收敛，但缺乏全局协调。
  \item
    当 (0 \textless{} \beta \textless{} 1) 时，(H\_\{eff\}) 是单位矩阵和
    (H) 的插值。单位矩阵项 ((1-\beta)I)
    赋予了每个节点一种``自保持（self-retention）''或``惯性''，使其更倾向于保留自己当前的状态，而不是无条件地将信息传递给邻居。这有效地\textbf{减缓了信息在整个环上均匀化的速度}，允许局部模式或梯度在收敛前维持更长时间。
  \end{itemize}
\end{enumerate}

\hypertarget{ux5e94ux7528ux573aux666fux4e0eux8c03ux4f18}{%
\subsubsection{15.3.3
应用场景与调优}\label{ux5e94ux7528ux573aux666fux4e0eux8c03ux4f18}}

自保持混合机制在以下场景中尤为有用： -
\textbf{处理具有强局部相关性数据}：例如，在图像分类中，相邻像素块（对应环上相邻节点组）的特征可能高度相关。过快的全局混合可能破坏这种局部结构。通过降低
(\beta)，可以让模型更关注局部特征的整合。 -
\textbf{防止梯度过早平滑}：在训练初期，(H)
可能尚未学到有意义的连接模式。使用较小的 (\beta)
可以提供一个更稳定、更保守的混合基线，有助于训练的稳定。 -
\textbf{作为正则化}：(\beta)
可以作为一个额外的超参数，用于控制模型的容量和泛化能力。较小的 (\beta)
可能起到正则化作用，防止过拟合。

在实践中，(\beta)
可以设置为固定值，也可以设计为可学习的参数（尽管需要确保其值域在{[}0,1{]}内），甚至可以作为门控机制的一部分，由输入自适应地决定。

\hypertarget{ux590dux6570ux9149ux63a8ux7406ux6a21ux5f0fux5229ux7528ux76f8ux4f4dux4fe1ux606f}{%
\subsection{15.4
复数酉推理模式：利用相位信息}\label{ux590dux6570ux9149ux63a8ux7406ux6a21ux5f0fux5229ux7528ux76f8ux4f4dux4fe1ux606f}}

基础MQR在实数域操作：酉矩阵 (U) 被参数化并用于生成实数的双随机矩阵 (H =
\textbar U\textbar\^{}2)，松弛动力学也在实数向量空间中进行。然而，酉矩阵
(U) 本身是复数域的，其元素包含幅度和相位信息。将 (H)
定义为逐元素幅度平方，实际上丢弃了 (U)
中所有的相位信息。相位在量子力学、信号处理和复数神经网络中被证明可以编码重要的相对关系和周期模式。

\hypertarget{ux590dux6570ux52a8ux529bux5b66ux5b9aux4e49}{%
\subsubsection{15.4.1
复数动力学定义}\label{ux590dux6570ux52a8ux529bux5b66ux5b9aux4e49}}

MQR的\textbf{复数酉推理模式}旨在利用完整的酉矩阵
(U)，在复数域执行环形动力学。其修改如下：

\begin{enumerate}
\def\labelenumi{\arabic{enumi}.}
\tightlist
\item
  \textbf{复数状态}：隐藏状态 (h \in C\^{}\{N \times d\})
  变为复数矩阵，其中 (d) 是特征维度。
\item
  \textbf{复数传播}：松弛迭代规则修改为： {[} h \leftarrow (1-\alpha) h
  U\^{}\dagger + \alpha J\_c(x) {]} 其中：

  \begin{itemize}
  \tightlist
  \item
    (U \in C\^{}\{N \times N\}) 是学习的酉矩阵。
  \item
    (U\^{}\dagger) 是 (U) 的共轭转置（Hermitian conjugate）。
  \item
    (J\emph{c(x): R\^{}\{d}\{in\}\}
    \to C\^{}\{N \times d\})
    是一个复数注入模块。这可以通过两个独立的LoRA分支分别生成实部和虚部实现：(J\emph{c(x)
    = J}\{re\}(x) + i \cdot J\_\{im\}(x))。
  \item
    迭代在复数域进行，加法和乘法均为复数运算。
  \end{itemize}
\item
  \textbf{收敛性分析}：

  \begin{itemize}
  \tightlist
  \item
    酉矩阵 (U) 是等距算子：对于任意复数向量 (v)，有 (\textbar U
    v\textbar{} = \textbar v\textbar)。因此，线性部分 ((1-\alpha) h
    U\^{}\dagger) 是 ((1-\alpha))-收缩的（在复数L2范数下）。
  \item
    如果注入是有限的，并且我们考虑实数域和复数域范数的兼容性，压缩映射论证依然成立。因此，复数迭代同样收敛到唯一的复数固定点
    (h\^{}*\_c \in C\^{}\{N \times d\})。
  \end{itemize}
\end{enumerate}

\hypertarget{ux8bfbux51faux4e0eux6d4bux91cf}{%
\subsubsection{15.4.2 读出与测量}\label{ux8bfbux51faux4e0eux6d4bux91cf}}

获得复数稳态 (h\^{}*\emph{c)
后，不能直接将其输入给通常为实数权重的读出层
(W}\{readout\})。需要进行\textbf{测量（Measurement）}，将复数信息转换为实数。MQR提供了几种测量方式：

\begin{enumerate}
\def\labelenumi{\arabic{enumi}.}
\tightlist
\item
  \textbf{幅度读出（默认）}：取每个复数元素的模（幅度）(\textbar h\^{}\emph{\emph{c\textbar)，得到一个实数值矩阵，然后进行局部采样和线性读出：(y
  = W}\{readout\} \cdot \textbar h\^{}}\_c\textbar\_\{S\})。
\item
  \textbf{实部/虚部读出}：分别取实部 (\Re(h\^{}*\_c)) 或虚部
  (\Im(h\^{}*\_c))。
\item
  \textbf{相位读出}：取相位角
  (\angle h\^{}*\_c)，但需注意相位是周期性的，直接使用可能不连续。
\item
  \textbf{联合特征}：将幅度和相位（或实部和虚部）拼接起来作为特征。
\end{enumerate}

\textbf{优势}： -
\textbf{更丰富的表示}：相位可以编码循环、振荡、干涉等模式，对于音频、时序信号、具有周期结构的数据可能特别有效。
- \textbf{参数效率}：一个 (N \times N)
的复数酉矩阵所蕴含的信息量理论上大于一个相同大小的实数双随机矩阵。 -
\textbf{理论连贯性}：直接操作酉矩阵 (U)，而非其衍生量
(H)，在数学上更为本质。

\textbf{挑战}： -
\textbf{计算复杂度}：复数运算（尤其是矩阵乘法）的计算成本通常是实数运算的2-4倍。
- \textbf{优化难度}：在复数流形（酉群）上进行优化需要专门的技巧。 -
\textbf{梯度流动}：复数反向传播需要遵循Wirtinger导数或CR（Cauchy-Riemann）演算，实现更复杂。

\hypertarget{section}{%
\subsection{15.}\label{section}}

\end{document}
